% L’agriculture de grandes plantations — Géographie Tle A — Cours signé OnBuch+
% Programme officiel MINESEC. Compiler avec : tectonic N07-agriculture-grandes-plantations.tex
\newcommand{\DOCMATIERE}{Géographie}
\newcommand{\DOCNIVEAU}{Tle A}
\newcommand{\DOCMODULE}{Module 1 : Le Cameroun}
\newcommand{\DOCLECON}{N07}
\newcommand{\DOCTITRE}{L’agriculture de grandes plantations}
% OnBuch+ — préambule « cours signés » ENRICHI (compilé avec tectonic / XeLaTeX)
% Système visuel « Pop » d'OnBuch+ : fond crème (jamais blanc), bordures encre
% épaisses, ombres « dures » décalées, titres Archivo Black en capitales.
% Version MATHÉMATIQUES : graphes (pgfplots/tikz), boîtes pédagogiques
% (définition, propriété, méthode, attention, le savais-tu, exercice résolu,
% exercice, TP, bilan), page de garde et signature OnBuch+ ancrée en bas.
%
% Macros attendues dans main.tex AVANT \input{preamble} :
%   \DOCMATIERE  \DOCNIVEAU  \DOCMODULE  \DOCLECON (numéro)  \DOCTITRE
\documentclass[11pt]{article}

\usepackage{fontspec}
\usepackage[french]{babel}
\usepackage[a4paper,margin=2cm,top=3.4cm,bottom=2.8cm,headheight=1.4cm,footskip=1.3cm]{geometry}
\usepackage{amsmath,amssymb,amsfonts}
\usepackage{mathtools} % \xmapsto, \coloneqq... souvent utilisés par les modèles
\usepackage{cancel} % simplifications barrées (bilans de forces)
% \abs{z} (module, valeur absolue) et \norm{u} (norme), utilisés par les modèles
\providecommand{\abs}{}\let\abs\relax\DeclarePairedDelimiter{\abs}{\lvert}{\rvert}
\providecommand{\norm}{}\let\norm\relax\DeclarePairedDelimiter{\norm}{\lVert}{\rVert}
\usepackage[table]{xcolor} % [table] : fournit aussi \rowcolors (alternance de lignes)
\usepackage{pagecolor}
\usepackage{tikz}
\usetikzlibrary{arrows.meta,positioning,calc,shapes.geometric,shapes.misc,shapes.arrows,decorations.pathmorphing,decorations.pathreplacing,decorations.markings,patterns,shadows,mindmap,trees,fit,backgrounds,matrix,intersections,angles,quotes}
\usepackage{pgfplots}
\usepackage[version=4]{mhchem} % équations nucléaires \ce{^{235}_{92}U}
\usepackage[european,siunitx]{circuitikz} % schémas électriques
\pgfplotsset{compat=1.18}
\usepgfplotslibrary{fillbetween,groupplots} % aires entre courbes (calcul intégral)
\usepackage{enumitem}
\usepackage{fancyhdr}
% [most] charge toutes les bibliothèques tcolorbox (listings, minted, raster,
% documentation...) et gonfle inutilement la mémoire TeX sur un document long
% (~300 boîtes) : on ne charge que ce qui est réellement utilisé.
\usepackage[skins,breakable]{tcolorbox}
\usepackage{graphicx}
\graphicspath{{/home/user/t-t/geo-onbuch/images/}}
\usepackage{colortbl}
\usepackage{booktabs}
\usepackage{tabularx}
\usepackage{multirow}
\usepackage{diagbox} % case d'en-tête diagonale des tableaux à double entrée
% Colonnes extensibles centrée / gauche / droite (C, L, R), souvent utilisées par les modèles
\newcolumntype{C}{>{\centering\arraybackslash}X}
\newcolumntype{L}{>{\raggedright\arraybackslash}X}
\newcolumntype{R}{>{\raggedleft\arraybackslash}X}
\usepackage{array}
\usepackage{multicol}
\usepackage{siunitx}
% parse-numbers=false : accepte aussi \SI{2,5 \times 10^{-3}}{...}
\sisetup{locale=FR,output-decimal-marker={,},group-minimum-digits=4,parse-numbers=false}
\DeclareSIUnit\ohm{\ensuremath{\symup{\Omega}}} % U+2126 absent de la police
\DeclareSIUnit\division{div} % réglages d'oscilloscope (ms/div, V/div)
\DeclareSIUnit\tr{tr} % tours (vitesse de rotation, ex. \SI{1500}{\tr\per\minute})
\DeclareSIUnit\molar{M} % concentration molaire (ex. \SI{3}{\molar}, \SI{1}{\micro\molar})
\DeclareSIUnit\an{an} % année (le modèle confond parfois avec \year, un compteur TeX) — ex. \SI{5}{\kilo\meter\per\an}
\DeclareSIUnit\FCFA{FCFA} % franc CFA (ex. \SI{500}{\FCFA\per\kilo\gram})
\DeclareSIUnit\degreeBrix{\textdegree Brix} % teneur en sucre d'un sirop/jus (réfractométrie)
\DeclareSIUnit\month{mois} % le modèle utilise parfois \month, un compteur TeX, comme unité de durée
\DeclareSIUnit\microliter{\micro\liter} % le modèle écrit parfois \microliter en un seul token
\DeclareSIUnit\million{millions} % dénombrement (ex. \SI{15}{\million\per\mL} spermatozoïdes)
\usepackage{qrcode}
% \degree : commande fréquemment utilisée par les modèles pour un angle ou une
% température en dehors de \SI{}{} (non fournie par les paquets chargés).
\providecommand{\degree}{^{\circ}}
% \qed (fin de démonstration), utilisé par les modèles sans amsthm
\providecommand{\qed}{\hfill$\square$}
% \barre : double barre des valeurs interdites dans un tableau de variations (tkz-tab)
\providecommand{\barre}{\ensuremath{\|}}
% environnement proof (amsthm non chargé) utilisé par les modèles
\ifdefined\proof\else\newenvironment{proof}[1][Démonstration]{\par\noindent\textit{#1.}\ }{\qed\par}\fi
\usepackage{lastpage}
\usepackage{hyperref}
\hypersetup{hidelinks}

% ============================================================
% Palette « Pop » OnBuch+ — mêmes hex que la classe P (plus_ui.dart)
% ============================================================
\definecolor{poporange}{HTML}{F59321}
\definecolor{popdark}{HTML}{E07A0C}
\definecolor{popcream}{HTML}{FFF6E8}
\definecolor{popink}{HTML}{1C1714}
\definecolor{poppurple}{HTML}{7A5AE0}
\definecolor{popgreen}{HTML}{2E9E6A}
\definecolor{poppink}{HTML}{E0457B}
\definecolor{popblue}{HTML}{2F7DE1}
\definecolor{popgold}{HTML}{C98A12}
\definecolor{popmuted}{HTML}{8A7F76}
\definecolor{popline}{HTML}{EADFD2}
\definecolor{popgray}{HTML}{8A7F76}
\colorlet{popgrey}{popgray}
% Alias tolérants (le modèle nomme parfois les couleurs différemment)
\colorlet{popwhite}{white}
\colorlet{popblack}{popink}
\colorlet{popred}{poppink}
\colorlet{popyellow}{popgold}
% Teintes claires (fonds de tableaux/encadrés) — le modèle nomme parfois
% les couleurs avec un suffixe L (ex. popblueL) sans les avoir définies.
\colorlet{poporangeL}{poporange!20}
\colorlet{popdarkL}{popdark!20}
\colorlet{poppurpleL}{poppurple!20}
\colorlet{popgreenL}{popgreen!20}
\colorlet{poppinkL}{poppink!20}
\colorlet{popblueL}{popblue!20}
\colorlet{popgoldL}{popgold!20}
\colorlet{popredL}{popred!20}
\colorlet{popgrayL}{popgray!20}
% Teintes claires pour fonds de tableaux / zones
\colorlet{popblueL}{popblue!10!white}
\colorlet{poporangeL}{poporange!14!white}
\colorlet{popgreenL}{popgreen!12!white}
\colorlet{poppurpleL}{poppurple!12!white}
\colorlet{poppinkL}{poppink!12!white}
\colorlet{popgoldL}{popgold!14!white}

\pagecolor{popcream}

% ============================================================
% Typographie
% ============================================================
\defaultfontfeatures{Ligatures=TeX}
\setmainfont{PlusJakartaSans-Medium.ttf}[Path=../../../fonts/,BoldFont=PlusJakartaSans-Bold.ttf]
% Maths en sans-serif (Fira Math) avec chiffres et lettres droites de Plus
% Jakarta, pour que les formules se fondent dans le texte.
\usepackage[math-style=ISO,bold-style=ISO]{unicode-math}
\setmathfont{FiraMath-Regular.otf}[Path=../../../fonts/]
\setmathfont{PlusJakartaSans-Medium.ttf}[Path=../../../fonts/,range={up/num,up/{Latin,latin},\mathup}]
\usepackage{icomma} % virgule décimale sans espace en mode math
% Caractères mathématiques Unicode absents des polices de texte (Plus Jakarta,
% Archivo Black) : rendus par leur équivalent mathématique, en texte comme en
% formule — sinon ils s'affichent en carré vide (ex. « Arithmétique dans ℤ »).
\usepackage{newunicodechar}
\newunicodechar{ℝ}{\ensuremath{\mathbb{R}}}
\newunicodechar{ℤ}{\ensuremath{\mathbb{Z}}}
\newunicodechar{ℂ}{\ensuremath{\mathbb{C}}}
\newunicodechar{ℕ}{\ensuremath{\mathbb{N}}}
\newunicodechar{ℚ}{\ensuremath{\mathbb{Q}}}
\newunicodechar{𝔻}{\ensuremath{\mathbb{D}}}
\newunicodechar{α}{\ensuremath{\alpha}}
\newunicodechar{β}{\ensuremath{\beta}}
\newunicodechar{γ}{\ensuremath{\gamma}}
\newunicodechar{δ}{\ensuremath{\delta}}
\newunicodechar{ε}{\ensuremath{\varepsilon}}
\newunicodechar{θ}{\ensuremath{\theta}}
\newunicodechar{λ}{\ensuremath{\lambda}}
\newunicodechar{μ}{\ensuremath{\mu}}
\newunicodechar{σ}{\ensuremath{\sigma}}
\newunicodechar{τ}{\ensuremath{\tau}}
\newunicodechar{φ}{\ensuremath{\varphi}}
\newunicodechar{ψ}{\ensuremath{\psi}}
\newunicodechar{ω}{\ensuremath{\omega}}
\newunicodechar{Σ}{\ensuremath{\Sigma}}
% majuscules produites par \MakeUppercase dans les titres (α-aminés -> Α)
\newunicodechar{Α}{\ensuremath{\alpha}}
\newunicodechar{Β}{\ensuremath{\beta}}
\newunicodechar{Γ}{\ensuremath{\Gamma}}
\newunicodechar{Δ}{\ensuremath{\Delta}}
\newunicodechar{Ε}{\ensuremath{\varepsilon}}
\newunicodechar{Θ}{\ensuremath{\Theta}}
\newunicodechar{Λ}{\ensuremath{\Lambda}}
\newunicodechar{Μ}{\ensuremath{\mu}}
\newunicodechar{Τ}{\ensuremath{\tau}}
\newunicodechar{Φ}{\ensuremath{\Phi}}
\newunicodechar{Ψ}{\ensuremath{\Psi}}
\newunicodechar{Ω}{\ensuremath{\Omega}}
\newunicodechar{↦}{\ensuremath{\mapsto}}
\newunicodechar{∈}{\ensuremath{\in}}
\newunicodechar{∉}{\ensuremath{\notin}}
\newunicodechar{∧}{\ensuremath{\wedge}}
\newunicodechar{⇔}{\ensuremath{\Leftrightarrow}}
\newunicodechar{⟺}{\ensuremath{\Longleftrightarrow}}
\newunicodechar{⇒}{\ensuremath{\Rightarrow}}
\newunicodechar{⟹}{\ensuremath{\Longrightarrow}}
\newunicodechar{∘}{\ensuremath{\circ}}
\newunicodechar{∪}{\ensuremath{\cup}}
\newunicodechar{∩}{\ensuremath{\cap}}
\newunicodechar{≡}{\ensuremath{\equiv}}
\newunicodechar{⊂}{\ensuremath{\subset}}
\newunicodechar{⊥}{\ensuremath{\perp}}
\newunicodechar{∃}{\ensuremath{\exists}}
\newunicodechar{∀}{\ensuremath{\forall}}
\newunicodechar{∛}{\ensuremath{\sqrt[3]{\,}}}
\newunicodechar{∜}{\ensuremath{\sqrt[4]{\,}}}
\newunicodechar{⃗}{\ensuremath{{}^{\to}}}
\newunicodechar{ⁿ}{\ifmmode^{n}\else\textsuperscript{\ensuremath{n}}\fi}
\newunicodechar{ˣ}{\ifmmode^{x}\else\textsuperscript{\ensuremath{x}}\fi}
\newunicodechar{ᵇ}{\ifmmode^{b}\else\textsuperscript{\ensuremath{b}}\fi}
\newunicodechar{ʳ}{\ifmmode^{r}\else\textsuperscript{\ensuremath{r}}\fi}
\newunicodechar{ᵘ}{\ifmmode^{u}\else\textsuperscript{\ensuremath{u}}\fi}
\newunicodechar{ʸ}{\ifmmode^{y}\else\textsuperscript{\ensuremath{y}}\fi}
\newunicodechar{ᵃ}{\ifmmode^{a}\else\textsuperscript{\ensuremath{a}}\fi}
\newunicodechar{ᵉ}{\ifmmode^{e}\else\textsuperscript{\ensuremath{e}}\fi}
\newunicodechar{ᵏ}{\ifmmode^{k}\else\textsuperscript{\ensuremath{k}}\fi}
\newunicodechar{ᵖ}{\ifmmode^{p}\else\textsuperscript{\ensuremath{p}}\fi}
\newunicodechar{ᴺ}{\ifmmode^{N}\else\textsuperscript{\ensuremath{N}}\fi}
\newunicodechar{ᶜ}{\ifmmode^{c}\else\textsuperscript{\ensuremath{c}}\fi}
\newunicodechar{ʰ}{\ifmmode^{h}\else\textsuperscript{\ensuremath{h}}\fi}
\newunicodechar{ᵠ}{\ifmmode^{q}\else\textsuperscript{\ensuremath{q}}\fi}
\newunicodechar{ₙ}{\ifmmode_{n}\else\textsubscript{\ensuremath{n}}\fi}
\newunicodechar{ₐ}{\ifmmode_{a}\else\textsubscript{\ensuremath{a}}\fi}
\newunicodechar{ᵢ}{\ifmmode_{i}\else\textsubscript{\ensuremath{i}}\fi}
\newunicodechar{ᵣ}{\ifmmode_{r}\else\textsubscript{\ensuremath{r}}\fi}
\newunicodechar{ₖ}{\ifmmode_{k}\else\textsubscript{\ensuremath{k}}\fi}
\newunicodechar{ᵦ}{\ifmmode_{\beta}\else\textsubscript{\ensuremath{\beta}}\fi}
\newunicodechar{ₓ}{\ifmmode_{x}\else\textsubscript{\ensuremath{x}}\fi}
\newfontfamily\popdisplay{ArchivoBlack-Regular.ttf}[Path=../../../fonts/]
\newfontfamily\poplogo{SpaceGrotesk-Bold.ttf}[Path=../../../fonts/]
\newfontfamily\popsemi{PlusJakartaSans-SemiBold.ttf}[Path=../../../fonts/]
\newfontfamily\popxbold{PlusJakartaSans-ExtraBold.ttf}[Path=../../../fonts/]
\newfontfamily\popxboldls{PlusJakartaSans-ExtraBold.ttf}[Path=../../../fonts/,LetterSpace=3.0]

\setlist[enumerate]{leftmargin=1.7em,itemsep=5pt,topsep=4pt}
\setlist[itemize]{leftmargin=1.7em,itemsep=4pt,topsep=4pt,label={\color{poporange}\textbullet}}
\setlength{\parindent}{0pt}
\setlength{\parskip}{5pt}
\renewcommand{\arraystretch}{1.25}

% pgfplots : style Pop par défaut
\pgfplotsset{
  every axis/.append style={axis line style={popink,line width=1pt},tick style={popink},
    grid style={popline,line width=0.5pt},label style={font=\small},tick label style={font=\footnotesize}},
  popcurve/.style={poporange,line width=1.8pt,no markers,smooth},
}
% Même style disponible dans un \draw TikZ simple (hors pgfplots)
\tikzset{popcurve/.style={poporange,line width=1.8pt}}
\tikzset{popgrid/.style={popline,very thin}}
\colorlet{popcurve}{poporange} % le nom du style est parfois utilisé comme couleur

% ============================================================
% Pill / squiggle
% ============================================================
\newcommand{\poppill}[3]{%
  \tikz[baseline=-0.55ex]\node[draw=popink,line width=1.2pt,rounded corners=2.6pt,%
    fill=#2,inner xsep=6.5pt,inner ysep=3pt]{%
    \popxboldls\fontsize{7}{7}\selectfont\color{#3}\MakeUppercase{#1}};%
}
\newcommand{\popsquiggle}[1]{%
  \tikz[baseline=-0.4ex]{
    \draw[#1,line width=2.6pt,line cap=round]
      plot [smooth] coordinates {(0,0.09) (0.34,0.01) (0.68,0.21) (1.02,0.01) (1.36,0.10)};
  }%
}
% Mot-clé surligné (marqueur orange sous le texte)
\newcommand{\cle}[1]{{\popxbold\color{popdark}#1}}

% ============================================================
% En-tête / pied
% ============================================================
\pagestyle{fancy}
\fancyhf{}
\renewcommand{\headrulewidth}{0pt}
\renewcommand{\footrulewidth}{0pt}
\lhead{\raisebox{-0.15ex}{\poplogo\fontsize{13}{13}\selectfont\color{popink}OnBuch\color{poporange}+}}
\rhead{\poppill{\DOCMATIERE\ \textbullet\ \DOCNIVEAU\ \textbullet\ Leçon \DOCLECON}{white}{popink}}
\lfoot{\popxbold\fontsize{7.6}{7.6}\selectfont\color{popmuted}\MakeUppercase{Cours signé OnBuch+}\ \textbullet\ {\popsemi\DOCTITRE}}
\rfoot{\tikz[baseline=-0.6ex]\node[rounded corners=8.5pt,fill=popink,minimum height=17pt,minimum width=17pt,inner xsep=5pt,inner sep=0]{\popxbold\fontsize{8}{8}\selectfont\color{popcream}\thepage\,/\,\pageref*{LastPage}};}

% ============================================================
% Titres
% ============================================================
\newcommand{\chaptitre}[2]{%
  \begin{tcolorbox}[enhanced,colback=poporange,colframe=popink,boxrule=2.6pt,arc=11pt,
      drop shadow=popink,left=15pt,right=15pt,top=12pt,bottom=11pt,before skip=3pt,after skip=18pt]
    \poppill{#2}{popink}{popcream}\par\vspace{9pt}
    {\raggedright\hyphenpenalty=10000\exhyphenpenalty=10000\popdisplay\fontsize{22}{24}\selectfont\color{popcream}\MakeUppercase{#1}\par}
  \end{tcolorbox}
}
% Section numérotée : pastille orange avec le numéro + titre Archivo
\newcounter{popsec}
\newcommand{\coursec}[1]{%
  \stepcounter{popsec}\setcounter{popsub}{0}%
  \par\vspace{16pt}\noindent
  \begin{minipage}{\linewidth}
  \tikz[baseline=-0.7ex]\node[draw=popink,line width=1.6pt,fill=poporange,rounded corners=4pt,minimum size=22pt,inner sep=0]{\popdisplay\fontsize{12}{12}\selectfont\color{popcream}\thepopsec};%
  \hspace{8pt}{\popdisplay\fontsize{15}{17}\selectfont\color{popink}\MakeUppercase{#1}}\par
  \vspace{4pt}\noindent{\color{poporange}\rule{36pt}{3.6pt}}
  \end{minipage}\par\nopagebreak\vspace{8pt}
}
\newcounter{popsub}
\newcommand{\courssub}[1]{%
  \stepcounter{popsub}%
  \par\vspace{9pt}\noindent{\popxbold\fontsize{11.5}{13}\selectfont\color{popink}{\color{poporange}\thepopsec.\thepopsub}\hspace{6pt}#1}\par\nopagebreak\vspace{4pt}
}

% ============================================================
% Boîtes pédagogiques — même recette : fond blanc, bordure épaisse colorée,
% ombre dure encre, badge plein. Titre optionnel [#1].
% ============================================================
\tcbset{popbase/.style={breakable,enhanced,colback=white,boxrule=2.3pt,arc=9pt,
  shadow={3pt}{-3pt}{0pt}{popink},left=14pt,right=14pt,top=11pt,bottom=11pt,before skip=11pt,after skip=15pt,
  fontupper=\popsemi\fontsize{10.6}{14.5}\selectfont\color{popink},
  fontlower=\popsemi\fontsize{10.6}{14.5}\selectfont\color{popink}}}
% Coche dessinée (le glyphe ✓ n'existe pas dans les polices du document)
\newcommand{\popcheck}[1]{\tikz[baseline=-0.5ex]\draw[#1,line width=1.6pt,line cap=round,line join=round](-0.13,0.0)--(-0.04,-0.09)--(0.13,0.1);}
\providecommand{\coche}{\popcheck{popgreen}}
\providecommand{\resultat}[1]{\par\smallskip\noindent\fcolorbox{popgreen}{popgreenL}{\parbox{\dimexpr\linewidth-2\fboxsep-2\fboxrule\relax}{\textbf{Résultat.} #1}}\par\smallskip}
\newcommand{\popboxhead}[3]{\poppill{#1}{#2}{white}\ifx\relax#3\relax\else\hspace{6pt}{\popxbold\fontsize{10.6}{12}\selectfont\color{popink}#3}\fi\par\vspace{7pt}}

\newtcolorbox{aretenir}[1][]{popbase,colframe=popgreen,before upper={\popboxhead{À retenir}{popgreen}{#1}}}
\newtcolorbox{exemplebox}[1][]{popbase,colframe=poporange,before upper={\popboxhead{Exemple}{poporange}{#1}}}
\newtcolorbox{definition}[1][]{popbase,colframe=popblue,before upper={\popboxhead{Définition}{popblue}{#1}}}
\newtcolorbox{propriete}[1][]{popbase,colframe=poppurple,before upper={\popboxhead{Propriété}{poppurple}{#1}}}
\newtcolorbox{methode}[1][]{popbase,colframe=popgold,before upper={\popboxhead{Méthode}{popgold}{#1}}}
\newtcolorbox{attention}[1][]{popbase,colframe=poppink,before upper={\popboxhead{Attention}{poppink}{#1}}}
\newtcolorbox{experience}[1][]{popbase,colframe=popblue,colback=popblueL,before upper={\popboxhead{Expérience}{popblue}{#1}}}
% « Le savais-tu ? » : carte violette pleine (contexte, culture, Cameroun)
\newtcolorbox{savaistu}[1][]{popbase,colback=poppurple,colframe=popink,
  fontupper=\popsemi\fontsize{10.4}{14.2}\selectfont\color{white},
  before upper={\poppill{Le savais-tu\,?}{white}{poppurple}\ifx\relax#1\relax\else\hspace{6pt}{\popxbold\color{white}#1}\fi\par\vspace{7pt}}}
% Exercice résolu : énoncé en haut, \tcblower puis la solution sur fond crème
\newcounter{popexo}\newcounter{popexr}
\newtcolorbox{exoresolu}[1][]{popbase,colframe=poporange,colbacklower=popcream,
  segmentation style={popink,line width=1.2pt,dashed},
  before upper={\stepcounter{popexr}\popboxhead{Exercice résolu \thepopexr}{poporange}{#1}},
  before lower={\poppill{Solution}{popink}{popcream}\par\vspace{6pt}}}
% Exercice (non résolu en place) — la correction est dans la partie Corrigés
\newtcolorbox{exercice}[1][]{popbase,colframe=popink,
  before upper={\stepcounter{popexo}\popboxhead{Exercice \thepopexo}{popink}{#1}}}
% Corrigé
\newtcolorbox{corrige}[1][]{popbase,colframe=popgreen,colback=popgreenL,
  before upper={\popboxhead{Corrigé}{popgreen}{#1}}}
% Objectifs / prérequis / bilan
\newtcolorbox{objectifs}{popbase,colframe=popink,colback=poporangeL,
  before upper={\popboxhead{Objectifs de la leçon}{popink}{}}}
\newtcolorbox{prerequis}{popbase,colframe=popink,colback=popblueL,
  before upper={\popboxhead{Prérequis}{popblue}{}}}
\newtcolorbox{bilan}{popbase,colframe=popink,colback=white,boxrule=2.8pt,
  before upper={\popboxhead{Fiche bilan}{poporange}{L'essentiel à connaître pour l'examen}}}
% Figure encadrée avec légende
\newtcolorbox{popfigure}[1][]{enhanced,breakable=false,colback=white,colframe=popline,boxrule=1.4pt,arc=8pt,
  left=10pt,right=10pt,top=10pt,bottom=8pt,before skip=10pt,after skip=12pt,halign=center,
  lower separated=false,halign lower=center,
  fontlower=\popsemi\fontsize{9}{11.5}\selectfont\color{popmuted},
  #1}
\newcommand{\legende}[1]{\par\vspace{6pt}{\popsemi\fontsize{9}{11.5}\selectfont\color{popmuted}#1}}

% ============================================================
% Signature OnBuch+
% ============================================================
\newcommand{\poplogoinline}[1]{{\poplogo\fontsize{#1}{#1}\selectfont\color{popink}OnBuch\color{poporange}+}}
\newcommand{\signatureonbuch}{%
  \begin{tcolorbox}[enhanced,colback=popink,colframe=popink,boxrule=2.2pt,arc=10pt,
      drop shadow=poporange,left=14pt,right=14pt,top=10pt,bottom=10pt,before skip=0pt,after skip=0pt]
    \tikz[baseline=-0.7ex]\node[circle,fill=popgreen,draw=popcream,line width=1.3pt,minimum size=22pt,inner sep=0]{\popcheck{white}};%
    \hspace{9pt}\begin{minipage}[c]{0.62\linewidth}
      {\popxbold\fontsize{12}{13}\selectfont\color{popcream}Signé\ \textbullet\ {\poplogo OnBuch\color{poporange}+}}\par\vspace{2pt}
      {\raggedright\popsemi\fontsize{8}{10.5}\selectfont\color{popline}\MakeUppercase{Équipe pédagogique OnBuch+}\\\MakeUppercase{Conforme au programme officiel MINESEC}\par}
    \end{minipage}\hfill
    {\poplogo\fontsize{22}{22}\selectfont\color{popcream}OnBuch\color{poporange}+}
  \end{tcolorbox}%
}

% ============================================================
% Page de garde
% #1 titre · #2 matière · #3 niveau · #4 module · #5 n° leçon · #6 durée ·
% #7 description / objectifs (texte ou itemize)
% ============================================================
\newcommand{\pagegarde}[7]{%
  \thispagestyle{empty}
  \newgeometry{a4paper,margin=1.7cm,top=1.6cm,bottom=1.4cm}%
  \begin{tikzpicture}[remember picture,overlay]
    \fill[poporange!18] ([xshift=-3.2cm,yshift=-3.6cm]current page.north east) circle (4.6cm);
    \fill[poppurple!14] ([xshift=2.4cm,yshift=7.6cm]current page.south west) circle (3.8cm);
  \end{tikzpicture}%
  {\poplogo\fontsize{30}{30}\selectfont\color{popink}OnBuch\color{poporange}+}\hfill
  \raisebox{0.6ex}{\poppill{Cours signé \textbullet\ Premium}{popink}{popcream}}\par
  \vspace{1pt}\popsquiggle{poppurple}\par
  \vspace{8pt}
  \begin{tcolorbox}[enhanced,colback=poporange,colframe=popink,boxrule=2.8pt,arc=13pt,
      drop shadow=popink,left=18pt,right=18pt,top=12pt,bottom=13pt,after skip=10pt]
    \poppill{#2 \textbullet\ #3}{popink}{popcream}\hspace{5pt}\poppill{#4}{white}{popink}\par\vspace{8pt}
    {\popdisplay\fontsize{14}{14}\selectfont\color{popink}LEÇON #5}\par\vspace{4pt}
    {\raggedright\hyphenpenalty=10000\exhyphenpenalty=10000\popdisplay\fontsize{25}{27}\selectfont\color{popcream}\MakeUppercase{#1}\par}
    \vspace{8pt}
    {\popxbold\fontsize{9.5}{9.5}\selectfont\color{popink}\MakeUppercase{Durée conseillée : #6}}
  \end{tcolorbox}
  \begin{tcolorbox}[enhanced,colback=white,colframe=popink,boxrule=2.2pt,arc=10pt,
      drop shadow=popink,left=16pt,right=16pt,top=10pt,bottom=10pt,after skip=8pt]
    \poppill{Dans cette leçon}{poppurple}{white}\par\vspace{5pt}
    \popsemi\fontsize{9.3}{12.6}\selectfont\color{popink}#7
  \end{tcolorbox}
  \noindent\begin{minipage}[t]{0.487\linewidth}
    \begin{tcolorbox}[enhanced,colback=popgreen,colframe=popink,boxrule=2.2pt,arc=10pt,
        drop shadow=popink,left=11pt,right=11pt,top=10pt,bottom=10pt,height=3.1cm,valign=center]
      \tcbox[colback=white,colframe=popink,boxrule=1.3pt,arc=5pt,boxsep=0pt,nobeforeafter,
        left=3pt,right=3pt,top=3pt,bottom=3pt,valign=center]{\qrcode[height=1.9cm]{https://play.google.com/store/apps/details?id=com.luvvix.onbuch&hl=fr}}%
      \hspace{8pt}\begin{minipage}[c]{0.5\linewidth}
        {\popdisplay\fontsize{11}{12.5}\selectfont\color{white}\MakeUppercase{Télécharge OnBuch}}\par\vspace{4pt}
        {\raggedright\popsemi\fontsize{8.4}{11}\selectfont\color{white}Gratuit \textbullet\ Google Play\\Tuteur IA, annales, concours, résultats\par}
      \end{minipage}
    \end{tcolorbox}
  \end{minipage}\hfill
  \begin{minipage}[t]{0.487\linewidth}
    \begin{tcolorbox}[enhanced,colback=poppurple,colframe=popink,boxrule=2.2pt,arc=10pt,
        drop shadow=popink,left=11pt,right=11pt,top=10pt,bottom=10pt,height=3.1cm,valign=center]
      \tikz[baseline=-0.7ex]\node[circle,fill=white,draw=popink,line width=1.3pt,minimum size=1.6cm,inner sep=0]{\popdisplay\fontsize{24}{24}\selectfont\color{poppurple}+};%
      \hspace{8pt}\begin{minipage}[c]{0.6\linewidth}
        {\popdisplay\fontsize{11}{12.5}\selectfont\color{white}\MakeUppercase{Passe à OnBuch+}}\par\vspace{4pt}
        {\raggedright\popsemi\fontsize{8.4}{11}\selectfont\color{white}Tous les cours signés, Tuteur IA illimité, fiches PDF — onglet OnBuch+ de l'app\par}
      \end{minipage}
    \end{tcolorbox}
  \end{minipage}
  \vspace{8pt}
  \popcontenu
  \vfill
  \signatureonbuch
  \restoregeometry
  \newpage
}

% Rangée « Ce cours contient » (page de garde)
\newcommand{\poptile}[3]{% #1 couleur #2 gros texte #3 légende
  \begin{tcolorbox}[enhanced,colback=#1,colframe=popink,boxrule=2pt,arc=9pt,shadow={2.5pt}{-2.5pt}{0pt}{popink},
      left=6pt,right=6pt,top=9pt,bottom=9pt,halign=center,valign=center,height=1.75cm,width=\linewidth,nobeforeafter]
    {\popdisplay\fontsize{12.5}{13}\selectfont\color{white}\MakeUppercase{#2}}\par\vspace{3pt}
    {\centering\popsemi\fontsize{7.8}{9.5}\selectfont\color{white}#3\par}
  \end{tcolorbox}}
\newcommand{\popcontenu}{%
  \par\noindent{\popxboldls\fontsize{7.5}{7.5}\selectfont\color{popmuted}CE COURS SIGNÉ CONTIENT}\par\vspace{7pt}
  \noindent\begin{minipage}[t]{0.235\linewidth}\poptile{popblue}{Cours}{complet, illustré, pas à pas}\end{minipage}\hfill
  \begin{minipage}[t]{0.235\linewidth}\poptile{poporange}{Méthodes}{et exercices résolus}\end{minipage}\hfill
  \begin{minipage}[t]{0.235\linewidth}\poptile{poppink}{Exercices}{type examen + corrigés}\end{minipage}\hfill
  \begin{minipage}[t]{0.235\linewidth}\poptile{popgreen}{Fiche bilan}{l'essentiel à retenir}\end{minipage}\par}

% Sommaire
\newtcolorbox{sommaire}{enhanced,colback=white,colframe=popink,boxrule=2.2pt,arc=10pt,
  drop shadow=popink,left=18pt,right=18pt,top=13pt,bottom=13pt,before skip=6pt,after skip=20pt,
  before upper={\poppill{Sommaire}{poppurple}{white}\par\vspace{9pt}\popsemi\fontsize{10.6}{14.5}\selectfont\color{popink}}}

% Signature de fin de cours — poussée tout en bas de la dernière page
\newcommand{\findecours}{%
  \par\vfill\nopagebreak
  \begin{center}\popsquiggle{poporange}\end{center}\vspace{4pt}
  \signatureonbuch
}
\AtBeginDocument{\def\llbracket{\mathopen{[\mkern-3.5mu[}}\def\rrbracket{\mathclose{]\mkern-3.5mu]}}} % intervalles d'entiers (glyphe ⟦ absent de la police)
% Glyphes absents de Fira Math : redéfinis à partir de glyphes disponibles
\AtBeginDocument{%
  \def\setminus{\mathbin{\backslash}}\let\smallsetminus\setminus
  \def\vdots{\mathord{\vbox{\baselineskip3.2pt\lineskiplimit0pt\kern4pt\hbox{.}\hbox{.}\hbox{.}}}}%
  \def\ddots{\mathinner{\mkern1mu\raise6.5pt\hbox{.}\mkern2mu\raise3.6pt\hbox{.}\mkern2mu\raise0.7pt\hbox{.}\mkern1mu}}%
  \def\triangle{\mathord{\scalebox{0.9}{$\symup{\Delta}$}}}%
  \def\nabla{\mathord{\raisebox{\depth}{\scalebox{1}[-1]{$\symup{\Delta}$}}}}%
  \def\checkmark{\popcheck{popdark}}%
  \def\star{\mathbin{\ast}}\def\bigstar{\mathord{\ast}}%
  \def\diamond{\mathbin{\scalebox{0.6}{$\Diamond$}}}%
  \def\bigcup{\mathop{\mathchoice{\raisebox{-0.3em}{\scalebox{1.7}{$\cup$}}}{\scalebox{1.25}{$\cup$}}{\cup}{\cup}}\displaylimits}%
  \def\bigcap{\mathop{\mathchoice{\raisebox{-0.3em}{\scalebox{1.7}{$\cap$}}}{\scalebox{1.25}{$\cap$}}{\cap}{\cap}}\displaylimits}%
  \def\lesseqqgtr{\mathrel{\lessgtr}}\def\bigoplus{\mathop{\oplus}\displaylimits}%
}
\providecommand{\ssi}{si et seulement si} % abréviation fréquente
\AtBeginDocument{\providecommand{\wideparen}{\overparen}} % arc de cercle

%% FIGFIX-BEGIN v3 — correctifs globaux des figures (docs/onbuch-generation/tools/figfix)
\usepackage{adjustbox}\usepackage{etoolbox}
% toute figure TikZ/pgfplots plus large que la ligne est réduite à la largeur disponible
\BeforeBeginEnvironment{tikzpicture}{\begin{adjustbox}{max width=\linewidth}}
\AfterEndEnvironment{tikzpicture}{\end{adjustbox}}
% légendes pgfplots lisibles par-dessus les courbes
\pgfplotsset{every axis/.append style={legend style={fill=white,fill opacity=0.92,text opacity=1,draw=black!15,inner sep=3pt}}}
% étiquettes d'arêtes (syntaxe edge["..."]) avec fond blanc
\tikzset{every edge quotes/.append style={fill=white,inner sep=1.2pt}}
% bulles de cartes mentales : texte à la ligne dans la bulle, bulle un peu plus grande
\tikzset{every concept/.append style={text width=2.0cm,align=center,minimum size=2.3cm,inner sep=2pt}}
%% FIGFIX-END
\begin{document}

\pagegarde{L’agriculture de grandes plantations}{Géographie}{Tle A}{Module 1 : Le Cameroun}{N07}{2 h + dossier 2 (2 h)}{Cette leçon étudie l'agriculture de grandes plantations au Cameroun : ses caractéristiques (superficies, méthodes, financement), ses types (plantations des grandes firmes comme la CDC, la SOCAPALM, la SODECOTON ou la SOSUCAM, et plantations des particuliers), ainsi que ses impacts socio-économiques. Elle s'achève par le dossier 2 consacré à la main-d'œuvre, vecteur d'intégration nationale.
\vspace{4pt}
\begin{multicols}{2}\small
\begin{itemize}
\item À la fin de cette leçon, je sais définir l'agriculture de grandes plantations et en dégager les caractéristiques (superficies, méthodes de travail, financement)
\item À la fin de cette leçon, je sais distinguer les types de grandes plantations (grandes firmes et particuliers)
\item À la fin de cette leçon, je sais localiser sur une carte du Cameroun les principales plantations des grandes firmes (CDC, SOCAPALM, SAFACAM, SODECOTON, SOSUCAM, HEVECAM) et les cultures associées
\item À la fin de cette leçon, je sais analyser l'impact socio-économique des grandes plantations (emplois, conditions de vie, accaparement des terres, pollution)
\item À la fin de cette leçon, je sais définir et illustrer les notions d'agriculture intensive et d'agropole
\item À la fin de cette leçon, je sais décrire l'organisation spatiale et sociale des camps de plantation (cadres, ouvriers, écoles, centres de santé)
\end{itemize}\end{multicols}}

\begin{sommaire}
\begin{multicols}{2}
\begin{enumerate}
\item Définition et caractéristiques de l'agriculture de grandes plantations
\item Les types de grandes plantations
\item L'impact socio-économique des grandes plantations
\item Dossier 2 : la main-d'œuvre dans les grandes plantations, vecteur d'intégration nationale
\item Travaux pratiques : cartographier et représenter les grandes plantations
\item Activité d'intégration
\item Méthodes et exercices résolus
\item Exercices
\item Corrigés des exercices
\item Fiche bilan
\end{enumerate}
\end{multicols}
\end{sommaire}

\begin{objectifs}
\begin{itemize}
\item À la fin de cette leçon, je sais définir l'agriculture de grandes plantations et en dégager les caractéristiques (superficies, méthodes de travail, financement)
\item À la fin de cette leçon, je sais distinguer les types de grandes plantations (grandes firmes et particuliers)
\item À la fin de cette leçon, je sais localiser sur une carte du Cameroun les principales plantations des grandes firmes (CDC, SOCAPALM, SAFACAM, SODECOTON, SOSUCAM, HEVECAM) et les cultures associées
\item À la fin de cette leçon, je sais analyser l'impact socio-économique des grandes plantations (emplois, conditions de vie, accaparement des terres, pollution)
\item À la fin de cette leçon, je sais définir et illustrer les notions d'agriculture intensive et d'agropole
\item À la fin de cette leçon, je sais décrire l'organisation spatiale et sociale des camps de plantation (cadres, ouvriers, écoles, centres de santé)
\item À la fin de cette leçon, je sais légender une carte, construire un croquis et un diagramme à partir de documents géographiques
\item À la fin de cette leçon, je sais inventorier, classer et interpréter des cartes, photographies et tableaux statistiques sur les plantations
\end{itemize}
\end{objectifs}

\begin{prerequis}
\begin{itemize}
\item Les grandes zones agro-écologiques du Cameroun (climats, sols) vues en classe de Première
\item La distinction entre agriculture vivrière et agriculture de rente (Première)
\item La notion de développement et d'intégration nationale (début d'année de Terminale)
\item La lecture et la légendation d'une carte thématique (savoir-faire de base)
\item La répartition de la population et des migrations intérieures au Cameroun
\end{itemize}
\end{prerequis}

\newpage

\coursec{L'agriculture de grandes plantations}

\noindent Imagine-toi sur la route nationale n°3, entre \textbf{Douala} et \textbf{Édéa}. Par la vitre du car, le paysage change brutalement : la forêt dense et la mosaïque des champs villageois laissent place à une géométrie parfaite. À perte de vue, des rangées d'arbres identiques — des palmiers à huile — dessinent un damier vert foncé sur des dizaines de kilomètres. Au détour d'un virage, tu aperçois un ensemble de bâtiments industriels : une huilerie d'où s'échappent des volutes de vapeur, des hangars de stockage, et, alignées le long de pistes rectilignes, des centaines de cases en tôle ou en dur : les \emph{camps} des travailleurs. Un panneau indique \textbf{SOCAPALM}, \emph{Société Camerounaise de Palmeraies}. Ton voisin de siège, un ancien ouvrier, te raconte que son père est venu du Nord dans les années 1980, que lui est né ici, que ses enfants vont à l'école de la plantation, se soignent à son dispensaire. Mais il te parle aussi des conflits avec les villages riverains qui réclament leurs terres, et de la couleur bizarre de l'eau dans le petit affluent de la Sanaga en aval de l'usine.

\medskip\noindent \textbf{Problématique :} Comment ces grandes plantations sont-elles organisées ? Quelles cultures y pratique-t-on et avec quels moyens techniques et financiers ? Contribuent-elles réellement au développement et à l'intégration nationale, ou créent-elles de nouveaux problèmes fonciers et environnementaux ?

\courssub{Qu'est-ce qu'une grande plantation ? : définition et opposition à l'exploitation paysanne}

\begin{definition}[Agriculture de grandes plantations]
L'agriculture de grandes plantations est une \cle{agriculture commerciale} de type \cle{intensif}, pratiquée en \cle{monoculture} sur de vastes superficies (plusieurs centaines à dizaines de milliers d'hectares), avec un fort investissement en \cle{capitaux} (machines, intrants, infrastructures) et en \cle{main-d'œuvre salariée}, orientée vers l'\cle{exportation} ou l'\cle{industrie de transformation}.
\end{definition}

\noindent Pour bien saisir la spécificité de ce modèle, opposons-le à l'agriculture paysanne qui domine encore le paysage camerounais. Le tableau ci-dessous résume les ruptures fondamentales :

\begin{center}
\begin{tabularx}{\linewidth}{|l|X|X|}
\hline
\rowcolor{popblueL}
\textbf{Critère} & \textbf{Petite exploitation paysanne (vivrière / rente)} & \textbf{Grande plantation (agro-industrielle)} \\
\hline
\textbf{Superficie} & 0,5 à 5 ha (moyenne \SI{1,5}{ha}) & \SI{500}{ha} à \SI{>50\,000}{ha} \\
\hline
\textbf{Objectif principal} & Autoconsommation + vente surplus (sécurité alimentaire) & Production marchande \textbf{intégrale} (export / usine) \\
\hline
\textbf{Assolement / Culture} & Polyculture, associations culturales, jachère & \textbf{Monoculture} spécialisée pérenne (palmier, hévéa, bananier, cotonnier, canne) \\
\hline
\textbf{Travail} & Familial, peu ou pas salarié, faible mécanisation & \textbf{Salarié} nombreux, hiérarchisé, \textbf{mécanisation} forte \\
\hline
\textbf{Capitaux / Intrants} & Faibles (fumure organique, semences paysannes) & \textbf{Massifs} (engrais chimiques, pesticides, variétés sélectionnées, irrigation) \\
\hline
\textbf{Statut foncier} & Droit coutumier, communauté / famille & \textbf{Concession} (bail emphytéotique État), titre foncier \\
\hline
\textbf{Transformation} & Quasi nulle sur place (vente brute) & \textbf{Intégrée} sur site (huilerie, cotonnière, sucrerie, usine de latex, conditionnement banane) \\
\hline
\end{tabularx}
\end{center}

\medskip\noindent \textbf{Observation documentée :} Une photographie aérienne de la zone de \textbf{Mbongo} (SOCAPALM, département de la Sanaga-Maritime) montre un quadrillage parfait de parcelles rectangulaires de \SI{100}{ha} chacune, séparées par des pistes carrossables de \SI{10}{m} de large. Aucune case, aucun arbre isolé ne rompt l'uniformité du couvert végétal sur des kilomètres carrés. À l'opposé, un cliché du village voisin de \textbf{Dibamba} révèle une mosaïque : cases, champs de manioc, bananiers, cacaoyers, jachères, haies vives, sentiers sinueux.

\begin{popfigure}
\begin{tikzpicture}[scale=0.85, every node/.style={font=\small}]
% --- Petite exploitation (gauche) ---
\begin{scope}[xshift=-7cm]
% Contour parcelle
\draw[popdark, thick, rounded corners=5pt] (-3,-2.5) rectangle (3,2.5);
% Maison
\draw[poporange, fill=poporangeL, thick] (-0.5,-0.2) rectangle (0.5,0.5);
\node at (0,0.15) {\tiny Case};
% Champs hétérogènes
\foreach \i/\c/\col in {1/manioc/popgreen!60, 2/maïs/popgold!80, 3/bananier/popgreen!40, 4/cacao/poporange!60, 5/jachère/popmuted} {
    \pgfmathsetmacro{\x}{-2.5 + (\i-1)*1.0}
    \draw[fill=\col, draw=popline] (\x,-2.3) rectangle (\x+0.8, -1.0);
    \node[rotate=90] at (\x+0.4, -1.65) {\tiny \c};
}
% Jachère / brousse
\draw[fill=popgreen!30!popmuted, draw=popline, pattern=north west lines, pattern color=popgreen!50!popmuted] (-2.5, 1.0) rectangle (2.5, 2.3);
\node at (0, 1.65) {\tiny Jachère / Brousse};
% Sentier
\draw[popline, double distance=2pt, double=popmuted] (0,-2.5) -- (0,-3.5);
\node[below] at (0,-3.5) {\tiny Sentier};
% Titre
\node[align=center, font=\bfseries\small, popdark] at (0, 3) {Petite exploitation vivrière\\ \textit{(environ 2 ha)}};
\end{scope}

% --- Grande plantation (droite) ---
\begin{scope}[xshift=7cm]
% Contour parcelle
\draw[popdark, thick, rounded corners=5pt] (-5,-3) rectangle (5,3);
% Quadrillage parcelles (ex: 4x3)
\foreach \x in {-4,-2,0,2} {
    \foreach \y in {-2,0,2} {
        \draw[fill=popgreen!70!popblue, draw=popgreen!50!popdark] (\x,\y) rectangle (\x+1.9, \y+1.9);
        % Symboles palmiers (points)
        \foreach \px in {0.2,0.9,1.6} {
            \foreach \py in {0.2,0.9,1.6} {
                \fill[popgreen!30!popdark] (\x+\px, \y+\py) circle (1.5pt);
            }
        }
    }
}
% Pistes
\draw[popmuted, line width=3pt] (-5,0) -- (5,0); % piste centrale E-W
\draw[popmuted, line width=3pt] (0,-3) -- (0,3); % piste centrale N-S
% Usine / Camp
\draw[poporange, fill=poporangeL, thick] (3.5, 1.5) rectangle (5.5, 2.8);
\node at (4.5, 2.15) {\tiny Usine / Huilerie};
\draw[popblue, fill=popblueL, thick] (3.5, -1.0) rectangle (5.5, 0.3);
\node at (4.5, -0.35) {\tiny Camp ouvriers};
% Flèche Nord
\draw[->, thick, popdark] (4, 2.5) -- (4, 3.2) node[above] {N};
% Échelle
\draw[|-|, thick] (-4.5, -3.3) -- (-2.5, -3.3) node[midway, below] {1 km};
% Titre
\node[align=center, font=\bfseries\small, popdark] at (0, 3.5) {Grande plantation (SOCAPALM type)\\ \textit{(bloc de 10 000 ha environ)}};
\end{scope}
\end{tikzpicture}
\legende{Schéma comparatif : structure spatiale d'une petite exploitation paysanne (hétérogène, polyculture, jachère) vs grande plantation (quadrillage rigoureux, monoculture, pistes rectilignes, infrastructures industrielles et sociales intégrées).}
\end{popfigure}

\begin{attention}[Piège fréquent : ne pas confondre « grande exploitation » et « grande plantation »]
Au Bac, ne dites jamais : « Le grand cultivateur de café de l'Ouest qui a 50 ha fait de la plantation. »
\begin{itemize}
    \item Une \cle{grande exploitation} paysanne reste une exploitation \textbf{familiale} (main-d'œuvre familiale prédominante, polyculture, faible capitalisation), même si elle dépasse 20 ha.
    \item La \cle{grande plantation} est une \textbf{entreprise capitaliste} (personne morale), salariée, monoculture, à forte intensité de capital par hectare.
\end{itemize}
La distinction est \textbf{juridique} (statut de la terre : concession vs droit coutumier), \textbf{économique} (logique de profit capitaliste vs reproduction sociale) et \textbf{technique} (itinéraire cultural standardisé).
\end{attention}

\courssub{Les trois piliers structurels : superficies, méthodes de travail, financement}

\paragraph{1. Des superficies d'échelle industrielle}
Les concessions accordées par l'État (baux emphytéotiques de 50 ans renouvelables, loi n°74/1 du 6 juillet 1974 modifiée) portent sur des surfaces massives. Cela permet l'\cle{économie d'échelle} : dilution des coûts fixes (usine, routes, recherche) sur un grand volume de production.

\begin{exemplebox}[La SOCAPALM : l'exemple-type de la palmeraie industrielle]
La \textbf{SOCAPALM} (filiale du groupe \emph{Socfinaf} / \emph{Bolloré}) exploite \textbf{environ 58 000 ha} de palmiers à huile matures (chiffres 2023, ordre de grandeur \SIrange{50\,000}{60\,000}{ha}) répartis sur \textbf{5 sites} : \textbf{Mbongo, Dibombari, Eseka, Kienké, Mbonjo}. Chaque site forme une \cle{unité de production} quasi autonome (plantation + huilerie + camps + école + dispensaire). La superficie totale concédée (incluant réserves foncières, infrastructures, zones non plantées) dépasse \SI{100\,000}{ha}.
\end{exemplebox}

\paragraph{2. Des méthodes de travail rationalisées : la monoculture intensive}
L'itinéraire technique est standardisé, issu de la recherche agronomique (IRAD, CIRAD, centres de recherche des firmes) :
\begin{itemize}
    \item \textbf{Matériel végétal sélectionné :} clones à haut rendement (ex. \emph{Tenera} pour le palmier, \emph{GT1} pour l'hévéa), précoces, résistants aux maladies.
    \item \textbf{Préparation du sol mécanique :} abattage, andryage, labour profond, billonnage (tracteurs \SI{150}{--}{250 ch}).
    \item \textbf{Densité de plantation optimisée :} \SI{143}{palmiers/ha} (triangle équilatéral \SI{9}{m}), \SI{500}{hévéas/ha}, \SI{2\,500}{bananiers/ha}.
    \item \textbf{Intrants chimiques massifs :} engrais NPK + Mg + B (formules spécifiques par stade), herbicides (glyphosate, pour tour de tronc), fongicides/insecticides (avion ou pulvérisateur porté pour hévéa).
    \item \textbf{Récolte organisée :} tournées hebdomadaires (palmier), saignée quotidienne (hévéa), coupe hebdomadaire (banane), ramassage mécanique (coton). Main-d'œuvre \textbf{spécialisée} (cueilleurs, saigneurs, conducteurs d'engins) encadrée par des \textbf{contremaîtres} et des \textbf{ingénieurs agronomes}.
\end{itemize}

\begin{definition}[Agriculture intensive]
L'agriculture intensive est un système de production qui obtient des \textbf{rendements élevés à l'hectare} grâce à d'importants apports de \textbf{capitaux} (machines, bâtiments, intrants), de \textbf{travail} (main-d'œuvre qualifiée nombreuse) et de \textbf{savoir-faire technique} (variétés, itinéraires), par opposition à l'agriculture extensive (faibles intrants, grands espaces, faible rendement/ha).
\end{definition}

\paragraph{3. Un financement lourd : capitaux mixtes et endettement}
L'investissement initial (mise en place) atteint \SI{5\,000}{--}{15\,000} €/ha selon la culture (palmier, hévéa, bananier). Le retour sur investissement prend 5 à 7 ans (immaturité). Le montage financier associe :
\begin{itemize}
    \item \textbf{Capitaux étrangers :} groupes multinationaux (Bolloré/Socfin pour SOCAPALM, SAFACAM ; GMG/Halcyon Agri pour HEVECAM ; Wilmar pour SODECOTON coton ; Somdiaa pour SOSUCAM ; CDC pour banane/palmier/hévéa).
    \item \textbf{Capitaux publics :} l'État camerounais détient souvent des parts minoritaires (ex. \SI{33}{\%} à la SOCAPALM, \SI{40}{\%} à la SODECOTON historiquement) via la \textbf{SN Investissement (SNI)}.
    \item \textbf{Emprunts bancaires :} banques commerciales (BICEC, Afriland, SCB) et institutions de développement (BEAC, BAD, Proparco, AFD).
\end{itemize}
\begin{propriete}[Société d'économie mixte (SEM) à capitaux majoritairement privés]
La plupart des grandes plantations camerounaises sont des \textbf{SEM} où le privé (souvent étranger) est actionnaire majoritaire et gestionnaire, l'État étant actionnaire minoritaire garant de l'intérêt national (emplois, recettes fiscales, aménagement du territoire).
\end{propriete}

\courssub{L'agropole : quand la plantation structure le territoire}

Au-delà de l'unité de production, les grandes plantations camerounaises fonctionnent souvent comme des \cle{agropoles} : des pôles territoriaux intégrant production, transformation de première main et services urbains.

\begin{definition}[Agropole]
Une \cle{agropole} est un \textbf{territoire organisé autour d'un pôle agro-industriel} qui articule : (1) la \textbf{production agricole} (plantation noyau + éventuels planteurs villageois sous contrat), (2) la \textbf{transformation industrielle} sur site (huilerie, cotonnière, sucrerie, usine de latex, conditionnement banane), (3) des \textbf{services urbains} (logement, santé, éducation, commerces, voirie, électricité, eau) profitant aux travailleurs et aux populations riveraines.
\end{definition}

\noindent \textbf{Exemple concret : l'agropole de \emph{Mbongo} (SOCAPALM / Sanaga-Maritime).}
\begin{itemize}
    \item \textbf{Noyau :} \SI{12\,000}{ha} de palmeraies + huilerie (\SI{60}{t/h} de régimes).
    \item \textbf{Population :} \SI{>10\,000}{habitants} (ouvriers, cadres, familles, commerçants attirés).
    \item \textbf{Équipements :} École primaire + secondaire (CEG), Centre de santé intégré (CSI) avec médecin, réseau eau/électricité, marché, stade, voirie bitumée interne.
    \item \textbf{Rayonnement :} Achats de régimes aux \emph{planteurs villageois} (PV) des environs (contrats d'approvisionnement), créant une zone d'influence de \SI{30}{km} de rayon.
\end{itemize}

\begin{savaistu}[L'huile de palme : un produit « tout-en-un »]
L'huile de palme brute (HPB) produite au Cameroun (\SI{\approx 450\,000}{t/an}, 2023) a trois débouchés majeurs :
\begin{enumerate}
    \item \textbf{Alimentation locale :} huile de cuisine (huile rouge raffinée), margarine (\emph{Oléine}, \emph{Stéarine}) — couvre \SI{>80}{\%} des besoins nationaux.
    \item \textbf{Industrie locale :} savonneries (\emph{Savonnerie du Cameroun}, \emph{Unilever}), cosmétiques, agro-alimentaire (biscuiteries, fritures industrielles).
    \item \textbf{Exportation :} vers la CEMAC (Tchad, RCA, Gabon), l'Europe (biocarburants, oléochimie) et l'Asie. Les recettes d'exportation oscillent autour de \SIrange{150}{200} milliards FCFA/an.
\end{enumerate}
C'est la \textbf{première culture d'exportation} du Cameroun en valeur (devant le cacao, le coton, le bois, la banane).
\end{savaistu}

\coursec{Les types de grandes plantations}

\courssub{Les plantations des grandes firmes agro-industrielles (noyaux dits « \emph{estates} »)}
Le programme cite explicitement six sociétés majeures. Les voici avec leurs cultures, localisations et ordres de grandeur (estimations 2023, surfaces plantées matures « noyaux ») :

\begin{center}
\begin{tabularx}{\linewidth}{|l|X|c|X|}
\hline
\rowcolor{popblueL}
\textbf{Société} & \textbf{Culture(s) principale(s)} & \textbf{Superf. plantée (ha)} & \textbf{Localisation principale (Région / Département)} \\
\hline
\textbf{CDC} (Cameroon Development Corporation) & Bananier, Palmier, Hévéa & \SI{\approx 40\,000}{ha} & Sud-Ouest (Fako, Meme) : Tiko, Muyuka, Mondoni, Lobe \\
\hline
\textbf{SOCAPALM} & Palmier à huile & \SI{\approx 58\,000}{ha} & Littoral (Sanaga-Maritime), Centre (Nyong-et-Kellé), Sud (Océan) \\
\hline
\textbf{HEVECAM} & Hévéa & \SI{\approx 38\,000}{ha} & Centre (Nyong-et-So'o, Méfou-et-Afamba) : Nkolfoulou, Obala, Mbandjock \\
\hline
\textbf{SODECOTON} (Noyaux / Fermes semencières) & Coton & \SI{\approx 15\,000}{ha} & Nord, Adamaoua, Extrême-Nord (Garoua, Maroua, Ngaoundéré) \\
\hline
\textbf{SOSUCAM} & Canne à sucre & \SI{\approx 10\,000}{ha} & Littoral (Moungo) : Nkoteng, Mbandjock \\
\hline
\textbf{SAFACAM} & Palmier, Hévéa & \SI{\approx 4\,000}{ha} & Littoral (Sanaga-Maritime) : Dizangué, Mbiako \\
\hline
\end{tabularx}
\end{center}

\begin{attention}[Spécificité de la CDC et de SODECOTON]
\begin{itemize}
    \item \textbf{CDC :} Entreprise publique (État 100 \%) en restructuration (crise anglophone depuis 2016). Historiquement 2e employeur public après l'État. Cultures diversifiées (banane dessert = 1er exportateur camerounais, palmier, hévéa). Siège à Bota (Limbe).
    \item \textbf{SODECOTON :} SEM (État \SI{60}{\%}, Wilmar \SI{40}{\%}). Son « noyau » (fermes semencières, centres de multiplication) est petit (\SI{\approx 15\,000}{ha}). Son emprise réelle est \textbf{diffuse} : elle encadre \SI{>200\,000}{ha} cultivés par des \textbf{planteurs villageois} (PV) sous contrat (intrants à crédit, encadrement technique, prix garanti). C'est de l'\cle{agriculture contractuelle}, pas de la plantation salariée classique.
\end{itemize}
\end{attention}

\courssub{Les plantations des particuliers (émergence d'une bourgeoisie agraire nationale)}
Depuis les années 1990 (lois de libéralisation foncière), des opérateurs économiques camerounais (commerçants, cadres, diaspora) acquièrent des titres fonciers et créent des plantations de taille moyenne (\SI{100}{--}{2\,000}{ha}).
\begin{itemize}
    \item \textbf{Cultures :} Palmier à huile (régions Littoral, Sud-Ouest, Centre), Hévéa (Centre, Sud), Cacao/Café (Ouest, Sud-Ouest - rénovation), Ananas (Moungo), Maïs/Soja (Adamaoua, Nord - fermes céréalières).
    \item \textbf{Exemples :} Plantations de la Mungo (palmier), Fermes de Mbandjock (canne/hévéa), nombreuses plantations individuelles autour de \textbf{Eseka, Dibombari, Muyuka, Ngaoundéré}.
    \item \textbf{Modèle :} Souvent « noyau + planteurs villageois » (outgrowers) pour alimenter une mini-huilerie ou un pressoir artisanal amélioré. Moins capitalisées que les SEM, elles dynamisent l'offre locale de régimes frais.
\end{itemize}

\courssub{Illustration par les données chiffrées : superficies et production (Noyaux)}

\begin{exoresolu}[Commentaire de document statistique : les superficies des principales plantations (estimations 2023)]
\textbf{Document :} Tableau récapitulatif des superficies plantées (hectares) des principales sociétés agro-industrielles (noyaux).

\begin{center}
\begin{tabularx}{\linewidth}{|l|X|c|c|}
\hline
\rowcolor{popblueL}
\textbf{Société} & \textbf{Culture principale} & \textbf{Superficie plantée (ha)} & \textbf{Part dans le total (\%)} \\
\hline
SOCAPALM & Palmier à huile & 58\,000 & 35,4 \\
CDC & Bananier / Palmier / Hévéa & 40\,000 & 24,4 \\
HEVECAM & Hévéa & 38\,000 & 23,2 \\
SODECOTON (Noyaux) & Coton & 15\,000 & 9,1 \\
SOSUCAM & Canne à sucre & 10\,000 & 6,1 \\
SAFACAM & Palmier / Hévéa & 4\,000 & 2,4 \\
\hline
\textbf{Total} & & \textbf{165\,000} & \textbf{100} \\
\hline
\end{tabularx}
\end{center}

\textbf{Consigne :} Construis un diagramme en barres représentant ces superficies et commente la hiérarchie.

\tcblower
\textbf{Solution détaillée :}

\textbf{1. Construction du diagramme (voir figure ci-dessous).}
\begin{itemize}
    \item Axe des abscisses : noms des sociétés (ordre décroissant de superficie).
    \item Axe des ordonnées : superficie en milliers d'hectares (échelle 0 à 65).
    \item Barres de largeur égale, couleur distincte par culture dominante.
    \item Titre précis, source, unité.
\end{itemize}

\textbf{2. Commentaire structuré :}
\begin{itemize}
    \item \textbf{Tripolarité :} Trois pôles dominent : Palmier (SOCAPALM + SAFACAM + part CDC \approx 50 \%), Hévéa (HEVECAM + part CDC \approx 30 \%), Banane (CDC \approx 15 \%). Le coton (noyau) et la canne restent marginaux en surface noyau.
    \item \textbf{Poids de la CDC :} 2e superficie noyau (\SI{40\,000}{ha}), mais 1er exportateur banane. Sa crise récente (2016-2023) a réduit l'activité effective.
    \item \textbf{Concentration :} 3 firmes (SOCAPALM, CDC, HEVECAM) = \SI{83}{\%} des surfaces noyau.
    \item \textbf{Nuance SODECOTON :} Sa faible surface noyau ne reflète pas son poids économique/social réel (encadrement \SI{>200\,000}{ha} villageois).
\end{itemize}
\textbf{Conclusion :} L'agro-industrie camerounaise est \textbf{très concentrée} et \textbf{spécialisée} sur trois pérennes d'exportation (palmier, hévéa, banane).
\end{exoresolu}

\begin{popfigure}
\begin{tikzpicture}
\begin{axis}[
    ybar,
    bar width=16pt,
    width=\linewidth,
    height=9cm,
    enlarge x limits=0.2,
    ylabel={Superficie (milliers d'hectares)},
    ylabel style={font=\small},
    xlabel={Société agro-industrielle},
    xlabel style={font=\small},
    xtick=data,
    xticklabels={SOCAPALM, CDC, HEVECAM, SODECOTON, SOSUCAM, SAFACAM},
    xticklabel style={rotate=30, anchor=east, font=\small},
    ymin=0, ymax=65,
    ymajorgrids=true,
    grid style={popline},
    legend style={at={(0.5,-0.28)}, anchor=north, legend columns=-1, font=\small, draw=popline},
    nodes near coords={\pgfmathprintnumber\pgfplotspointmeta},
    nodes near coords style={font=\tiny, anchor=south, yshift=2pt},
    title={Superficies plantées des principales grandes plantations - Noyaux (estimations 2023)},
    title style={font=\bfseries\small, yshift=5pt},
]
\addplot[fill=popgreen!70!popblue, draw=popgreen!50!popdark] coordinates {(1,58)};
\addplot[fill=popgreen!50!popgold, draw=popgreen!50!popdark] coordinates {(2,40)};
\addplot[fill=poporange!70!popdark, draw=poporange!50!popdark] coordinates {(3,38)};
\addplot[fill=popgold!80!poporange, draw=popgold!50!popdark] coordinates {(4,15)};
\addplot[fill=poppurple!70, draw=poppurple!50!popdark] coordinates {(5,10)};
\addplot[fill=poppink!70, draw=poppink!50!popdark] coordinates {(6,4)};
\legend{Palmier, Banane/Palmier/Hévéa, Hévéa, Coton, Canne, Palmier/Hévéa}
\end{axis}
\end{tikzpicture}
\legende{Diagramme en barres des superficies plantées des noyaux (en milliers d'hectares). Source : estimations compilées rapports d'activité, MINADER, SNI (2023). Ordres de grandeur.}
\end{popfigure}

\coursec{L'impact socio-économique des grandes plantations}

\courssub{Apports positifs : moteurs de l'intégration et du développement}

\begin{propriete}[Effets structurants positifs]
\begin{enumerate}
    \item \textbf{Emploi salarié massif :} \SI{\approx 60\,000}{--}{80\,000} emplois directs permanents (ouvriers + cadres) + saisonniers. 1er employeur privé formel du pays. Salaires \num{> SMIG} (\SI{36\,270}{FCFA} 2023), primes, logement, soins, scolarité gratuits.
    \item \textbf{Recettes fiscales et d'exportation :} \SI{> 500}{milliards FCFA/an} de CA cumulé. Droits d'exportation, IS, TVA, patentes. Huile de palme \approx \SI{180}{milliards}, Banane \approx \SI{150}{milliards}, Caoutchouc \approx \SI{100}{milliards}, Coton \approx \SI{80}{milliards}, Sucre \approx \SI{50}{milliards}.
    \item \textbf{Aménagement du territoire / Agropoles :} Création de pôles urbains ruraux (Mbongo, Tiko, Nkolfoulou, Mbandjock, Garoua) : routes, électricité, eau, santé, éducation, marchés. Désenclavement local.
    \item \textbf{Transfert de technologie / Encadrement :} Variétés améliorées, itinéraires techniques diffusés aux planteurs villageois (PV) via contrats (SOCAPALM, SODECOTON, CDC, SOSUCAM). Appui à la recherche (IRAD, stations propres).
    \item \textbf{Souveraineté alimentaire / Industrielle :} Couverture \SI{>80}{\%} besoins huile, \SI{100}{\%} sucre, \SI{>90}{\%} coton-fibre local. Base pour industries aval (savonneries, huileries artisanales, textile CICAM, agro-alimentaire).
\end{enumerate}
\end{propriete}

\courssub{Problèmes et limites : les revers de la médaille}

\begin{attention}[Griefs majeurs récurrents (sujets de dissertation / oral Bac)]
\begin{enumerate}
    \item \textbf{Accaparement foncier \& conflits coutumiers :} Concessions octroyées sans consentement libre, préalable et éclairé (CLPE) des communautés. Emprise sur terres de chasse, jachères, sites sacrés. Ex. : Conflits SOCAPAM/Kienké (Océan), HEVECAM/Nkolfoulou, CDC/Bakweri (Sud-Ouest). Procédures d'indemnisation jugées insuffisantes ou opaques.
    \item \textbf{Pollution \& dégradation environnementale :} Rejets industriels (effluents huileries/sucreries : DBO5 élevée, température, pH) dans cours d'eau (Sanaga, Mungo, Wouri). Épandages aériens (hévéa, banane) : dérive pesticides sur villages, eau de source. Déforestation initiale (conversion forêt \rightarrow monoculture), érosion sols nus (jeunes plantations), perte biodiversité (corridors rompus).
    \item \textbf{Précarité sociale \& conditions de travail :} Pénibilité (saignée hévéa nuit, port charges lourdes régimes), exposition pesticides (EPI souvent insuffisants), contrats journaliers/CDD renouvelés (éviter CDI), syndicats réprimés ou cooptés. Logements vétustes (camps « indigènes » vs camps « cadres »).
    \item \textbf{Enclavement relatif / Dépendance :} Villes d'entreprise mono-activités. Si usine ferme (crise CDC), effondrement local. Prix mondiaux volatils (caoutchouc, coton, huile) \rightarrow instabilité revenus État/ouvriers.
    \item \textbf{Rapatriement des bénéfices :} Majorité dividendes vers actionnaires étrangers (Bolloré, GMG, Wilmar, Somdiaa). Réinvestissement local limité (maintenance vs extension).
\end{enumerate}
\end{attention}

\courssub{Synthèse : Bilan contrasté}
Les grandes plantations sont le \textbf{squelette de l'économie d'exportation} et des \textbf{vecteurs puissants d'intégration nationale} (brassage populations, infrastructures). Mais leur modèle « concessionnaire colonial » hérité pose des problèmes aigus de \textbf{justice foncière}, \textbf{justice environnementale} et \textbf{justice sociale} non résolus. La réforme foncière (loi 2022) et la RSE (Responsabilité Sociétale des Entreprises) sont les chantiers actuels.

\begin{aretenir}[Fiche Bac — Impacts socio-économiques]
\textbf{+ Positifs :} Emplois formels (\SI{\approx 70\,000}{directs}), Recettes export/fiscales (\SI{>500}{Mds FCFA}), Agropoles (services, désenclavement), Transfert tech (PV), Souveraineté (huile, sucre, coton).
\textbf{- Négatifs :} Conflits fonciers (terres coutumières), Pollution (effluents, pesticides), Précarité ouvrière (CDD, EPI, camps ségrégés), Dépendance cours mondiaux, Fuite capitaux (dividendes).
\textbf{Mots-clés :} \cle{Accaparement foncier}, \cle{Agropole}, \cle{Agriculture contractuelle}, \cle{RSE}, \cle{CLPE}.
\end{aretenir}

\coursec{Dossier 2 : La main-d'œuvre dans les grandes plantations, vecteur d'intégration nationale}

\noindent \textbf{Situation déclenchante :} Au camp « Cameroun » de la SOCAPALM Mbongo, au réfectoire du midi, on entend du \emph{Foulfuldé} (Nord), du \emph{Bamiléké} (Ouest), du \emph{Bassa} (Littoral), du \emph{Ewondo} (Centre), du \emph{Pidgin English} (Sud-Ouest) et du \emph{Français}. Les ouvriers partagent le même plat de \emph{ndolé} ou \emph{koki}, vont à la même mosquée/église, leurs enfants jouent ensemble à l'école de la plantation. Mais le soir, les cadres rentrent dans leurs villas climatisées (camp « Europe »), les ouvriers dans leurs cases alignées (camp « Cameroun »).

\medskip\noindent \textbf{Problématique du dossier :} En quoi la main-d'œuvre des grandes plantations, par ses origines, son organisation spatiale et ses structures sociales, illustre-t-elle l'intégration nationale ? Quelles limites et tensions cette intégration recèle-t-elle ?

\courssub{Origine des travailleurs : un brassage national (et sous-régional)}

\begin{itemize}
    \item \textbf{Migration Nord \rightarrow Sud (historique, années 1950-1990) :} Main-d'œuvre « contrat » (contrat ONAREC/Office du Travail) recrutée au Nord (Adamaoua, Nord, Extrême-Nord) pour palmeries (SOCAPALM, SAFACAM), hévéa (HEVECAM), banane (CDC). \textbf{Flux massifs :} \num{>100\,000} travailleurs migrants sur 40 ans. Installation définitive : « \emph{les enfants de la plantation} » (2e, 3e génération).
    \item \textbf{Migration Ouest \rightarrow Sud-Ouest/Littoral :} Bamiléké, Bamoun vers CDC (Tiko, Muyuka), SOSUCAM (Mbandjock), SOCAPALM (Eseka). Commerçants, artisans, puis ouvriers qualifiés.
    \item \textbf{Main-d'œuvre locale (autochtones) :} Bassa, Bakoko, Duala, Ewondo, Bakweri. Souvent réticents au travail salarié pénible (préférence cultures propres, commerce). Majoritaires dans les camps « cadres » et emplois qualifiés.
    \item \textbf{Main-d'œuvre étrangère (CEMAC) :} Tchadiens, Centrafricains, Nigérians (pêcheurs/agriculteurs saisonniers), de plus en plus nombreux pour les tâches les plus dures (saignée, ramassage coton, coupe canne).
\end{itemize}

\begin{exemplebox}[Le « contrat ONAREC » (années 60-90)]
L'\emph{Office National de Recrutement de la Main-d'Œuvre} organisait le départ groupé de jeunes du Nord (Garoua, Maroua, Ngaoundéré) vers les plantations du Sud. Transport payé, contrat 2 ans renouvelable, rapatriement garanti. C'était une \textbf{politique d'État d'intégration par le travail}. Aujourd'hui, recrutement local direct ou sous-traitance.
\end{exemplebox}

\courssub{Organisation spatiale : la ségrégation résidentielle (camps cadres vs camps ouvriers)}

\begin{popfigure}
\begin{tikzpicture}[scale=0.9, every node/.style={font=\small}]
% Plan type camp plantation
% Usine centre
\draw[poporange, fill=poporangeL, thick] (-1, -0.5) rectangle (1, 1.5);
\node[align=center] at (0, 0.5) {USINE\\HUILERIE};
% Camp Cadres (Nord-Ouest - vent dominant)
\begin{scope}[xshift=-4cm, yshift=2cm]
\draw[popblue, fill=popblueL, thick] (-2.5, -1.5) rectangle (2.5, 1.5);
\foreach \x in {-2, -1, 0, 1, 2} {
    \foreach \y in {-1, 0, 1} {
        \draw[popblue, fill=white] (\x-0.3, \y-0.3) rectangle (\x+0.3, \y+0.3);
    }
}
\node[align=center, popblue] at (0, -2) {CAMP CADRES\\ (Villas, élec, eau, clim)};
\draw[->, popline] (0, -1.5) -- (0, -0.5);
\end{scope}
% Camp Ouvriers (Sud-Est)
\begin{scope}[xshift=4cm, yshift=-2.5cm]
\draw[popgreen, fill=popgreenL, thick] (-3, -2) rectangle (3, 2);
\foreach \x in {-2.5, -1.5, -0.5, 0.5, 1.5, 2.5} {
    \foreach \y in {-1.5, -0.5, 0.5, 1.5} {
        \draw[popgreen, fill=popmuted!30] (\x-0.4, \y-0.4) rectangle (\x+0.4, \y+0.4);
    }
}
\node[align=center, popgreen] at (0, -2.5) {CAMP OUVRIERS\\ (Cases alignées, collectifs)};
\draw[->, popline] (0, 2) -- (0, 1);
\end{scope}
% Village riverain (Ouest)
\begin{scope}[xshift=-6cm, yshift=-1cm]
\draw[poppurple, fill=poppurpleL, thick] (-2, -1.5) rectangle (2, 1.5);
\draw[poppurple, fill=poppurple!20] (-1.5, -1) rectangle (-0.5, 0); \node at (-1, -0.5) {Chefferie};
\draw[poppurple, fill=poppurple!20] (0.5, -1) rectangle (1.5, 0); \node at (1, -0.5) {École vill.};
\foreach \x in {-1.5, -0.5, 0.5, 1.5} {
    \foreach \y in {0.2, 1} {
        \draw[poppurple, fill=popmuted!30] (\x-0.3, \y-0.3) rectangle (\x+0.3, \y+0.3);
    }
}
\node[align=center, poppurple] at (0, -2) {VILLAGE RIVERAIN};
\end{scope}
% Flèche Nord
\draw[->, thick, popdark] (5, 3) -- (5, 3.7) node[above] {N};
% Légende interne
\node[align=left, font=\tiny, draw=popline, fill=popcream, anchor=north east] at (6, 3.5) {
\textbf{Légende :}\\
\color{poporange}\rule{2mm}{2mm} Usine / Transformation\\
\color{popblue}\rule{2mm}{2mm} Camp Cadres (haut standing)\\
\color{popgreen}\rule{2mm}{2mm} Camp Ouvriers (standard)\\
\color{poppurple}\rule{2mm}{2mm} Village riverain
};
\end{tikzpicture}
\legende{Croquis schématique : organisation spatiale type d'une agropole (ex. Mbongo). Séparation nette camps cadres / camps ouvriers (vent dominant, distance usine), village riverain en marge.}
\end{popfigure}

\begin{propriete}[Dualité spatiale structurelle]
\begin{itemize}
    \item \textbf{Camp Cadres (\"Europe\" / \"Résidentiel\") :} Ingénieurs, contremaîtres, administratifs. Villas individuelles, électricité 24/7, eau courante, climatisation, jardins, club house, tennis. Localisation : vent dominant (amont odeurs/bruit), vue sur plantation.
    \item \textbf{Camp Ouvriers (\"Cameroun\" / \"Indigènes\" / \"Quartiers\") :} Cueilleurs, saigneurs, manutentionnaires, chauffeurs. Cases alignées (\"longères\") en dur ou tôle, \SIrange{12}{18}{\meter\squared}, latrines collectives, bornes-fontaines, éclairage public récent. Localisation : vent dominant aval (proximité usine, bruit, odeurs effluents).
    \item \textbf{Ségrégation fonctionnelle :} Écoles distinctes (école cadres / école ouvriers), dispensaires distincts (médecin cadres / infirmier ouvriers), marchés distincts. Reproduction des inégalités sociales dans l'espace.
\end{itemize}
\end{propriete}

\courssub{Structures sociales : l'entreprise « État dans l'État » (paternalisme industriel)}

La plantation fournit un \textbf{package social complet} (logement, santé, éducation, loisirs, culte, rations alimentaires subventionnées) qui soude la communauté ouvrière mais la rend dépendante.

\begin{center}
\begin{tabularx}{\linewidth}{|l|X|X|}
\hline
\rowcolor{popblueL}
\textbf{Structure} & \textbf{Fonction d'intégration} & \textbf{Limites / Tensions} \\
\hline
\textbf{Écoles (Primaire + CEG)} & Scolarisation massive (taux \SI{>90}{\%}), mixité ethnique/religieuse, programme national. Ascenseur social (enfants ouvriers \rightarrow cadres). & Surcharge classes, double vacation, disparités moyens (école cadres mieux dotée). Grèves enseignants (salaires). \\
\hline
\textbf{Centres de santé (CSI / Hôpital)} & Soins gratuits travailleurs + familles + riverains (parfois). Vaccination, maternité, urgences. Réduction mortalité. & Pénurie médecins/médicaments, évacuations difficiles (Douala/Yaoundé loin), prise en charge accidents travail parfois contestée. \\
\hline
\textbf{Logement / Voirie / Eau / Élec} & Cadre de vie stable, urbanisation rurale. Accès électricité/eau \SI{>}{moyenne nationale rurale}. & Vétusté réseaux (fuites, coupures), surpeuplement cases (familles élargies), insalubrité (déchets, eaux usées). \\
\hline
\textbf{Commerces / Marchés / Cultes} & Vie sociale intense, mixité (ouvriers, cadres, villageois, commerçants ambulants). Mosquées, églises, temples côte à côte. & Économie informelle dominante, endettement ouvriers (crédits « tontines », usure), tensions interethniques latentes (ex. éleveurs vs cultivateurs). \\
\hline
\textbf{Syndicats / Comités d'entreprise} & Négociation salaires/primes, activités culturelles/sportives (foot inter-camps). & Syndicats souvent « jaunes » (pro-direction), répression grèves (licenciements, expulsion camps), division ethnique syndicats. \\
\hline
\end{tabularx}
\end{center}

\begin{savaistu}[Le football : ciment de l'intégration au camp]
Chaque grande plantation a son équipe de football (ex. \textbf{Canon de Mbongo}, \textbf{Panther de Nkolfoulou}, \textbf{Coton Sport de Garoua} né de SODECOTON). Le stade du camp est le lieu unique où cadres, ouvriers, villageois, Nordistes, Sudistes, étrangers se mélangent sans hiérarchie le dimanche. C'est un \textbf{rituel d'intégration par le sport} étudié par les sociologues (ex. travaux sur Coton Sport).
\end{savaistu}

\courssub{Bilan du dossier : Intégration réelle mais inachevée et conflictuelle}

\begin{aretenir}[Fiche Bac — Dossier 2 : Main-d'œuvre \& Intégration]
\textbf{1. Origines :} Brassage Nord/Sud/Ouest/Étranger \rightarrow \textbf{Creuset national}. 2e/3e générations = « Camerounais de la plantation » (identité hybride).
\textbf{2. Organisation spatiale :} \textbf{Ségrégation résidentielle} (Camps Cadres vs Ouvriers) = reproduction inégalités. Village riverain en marge (conflits fonciers).
\textbf{3. Structures sociales :} \textbf{Paternalisme industriel} (logement, santé, école, culte gratuits) \rightarrow cohésion sociale, dépendance, contrôle.
\textbf{4. Intégration :} \textbf{Réelle} (mariages mixtes, langue commune pidgin/français, culture partagée, football), mais \textbf{inachevée} (ségrégation habitat, précarité, conflits fonciers/travail, syndicats faibles).
\textbf{Notion clé :} \cle{Intégration par le bas} (vécu quotidien) vs \cle{Intégration par le haut} (politique État/Entreprise).
\end{aretenir}

\begin{exercice}[Sujet type Bac - Dossier 2]
\textbf{Document 1 :} Extrait d'entretien avec un ouvrier de 3e génération à Mbongo (2022) : « Mon grand-père est venu de Maroua en 1972. Mon père est né ici. Moi je suis né ici. Je suis Bassa maintenant, je parle bassa, je mange ndolé. Mais ma terre, c'est au Nord. Ici, on est chez SOCAPALM. Si l'usine ferme, on fait quoi ? »
\textbf{Document 2 :} Carte schématique du camp (cf. croquis ci-dessus).
\textbf{Consigne :} En vous appuyant sur les documents et vos connaissances, montrez que la main-d'œuvre des grandes plantations est un vecteur d'intégration nationale, tout en relevant les limites de cette intégration.
\end{exercice}

\begin{corrige}[Exercice n°1 - Dossier 2]
\textbf{Introduction :} Accroche (brassage), définition intégration nationale, présentation docs, problématique, annonce plan.
\textbf{I. Un puissant vecteur d'intégration nationale (le « creuset ») :}
\begin{itemize}
    \item A. Origines diverses \rightarrow brassage ethnique, linguistique, religieux (Doc 1 : grand-père Nord, père/fils nés sur place, acculturation Bassa).
    \item B. Structures sociales communes \rightarrow école unique (mixité), santé, culte, marché, football (Doc 2 : stade central). Création identité commune « enfants de la plantation ».
    \item C. Mobilité sociale \rightarrow fils ouvriers deviennent cadres/instituteurs/infirmiers (ascenseur).
\end{itemize}
\textbf{II. Une intégration limitée, inégalitaire et conflictuelle :}
\begin{itemize}
    \item A. Ségrégation spatiale persistante (Doc 2 : camps cadres/ouvriers distincts, inégalités habitat/services). Reproduction classes sociales.
    \item B. Précarité statutaire \rightarrow CDD, sous-traitance, bas salaires, pénibilité, syndicats muselés \rightarrow sentiment d'exploitation, pas de citoyenneté pleine.
    \item C. Conflits fonciers avec riverains (Doc 2 : village en marge) \rightarrow « autochtones » vs « allogènes » (ouvriers), menaces expulsion, insécurité foncière.
    \item D. Dépendance à l'entreprise (paternalisme) \rightarrow pas d'autonomie politique/économique locale.
\end{itemize}
\textbf{Conclusion :} Intégration vécue (quotidien, culture) vs intégration institutionnelle (droits, foncier, pouvoir) inachevée. Avenir lié à réforme foncière, RSE, diversification économique locale.
\end{corrige}

\coursec{Travaux pratiques : Cartographier et représenter les grandes plantations}

\noindent Les TP de cette leçon portent sur deux compétences clés du programme : \textbf{localiser} (carte de localisation) et \textbf{représenter l'organisation spatiale} (croquis d'agropole).

\courssub{TP 1 : Carte de localisation des grandes plantations du Cameroun (1 h)}

\begin{methode}[Réaliser une carte de localisation au Bac (3-4 pts)]
\begin{enumerate}
    \item \textbf{Fond de carte :} Contour du Cameroun (simplifié), limites régionales, capitale Yaoundé, métropole Douala, port de Kribi, grands fleuves (Sanaga, Wouri, Bénoué, Logone).
    \item \textbf{Localisation des 6 firmes majeures :} Point (●) ou surface (○ proportionnelle) au siège / cœur de plantation. Nom de la firme + culture(s) en légende.
    \item \textbf{Flux d'exportation :} Flèches (→) des zones de production vers Douala/Kribi (port). Épaisseur proportionnelle au tonnage/valeur.
    \item \textbf{Légende organisée :} 3 rubriques : \textit{Implantations}, \textit{Réseaux}, \textit{Flux}. Hiérarchie visuelle (taille, couleur).
    \item \textbf{Titre complet :} « Le Cameroun : localisation des grandes plantations agro-industrielles et flux d'exportation (2023) ».
\end{enumerate}
\end{methode}

\begin{experience}[TP 1 guidé : Construction de la carte]
\textbf{Matériel :} Fond de carte muet A4 (contour + régions + fleuves + Yaoundé/Douala/Kribi), crayons de couleur, règle, compas.
\textbf{Données à reporter (coordonnées approximatives Lambert / repères visuels) :}
\begin{itemize}
    \item \textbf{CDC} : Sud-Ouest, axe Tiko-Limbé-Muyuka-Mondoni (Banane, Palmier, Hévéa). ● Vert/Jaune/Brun.
    \item \textbf{SOCAPALM} : 5 sites. ● Vert : Mbongo/Dibombari (Littoral, près Édéa), Eseka (Centre, RN1), Kienké/Mbonjo (Sud, près Kribi/Campo).
    \item \textbf{HEVECAM} : Centre, Nord Yaoundé. ● Brun : Nkolfoulou (RN1), Obala, Mbandjock.
    \item \textbf{SOSUCAM} : Littoral, Moungo. ● Violet : Nkoteng / Mbandjock (RN1, ligne chemin de fer).
    \item \textbf{SAFACAM} : Littoral, Sanaga-Maritime. ● Vert/Brun : Dizangué / Mbiako (bord Sanaga).
    \item \textbf{SODECOTON (Noyaux)} : Nord/Adamaoua. ● Blanc : Garoua (Ladje), Maroua, Ngaoundéré (fermes semencières).
\end{itemize}
\textbf{Flux :} Flèches grosses vers Douala (Banane, Palmier, Hévéa, Coton, Sucre). Flèche vers Kribi (Palmier Sud, Bois, Minerai). Flèche Nord \rightarrow Douala (Coton fibre).
\textbf{Attendu :} Carte propre, légende encadrée, titre, orientation, échelle graphique, source.
\end{experience}

\begin{popfigure}
\begin{center}
\includegraphics[width=\linewidth,height=9.5cm,keepaspectratio]{N01/cmr_map.jpg}
\end{center}
\legende{Fond de carte : carte générale du Cameroun (TP 1). Le croquis attendu situe les noyaux de plantation (Sud-Ouest autour de Buea, Littoral autour de Douala et d'Édéa, Centre au nord de Yaoundé, Sud, Nord autour de Garoua et de Maroua) et trace en flèches épaisses l'orientation des flux vers les ports d'exportation de Douala et de Kribi. \\ Source : Wikimedia Commons, CIA, domaine public.}
\end{popfigure}

\courssub{TP 2 : Croquis de l'organisation spatiale d'une agropole (ex. Mbongo / SOCAPALM) (1 h)}

\begin{methode}[Réaliser un croquis d'agropole au Bac (4-5 pts)]
\begin{enumerate}
    \item \textbf{Contour \& Cadre :} Rectangle arrondi (concession). Flèche Nord (haut droite). Échelle graphique (ex. 1 cm = 2 km). Titre complet en haut.
    \item \textbf{Zonage interne (3 zones minimum, hachurées/colorées) :}
        \begin{itemize}
            \item \textit{Zone de production} (vert) : Quadrillage parcelles + symbole culture (● palmier, ▲ hévéa, ■ canne). Légende : « Palmeraie industrielle (monoculture intensive) ».
            \item \textit{Zone industrielle} (gris/rouge) : Usine (huilerie), hangars, citernes, quai de chargement. Symbole ⚙ ou 🏭.
            \item \textit{Zone d'habitat et de services} (jaune/rose/bleu) : Camp Cadres (villas, dispersé), Camp Ouvriers (aligné, dense), École (🏫), Dispensaire (⛑), Marché (🛒), Réseau eau/élec.
        \end{itemize}
    \item \textbf{Réseaux de communication :} Pistes principales plantation (traits pleins épais), Route nationale (double trait), Voie ferrée si existante (SOSUCAM : trait pointillé double), Cours d'eau (bleu).
    \item \textbf{Flux internes et externes :} Flèches fines vertes : « Régimes frais » (champ → usine). Flèches rouges : « Huile brute » (usine → route/rail → port). Flèches bleues : « Eau / Énergie » (fleuve/barrage → usine/camps).
    \item \textbf{Légende structurée (4 rubriques) :} 1. Occupation du sol, 2. Infrastructures, 3. Réseaux, 4. Flux. Ordre logique, symboles cohérents avec carte.
\end{enumerate}
\end{methode}

\begin{experience}[TP 2 guidé : Croquis de l'agropole de Mbongo (SOCAPALM)]
\textbf{Consigne :} À partir de la description ci-dessous, réalise le croquis légendé de l'agropole de Mbongo (Sanaga-Maritime) sur une feuille A4 paysage.
\textbf{Description de base (données terrain simplifiées) :}
\begin{itemize}
    \item Concession : \SI{12\,000}{ha} (rectangle \SI{12}{km} x \SI{10}{km} env.). Bordée au Nord par la Sanaga.
    \item Production : \SI{10\,000}{ha} palmiers (quadrillage \SI{100}{ha}, pistes \SI{10}{m}). \SI{2\,000}{ha} réserves/non planté.
    \item Industrie : Huilerie \SI{60}{t/h} au centre-Ouest, près piste principale. Citerne stockage. Quai de chargement camions.
    \item Habitat/Services :
        \begin{itemize}
            \item Camp Cadres (Nord-Ouest, vent dominant) : \SI{50}{villas}, club, tennis, école privée, dispensaire médecin.
            \item Camp Ouvriers (Sud-Est) : \SI{500}{cases} alignées (rues orthogonales), 2 écoles publiques, 1 CSI (infirmier), 1 marché, 1 stade, 1 mosquée, 1 église.
        \end{itemize}
    \item Réseaux : RN3 traverse au Sud (Douala-Édéa). Piste principale N-S relie camps/usine/RN3. Électricité (ligne HT depuis Édéa). Eau (forages + pompage Sanaga traité).
    \item Flux : Régimes (champs → usine). Huile (usine → RN3 → Douala/Kribi). Ouvriers (camps → champs). Village riverain Dibamba (Ouest, hors concession, conflit foncier).
\end{itemize}
\textbf{Attendu :} Croquis appliquant la méthode ci-dessus. Légende en 4 rubriques. Titre : « Croquis de l'agropole de Mbongo (SOCAPALM) : organisation spatiale et flux ».
\end{experience}

\begin{popfigure}
\begin{tikzpicture}[scale=0.7, every node/.style={font=\small}]
% Contour concession
\draw[popdark, thick, rounded corners=10pt] (-6,-4) rectangle (6,4);
% Sanaga (Nord)
\draw[popblue, thick, decorate, decoration={snake, amplitude=1pt, segment length=10pt}] (-6.5, 4) -- (6.5, 4) node[above, midway, popblue] {Sanaga};
% Quadrillage palmeraie (zone production)
\foreach \x in {-5.5,-3.5,-1.5,0.5,2.5,4.5} {
    \foreach \y in {-3.5,-1.5,0.5,2.5} {
        \draw[fill=popgreen!50!popblue, draw=popgreen!30!popdark, opacity=0.6] (\x,\y) rectangle (\x+1.8, \y+1.8);
        % Symboles palmiers
        \foreach \px in {0.3,0.9,1.5} {
            \foreach \py in {0.3,0.9,1.5} {
                \fill[popgreen!30!popdark] (\x+\px, \y+\py) circle (1pt);
            }
        }
    }
}
% Pistes principales
\draw[popmuted, line width=2.5pt] (-6,0) -- (6,0); % E-W centrale
\draw[popmuted, line width=2.5pt] (0,-4) -- (0,4); % N-S centrale
% Usine (Zone industrielle)
\draw[poporange, fill=poporangeL, thick] (-1, -1) rectangle (2, 1);
\node[align=center] at (0.5, 0) {HUILERIE\\60 t/h};
\draw[poporange, fill=poporangeL] (2, -0.3) rectangle (2.5, 0.3); % Quai
\node[anchor=west, font=\tiny] at (2.5, 0) {Quai charg.};
% Camp Cadres (Nord-Ouest)
\begin{scope}[xshift=-4cm, yshift=1.5cm]
\draw[popblue, fill=popblueL, thick] (-1.5, -1) rectangle (1.5, 1);
\foreach \x in {-1, 0, 1} {
    \foreach \y in {-0.5, 0.5} {
        \draw[popblue, fill=white, thick] (\x-0.3, \y-0.3) rectangle (\x+0.3, \y+0.3);
    }
}
\node[popblue, align=center] at (0, -1.3) {Camp Cadres};
\node[popblue, font=\tiny] at (0, 1.3) {🏫 🏥 ⛳};
\end{scope}
% Camp Ouvriers (Sud-Est)
\begin{scope}[xshift=3cm, yshift=-2.5cm]
\draw[popgreen, fill=popgreenL, thick] (-2, -1.5) rectangle (2, 1.5);
\foreach \x in {-1.5, -0.5, 0.5, 1.5} {
    \foreach \y in {-1, 0, 1} {
        \draw[popgreen, fill=popmuted!30] (\x-0.35, \y-0.35) rectangle (\x+0.35, \y+0.35);
    }
}
\node[popgreen, align=center] at (0, -1.8) {Camp Ouvriers};
\node[popgreen, font=\tiny] at (0, 1.8) {🏫 🏫 ⛑ 🛒 ⚽};
\end{scope}
% Village riverain (Ouest)
\begin{scope}[xshift=-7cm]
\draw[poppurple, fill=poppurpleL, thick, dashed] (-1.5, -1.5) rectangle (1.5, 1.5);
\node[poppurple, align=center] at (0, 0) {Village\\Dibamba};
\node[poppurple, font=\tiny] at (0, -1.8) {Conflit foncier};
\end{scope}
% RN3 (Sud)
\draw[popdark, double distance=3pt, double=popgray] (-6.5, -4.5) -- (6.5, -4.5) node[below, midway] {RN3 Douala - Édéa};
% Flèches flux
\draw[->, popgreen!70!popdark, thick] (0, -2) -- (0, -0.8) node[midway, right, font=\tiny] {Régimes};
\draw[->, poporange, very thick] (2, 0) -- (4, 0) node[midway, above, font=\small] {Huile brute};
\draw[->, poporange, very thick] (4, 0) -- (4, -4.5) node[midway, right, font=\tiny] {→ Douala/Kribi};
% Nord
\draw[->, thick, popdark] (5, 3.5) -- (5, 4.2) node[above] {N};
% Échelle
\draw[|-|, thick] (-5.5, -4.5) -- (-3.5, -4.5) node[midway, below] {2 km};
% Titre
\node[align=center, font=\bfseries\small, popdark] at (0, 5) {Croquis de l'agropole de Mbongo (SOCAPALM) : organisation spatiale et flux};
\end{tikzpicture}
\legende{Croquis type attendu pour le TP 2 (et Bac). Respect de la méthode : contour, zonage 3 zones, réseaux, flux, légende structurée, Nord, échelle, titre.}
\end{popfigure}

\courssub{Synthèse générale de la leçon}

\begin{aretenir}[Fiche Bac — Leçon 7 complète : L'agriculture de grandes plantations]
\textbf{1. Définition :} Agriculture commerciale, intensive, monoculture pérenne, vastes concessions (\SI{>500}{ha}), capitalistique, salariée, intégrée (usine), export/industrie.
\textbf{2. Types :}
\begin{itemize}
    \item \textbf{Grandes firmes (SEM) :} CDC (Banane/Palmier/Hévéa, Sud-Ouest), SOCAPALM (Palmier, Littoral/Centre/Sud), HEVECAM (Hévéa, Centre), SODECOTON (Coton, Nord - noyaux + encadrement villageois massif), SOSUCAM (Canne, Littoral), SAFACAM (Palmier/Hévéa, Littoral). \SI{\approx 165\,000}{ha} noyaux.
    \item \textbf{Particuliers :} Bourgeoisie agraire nationale, \SI{100}{--}{2\,000}{ha}, palmier/hévéa/céréales, modèle « noyau + PV ».
\end{itemize}
\textbf{3. Agropole :} Modèle territorial intégré (Production + Transformation + Services urbains). Ex. Mbongo, Tiko, Nkolfoulou, Mbandjock, Garoua.
\textbf{4. Impacts :} \textbf{+} Emplois (\SI{\approx 70\,000}{directs}), Recettes (\SI{>500}{Mds FCFA}), Aménagement, Souveraineté (huile, sucre, coton), Transfert tech. \textbf{-} Accaparement foncier, Pollution (effluents, pesticides), Précarité/Ségrégation camps, Dépendance cours mondiaux, Fuite capitaux.
\textbf{5. Dossier 2 - Main-d'œuvre :} Brassage national (Nord→Sud, Ouest→Sud, Étrangers) = \textbf{intégration vécue} (écoles, santé, football, mariages mixtes, pidgin). Mais \textbf{ségrégation spatiale} (camps cadres/ouvriers), \textbf{paternalisme}, \textbf{conflits fonciers} riverains, \textbf{syndicats faibles}. Intégration réelle mais inachevée.
\textbf{6. TP :} Carte localisation (6 firmes + flux ports) + Croquis agropole (zonage, réseaux, flux, légende 4 rubriques).
\end{aretenir}

\coursec{Les types de grandes plantations}

Après avoir défini l'agriculture de grandes plantations et en avoir décrit les caractéristiques générales (superficies, méthodes de travail, financement), nous allons maintenant distinguer les deux grandes catégories qui structurent ce secteur au Cameroun : les plantations détenues par de grandes firmes et celles qui appartiennent à des particuliers. Cette distinction est essentielle pour comprendre l'organisation de l'espace rural, les flux de main-d'œuvre et les enjeux d'intégration nationale.

\courssub{Les plantations des grandes firmes}

Les plantations des grandes firmes sont des exploitations agro-industrielles à capital important, souvent créées à l'époque coloniale ou dans les premières décennies de l'indépendance. Elles sont gérées par des sociétés d'État, des entreprises privées ou des structures mixtes (capitaux publics et privés). Leur taille dépasse généralement plusieurs milliers d'hectares, elles disposent d'un encadrement technique permanent, d'usines de transformation sur place ou à proximité, et leurs productions sont orientées vers l'exportation ou vers le marché national organisé.

\begin{definition}[Plantation de grande firme]
Une plantation de grande firme est une exploitation agricole de grande superficie (souvent > 1\,000 ha), capitalisée, gérée par une société structurée (publique, privée ou mixte), qui pratique une agriculture intensive, mécanisée et encadrée, avec transformation industrielle et commercialisation organisée.
\end{definition}

\courssub{Présentation des six grandes firmes et de leurs cultures}

Le tableau ci-dessous synthétise les principales firmes, leur localisation régionale, les cultures qu'elles développent, leurs superficies approximatives et la nature de leurs capitaux. Il constitue un outil de révision indispensable pour le Bac.

\begin{popfigure}
\begin{center}
\includegraphics[width=\linewidth,height=9.5cm,keepaspectratio]{N01/cmr_map.jpg}
\end{center}
\legende{Fond de carte : carte générale du Cameroun (régions, villes, routes, chemin de fer, barre d'échelle). Les six grandes firmes agro-industrielles (CDC, SOCAPALM, PHP, HEVECAM, SOSUCAM, SODECOTON) se localisent d'après le tableau ci-dessous : cinq sont dans le Sud forestier (Sud-Ouest, Littoral, Centre, Sud), une dans le Nord (coton, région de Garoua). \\ Source : Wikimedia Commons, CIA, domaine public.}
\end{popfigure}

\begin{center}
\begin{tabularx}{\linewidth}{|l|X|X|X|l|}
\hline
\rowcolor{popblueL}
\textbf{Firme} & \textbf{Localisation (Région, Villes clés)} & \textbf{Cultures développées} & \textbf{Superficie approx.} & \textbf{Type de capitaux} \\ \hline
CDC (Cameroon Development Corporation) & Sud-Ouest (Buea, Limbé, Tiko, Mbonge) & Banane, Palmier à huile, Hévéa & $\sim$ 40\,000 ha & Public (État) \\ \hline
SOCAPALM (Société Camerounaise de Palmeraies) & Littoral (Édéa, Dibombari, Mbongo, Apouh) & Palmier à huile & $\sim$ 35\,000 ha & Mixte (État + privé \emph{Siat}) \\ \hline
SAFACAM (Société Africaine Forestière et Agricole du Cameroun) & Littoral (Souza, Mbonjo, près de Douala) & Banane, Palmier à huile & $\sim$ 8\,000 ha & Privé (capitaux étrangers, groupe \emph{Socfin}) \\ \hline
SODECOTON (Société de Développement du Coton) & Nord, Adamaoua, Extrême-Nord (Garoua, Maroua, Ngaoundéré) & Coton & $\sim$ 15\,000 ha (propres) + encadrement paysan & Public (État) \\ \hline
SOSUCAM (Société Sucrière du Cameroun) & Centre (Nkoteng, Mbandjock) & Canne à sucre & $\sim$ 10\,000 ha & Mixte (État + privé \emph{Somdiaa}) \\ \hline
HEVECAM (Hévéa du Cameroun) & Sud (Kribi, Campo, Niete) & Hévéa (caoutchouc naturel) & $\sim$ 40\,000 ha & Privé (capitaux internationaux, \emph{Corrie MacColl / Siat}) \\ \hline
\end{tabularx}
\end{center}

\courssub{Les plantations des particuliers}

À côté de ces géants, existent des plantations appartenant à des planteurs camerounais (élites, cadres, commerçants, parfois des coopératives ou groupements d'intérêt économique). Elles portent sur les mêmes cultures d'exportation (café, cacao, palmier à huile, hévéa) mais leurs superficies sont plus modestes (généralement entre 10 et 500 ha). Elles sont gérées de façon plus familiale, avec un recours plus important à la main-d'œuvre locale et moins de mécanisation. Leur production est souvent vendue aux firmes transformatrices ou exportée via des circuits organisés.

\begin{definition}[Plantation de particulier]
Une plantation de particulier est une exploitation agricole commerciale appartenant à une personne physique ou à un groupe familial, de taille inférieure à celle des grandes firmes (généralement 10--500 ha), qui pratique une culture pérenne d'exportation (café, cacao, palmier, hévéa) avec un encadrement technique moins intensif et une main-d'œuvre souvent locale.
\end{definition}

\courssub{Comparaison des deux types}

Le tableau suivant met en évidence les différences structurelles entre les deux modèles. Il permet de répondre aux questions de type « comparer » au Bac.

\begin{center}
\begin{tabularx}{\linewidth}{|l|X|X|}
\hline
\rowcolor{popblueL}
\textbf{Critère} & \textbf{Plantations des grandes firmes} & \textbf{Plantations des particuliers} \\ \hline
Superficie moyenne & > 1\,000 ha (souvent 5\,000--50\,000 ha) & 10--500 ha \\ \hline
Financement & Capitaux publics, privés ou mixtes ; accès au crédit bancaire international & Capitaux propres, crédit agricole local (CFC, banques), parfois appui de l'État (PAD, PEA-Jeunes) \\ \hline
Encadrement technique & Ingénieurs agronomes, techniciens permanents, recherche variétale propre (IRAD, centres internes) & Conseillers agricoles ponctuels (MINADER, ONG), formation sur le tas, appui projets \\ \hline
Mécanisation & Élevée (tracteurs, avions épandeurs, usines de transformation sur site) & Faible à moyenne (tracteurs loués, travail manuel prédominant) \\ \hline
Main-d'œuvre & Salariés permanents + saisonniers (camps d'ouvriers) ; origine diverse (facteur d'intégration) & Familiale + journaliers locaux ; moins de brassage ethnique \\ \hline
Débouchés & Exportation directe, transformation industrielle intégrée, contrats d'État & Vente aux firmes transformatrices, coopératives, marchés locaux, filières certifiées \\ \hline
Impact environnemental & Risque pollution (intrants), déforestation initiale, mais plans de gestion certifiés (RSPO, ISO 14001) & Moins d'intrants chimiques, mais pression sur les terres communautaires et forêts secondaires \\ \hline
\end{tabularx}
\end{center}

\begin{methode}[Localiser une plantation sur une carte muette (Croquis Bac)]
\textbf{Étape 1 :} Repérer la région administrative concernée (ex. : Sud-Ouest pour la CDC). \\
\textbf{Étape 2 :} Choisir un figuré de surface (carré, cercle, triangle) et une couleur distincte pour chaque firme. \\
\textbf{Étape 3 :} Placer le symbole au centre approximatif de la zone de production (ville principale ou chef-lieu de la plantation). \\
\textbf{Étape 4 :} Rédiger une légende ordonnée : titre de la carte, échelle indicative, orientation (flèche Nord), liste des figurés avec leur signification (firme + culture principale). \\
\textbf{Étape 5 :} Vérifier la lisibilité : pas de chevauchement, symboles assez grands (min 4--5 mm), contraste des couleurs, trait soigné.
\end{methode}

\begin{exemplebox}[Exemple chiffré : la SODECOTON, un modèle hybride unique]
La SODECOTON encadre environ 250\,000 producteurs de coton dans le Grand-Nord (Nord, Adamaoua, Extrême-Nord). Elle gère aussi ses propres plantations industrielles (environ 15\,000 ha) et dispose de 5 usines d'égrenage (Garoua, Maroua, Kaélé, Guider, Figuil). En campagne 2022/2023, la production nationale de coton-graine a atteint près de 300\,000 tonnes, dont plus de 80 \% issues de l'encadrement paysan (culture contractuelle). Ce modèle hybride (plantation propre + culture contractuelle paysanne) est unique en Afrique centrale.
\end{exemplebox}

\begin{attention}[Piège fréquent : SODECOTON $\neq$ plantation unique]
Ne pas confondre la SODECOTON avec une seule plantation géante. C'est une société de développement qui combine : (1) des plantations industrielles (blocs de 100--500 ha), (2) un vaste réseau de paysans cotonniers encadrés par des contrats (intrants à crédit, conseil technique, prix garanti à la récolte), (3) la transformation (égrenage, huilerie de coton). Au Bac, précisez toujours « encadrement paysan » et « culture contractuelle » pour montrer que vous maîtrisez la nuance.
\end{attention}

\begin{savaistu}[La CDC, géant historique et vecteur d'urbanisation]
La Cameroon Development Corporation (CDC) a été créée en 1947 sous tutelle britannique (Commonwealth Development Corporation), puis reprise par l'État camerounais en 1967 (décret n°67/231). Avec plus de 40\,000 ha exploités (bananes, palmiers, hévéas) et un effectif nominal d'environ 20\,000 employés permanents (chiffre d'avant-crise anglophone), elle est l'un des plus grands employeurs du secteur formel. Ses camps (Tiko, Bota, Mbonge, Mukonje) sont de véritables villes intégrées avec écoles, hôpitaux, marchés, réseaux d'eau/électricité, illustrant le rôle intégrateur et structurant des plantations.
\end{savaistu}

\begin{savaistu}[L'agropole : le cas de SOSUCAM à Mbandjock/Nkoteng]
Le concept d'\textbf{agropole} (pôle agro-industriel intégré) s'illustre parfaitement avec SOSUCAM. Autour de l'usine sucrière, se sont développés : la plantation de canne, l'habitat ouvrier et cadre, les services sociaux (école, hôpital, marché), mais aussi des activités annexes (boulangerie, élevage, cultures vivrières des ouvriers). C'est un espace polarisé par l'usine, générateur d'intégration locale et de peuplement.
\end{savaistu}

\begin{exoresolu}[Exercice : Identifier la firme productrice de caoutchouc dans le Sud]
\textbf{Énoncé :} À l'aide de la carte schématique ci-dessus, nommez la firme qui produit de l'hévéa (caoutchouc naturel) dans la région du Sud et indiquez sa localisation précise. \\
\tcblower
\textbf{Solution :} La firme est \textbf{HEVECAM}. Elle est implantée dans la région du \textbf{Sud}, autour de \textbf{Kribi}, \textbf{Campo} et \textbf{Niete}, près de la côte atlantique. Sur la carte, son symbole (carré doré) se trouve au sud-ouest du contour, au voisinage de la ville de Kribi. La culture principale est l'hévéa, dont le latex est transformé en caoutchouc brut (TSR, RSS) pour l'exportation vers l'Europe et l'Asie.
\end{exoresolu}

\begin{exercice}[Entraînement : Construire un croquis de localisation (Type Bac)]
À partir du tableau de synthèse, réalisez un croquis légendé (format A4) représentant le Cameroun avec :
\begin{itemize}
\item Le contour du pays, les 10 régions (traits fins, noms en minuscules).
\item Les 6 firmes localisées par un figuré de surface (carré, cercle, triangle) de couleur distincte.
\item Une légende complète structurée (titre, échelle, orientation, figurés classés par firme avec culture).
\item Une flèche bleue indiquant le principal port d'exportation (Douala) et une flèche rouge pour le port en eau profonde de Kribi.
\end{itemize}
\end{exercice}

\begin{corrige}[Exercice n°1 : Croquis de localisation des grandes plantations]
Le croquis doit présenter :
\begin{enumerate}
\item Un contour simplifié mais reconnaissable du Cameroun (forme « triangle » allongé Nord-Sud).
\item Les 10 régions esquissées par des traits pointillés fins et nommées (Extrême-Nord, Nord, Adamaoua, Est, Centre, Sud, Littoral, Sud-Ouest, Nord-Ouest, Ouest).
\item Six figurés de surface placés aux coordonnées approximatives du tableau :
      \begin{itemize}
      \item CDC : Sud-Ouest (zone Buea/Tiko/Limbé) -- Carré vert.
      \item SOCAPALM : Littoral (zone Edéa/Dibombari) -- Carré orange.
      \item SAFACAM : Littoral (zone Souza, Ouest Douala) -- Carré violet.
      \item SODECOTON : Grand-Nord (Nord, Adamaoua, Extrême-Nord) -- Carré bleu (grand ou multiple).
      \item SOSUCAM : Centre (Nkoteng/Mbandjock, Nord Yaoundé) -- Carré rose.
      \item HEVECAM : Sud (Kribi/Campo/Niete) -- Carré or.
      \end{itemize}
\item Une légende en encadré (bas à droite ou gauche) en 3 colonnes : \textbf{Figuré} | \textbf{Firme} | \textbf{Culture principale}.
\item Échelle graphique (ex. 1 cm = 200 km) et flèche Nord.
\item Flèche bleue vers Douala (Wouri) : « Port de Douala - Exportation banane, coton, sucre, bois » ; Flèche rouge vers Kribi : « Port de Kribi - Exportation bois, hévéa, gaz (GNL) ».
\end{enumerate}
\textbf{Critères de réussite :} Propreté du trait, lisibilité des noms, légende complète et ordonnée, respect de l'échelle indicative.
\end{corrige}

\coursec{L'impact socio-économique des grandes plantations}

Après avoir identifié les acteurs majeurs que sont la CDC, la SOCAPALM, la SAFACAM, la SODECOTON, la SOSUCAM et l'HEVECAM, ainsi que les plantations particulières, nous devons maintenant évaluer ce que ces vastes domaines représentent concrètement pour le Cameroun. Derrière les hectares de palmiers à huile, d'hévéas ou de bananiers se joue une réalité complexe : ces plantations sont à la fois des moteurs puissants de l'économie nationale et des sources de tensions foncières et environnementales. Comprendre cette dualité, c'est saisir l'un des enjeux centraux de l'intégration nationale.

\courssub{Des retombées positives majeures pour l'économie et les populations}

\begin{exemplebox}[Le poumon économique de la région du Sud-Ouest : le cas de la CDC]
Imagine-toi au pied du mont Cameroun, à Tiko ou à Limbé. L'air sent l'embrun et la banane mûre. La \cle{CDC} (Cameroon Development Corporation), créée en 1947, est le plus grand employeur privé du pays après l'État. Elle gère environ \num{38 000} hectares de bananiers, de palmiers à huile et d'hévéas. Chaque matin, des milliers de travailleurs — coupeurs, chargeurs, conducteurs d'engins, techniciens de laboratoire, ingénieurs agronomes — convergent vers les \cle{camps} (Mile 4, Ekona, Idenau…). Sans cette activité, toute l'économie locale (commerces, transports, artisans) s'effondrerait. C'est ce que l'on appelle un \cle{effet d'entraînement} ou un \cle{multiplicateur économique}.
\end{exemplebox}

\noindent
L'agriculture de plantation génère d'abord un \textbf{emploi massif et structurant}. On distingue :
\begin{itemize}
    \item Les \textbf{emplois directs permanents} : ouvriers agricoles (coupe, ramassage, entretien), cadres moyens (contremaîtres, chefs de section), cadres supérieurs (ingénieurs, gestionnaires, comptables). Ensemble, les grandes firmes de plantation (CDC, SOCAPALM, SAFACAM, HEVECAM, SOSUCAM, PHP, Boh Plantations) emploient directement \textbf{environ 45 000 à 55 000 salariés permanents} (effectifs variables selon les plans sociaux et les replantations).
    \item Les \textbf{emplois directs saisonniers} : lors des pics de récolte (banane, coton, canne à sucre), le recours à la main-d'œuvre occasionnelle fait gonfler les effectifs de \num{20 000} à \num{30 000} travailleurs supplémentaires.
    \item Les \textbf{emplois induits (indirects)} : transporteurs routiers (camions citernes pour le latex, porte-conteneurs pour la banane), commerçants des camps, artisans (mécanique, menuiserie), personnel des sous-traitants (sécurité, restauration, maintenance industrielle).
\end{itemize}

\begin{attention}[Précision cruciale sur la SODECOTON]
Contrairement aux autres firmes citées, la \textbf{SODECOTON} (Société de Développement du Coton) ne gère \textbf{pas} de grandes plantations industrielles de coton. Elle assure l'encadrement technique, la fourniture d'intrants (semences, engrais, pesticides) et la collecte/égrenage du coton produit par \textbf{plus de 200 000 paysans cotonniers} (exploitations familiales de 1 à 5 ha) dans le Nord (Bénoué, Mayo-Rey, Faro, Diamaré, Mayo-Danay). Son « impact plantation » concerne ses fermes de multiplication de semences (ex: Maroua, Garoua) et ses usines d'égrenage, mais son modèle est l'\textbf{agriculture contractuelle paysanne}, et non la plantation salariée. Ne pas confondre les deux modèles au Bac.
\end{attention}

\begin{aretenir}[Chiffres clés de l'emploi agro-industriel]
Au total, le secteur des grandes plantations (hors encadrement paysan SODECOTON) assure la subsistance directe de \textbf{plus de 300 000 à 500 000 personnes} (salariés permanents + saisonniers + ayants droit) au Cameroun, soit un pilier démographique non négligeable dans les zones rurales concernées (Moungo, Sanaga-Maritime, Wouri, Nkam, Haut-Nyong, Mayo-Danay).
\end{aretenir}

\noindent
Ces emplois s'accompagnent d'une \textbf{amélioration tangible des conditions de vie} dans les camps, véritables « villes-usines » :
\begin{itemize}
    \item \textbf{Logement :} cases en dur (parpaings, tôles) attribuées aux familles d'ouvriers, souvent avec électricité et eau courante (bornes-fontaines ou branchement individuel pour les cadres).
    \item \textbf{Santé :} dispensaires ou hôpitaux d'entreprise (ex: Hôpital de la CDC à Tiko, Centre de santé de la SOCAPALM à Dibombari, Hôpital de la SOSUCAM à Mbandjock) offrant des soins gratuits ou subventionnés aux employés et leurs familles.
    \item \textbf{Éducation :} écoles primaires et maternelles dans les camps (ex: École primaire de la SOSUCAM à Mbandjock), parfois collèges ; prise en charge partielle des frais de scolarité.
    \item \textbf{Vie sociale :} marchés, lieux de culte, terrains de sport, centres sociaux.
\end{itemize}

\begin{definition}[Camp de plantation]
Un \cle{camp de plantation} est un habitat groupé, planifié et géré par l'entreprise agro-industrielle pour loger sa main-d'œuvre permanente et sa famille à proximité immédiate des champs et des usines. Il concentre des équipements collectifs (santé, éducation, commerce, loisirs) financés par la firme.
\end{definition}

\noindent
Sur le plan \textbf{macroéconomique}, ces firmes sont des pourvoyeuses essentielles de \textbf{devises} et de \textbf{recettes fiscales} :
\begin{itemize}
    \item \textbf{Exportations :} Banane (CDC, PHP, Boh Plantations), Coton (SODECOTON / paysans), Huile de palme (SOCAPALM, SAFACAM, CDC), Caoutchouc (HEVECAM, CDC, SAFACAM), Sucre (SOSUCAM). Ces produits représentent la part dominante des \textbf{exportations agricoles non-pétrolières} (hors cacao et café).
    \item \textbf{Fiscalité :} Impôts sur les sociétés, taxes à l'exportation, cotisations sociales (CNPS), taxes foncières. La contribution au budget de l'État se chiffre en \textbf{dizaines de milliards de FCFA par an}.
\end{itemize}

\begin{popfigure}
\begin{tikzpicture}
\begin{axis}[
    ybar,
    width=\linewidth,
    height=0.6\linewidth,
    bar width=14pt,
    ymin=0, ymax=40,
    ylabel={Part dans la valeur des exportations agricoles (\%)},
    xtick={1,2,3,4,5,6},
    xticklabels={Cacao/Café, Banane, Coton, Huile de palme, Caoutchouc, Sucre/Autres},
    xticklabel style={align=center, font=\small},
    ytick style={draw=none},
    ymajorgrids=true,
    grid style={popline},
    nodes near coords={\pgfmathprintnumber{\pgfplotspointmeta}\%},
    nodes near coords align={vertical},
    enlarge x limits=0.2,
    legend style={at={(0.5,-0.18)}, anchor=north, legend columns=3, font=\footnotesize, /tikz/every even column/.append style={column sep=5pt}},
]
\addplot[fill=popblue!80, draw=popblue, forget plot] coordinates {(1,35)};
\addplot[fill=poporange!80, draw=poporange, forget plot] coordinates {(2,25)};
\addplot[fill=popgreen!80, draw=popgreen, forget plot] coordinates {(3,15)};
\addplot[fill=poppurple!80, draw=poppurple, forget plot] coordinates {(4,12)};
\addplot[fill=poppink!80, draw=poppink, forget plot] coordinates {(5,8)};
\addplot[fill=popgold!80, draw=popgold, forget plot] coordinates {(6,5)};
\end{axis}
\end{tikzpicture}
\legende{Part indicative des principaux produits d'exportation agricole du Cameroun (valeurs arrondies, sources : MINADER, Douanes, rapports annuels firmes 2020-2023). Les produits de grandes plantations (banane, coton, palme, hévéa, sucre) totalisent ensemble \textbf{environ 60\%} de la valeur des exportations agricoles.}
\end{popfigure}

\noindent
Enfin, l'\textbf{aménagement du territoire} profite de ces implantations : ouverture et entretien de \textbf{routes et pistes} (ex: route Mbandjock-Nkoteng par SOSUCAM, axes Dibombari-Édéa par SOCAPALM), construction de \textbf{ponts}, électrification rurale étendue depuis les usines (turbo-alternateurs des sucreries, huileries).

\courssub{Des impacts négatifs lourds : foncier, environnemental, social}

\begin{savaistu}[Conflits fonciers à Dibombari et Édéa : la parole aux riverains]
« Nos arrière-grands-parents ont défriché cette forêt. Aujourd'hui, la SOCAPALM nous interdit d'y cultiver le manioc sous prétexte que c'est leur concession. On nous a payé une misère à l'époque, sans titre foncier clair. Nos cimetières sont dans la plantation. » — Extrait d'une pétition d'associations de riverains (Collectif des victimes de la SOCAPALM, 2021). Depuis les années 2010, des sit-in, des mémorandums au Premier Ministre et des procès opposent régulièrement les communautés Bakoko, Bassa et Douala aux concessions de la SOCAPALM (ex: concession de la Mbo, de Dibombari, de Kienké). Elles réclament la \textbf{rétrocession des terres non exploitées} ou une \textbf{revalorisation des redevances foncières} (souvent fixées à des montants dérisoires : \SI{500}{FCFA} à \SI{2000}{FCFA} l'hectare/an selon les baux anciens).
\end{savaistu}

\noindent
Le revers de la médaille est tout aussi structurant. Le premier contentieux est \textbf{foncier} : l'\textbf{accaparement des terres communautaires}.

\begin{definition}[Accaparement des terres (\emph{land grabbing})]
L'\cle{accaparement des terres} désigne la prise de contrôle — par achat, location à long terme (bail emphytéotique 50-99 ans) ou concession d'État — de vastes superficies de terres communautaires ou forestières par des firmes agro-industrielles ou des investisseurs privés, souvent sans \textbf{consentement libre, préalable et éclairé (CLPE)} des populations autochtones, et avec une compensation jugée insuffisante.
\end{definition}

\noindent
Concrètement, les concessions coloniales (CDC, HEVECAM) ou post-indépendance (SOCAPALM, SOSUCAM, SAFACAM) ont été attribuées sur des \textbf{terres coutumières} considérées comme « vacantes et sans maître » par l'État, ignorant les droits d'usage (chasse, cueillette, agriculture itinérante, lieux sacrés). Les villages se retrouvent \textbf{enclavés} au cœur des plantations (ex: villages « îlots » dans la concession HEVECAM à Nkongsamba, ou riverains de la SOSUCAM à Mbandjock/Nkoteng), sans espace pour leur agriculture vivrière, provoquant \textbf{insécurité alimentaire locale} et \textbf{conflits latents}.

\begin{attention}[Piège du vocabulaire : « Terres vacantes » $\neq$ « Terres vides »]
Au Bac, n'écris jamais que les plantations se sont installées sur des « terres vides ». Le Code foncier de 1974 (modifié 2005) définit les « terres vacantes et sans maître » comme des terres non immatriculées, mais elles sont \textbf{toujours soumises à des droits coutumiers d'usage}. Les populations y pratiquent l'agriculture de jachère, la chasse, la pharmacopée. Les ignorer, c'est commettre une erreur géographique et juridique grave.
\end{attention}

\noindent
Le deuxième contentieux est \textbf{environnemental}. L'agriculture intensive de plantation repose sur une \textbf{monoculture} à haut rendement, qui impose :
\begin{itemize}
    \item \textbf{Déforestation massive} : conversion de forêt dense humide (CDC, HEVECAM, SOCAPALM, SAFACAM) ou de savane boisée (SOSUCAM, fermes SODECOTON) en peuplements monospecifiques. Perte de biodiversité, fragmentation des habitats (éléphants, primates dans le Moungo, corridor Campo-Ma'an).
    \item \textbf{Pollution chimique} : usage intensif d'\textbf{herbicides} (glyphosate, paraquat — interdit dans l'UE, usage encadré mais persistant au Cameroun), d'\textbf{insecticides}, de \textbf{fongicides} et d'\textbf{engrais minéraux}. Ruissellement vers les rivières (Mungo, Wouri, Sanaga, Bénoué) $\rightarrow$ eutrophisation, mortalité piscicole, contamination nappes phréatiques.
    \item \textbf{Pollution industrielle} : rejets des \textbf{huileries} (effluents acides, riches en matière organique — \cle{POME} : \emph{Palm Oil Mill Effluent}) et des \textbf{savonneries} (soude, graisses) souvent insuffisamment traités avant déversement en milieu naturel. Odeurs nauséabondes pour les riverains.
    \item \textbf{Érosion des sols} : labour profond pour la canne à sucre (SOSUCAM), sols nus entre les rangées de jeunes palmiers/hévéas.
    \item \textbf{Pression sur la ressource en eau} : irrigation massive pour la canne à sucre (SOSUCAM : pompage dans la Sanaga), conflits d'usage en saison sèche.
\end{itemize}

\noindent
Le troisième contentieux est \textbf{social et salarial} :
\begin{itemize}
    \item \textbf{Bas salaires :} Le SMIG agricole (\SI{41 875}{FCFA} en 2024, revalorisé) est souvent le salaire de base des ouvriers de champ. Les primes (rendement, ancienneté, insalubrité) restent modestes. Difficulté d'épargne, endettement auprès des coopératives des camps.
    \item \textbf{Dépendance alimentaire :} Les camps sont des « déserts alimentaires » : pas d'espace pour le champ vivrier, prix élevés au marché du camp (monopole de la coopérative ou revendeurs). Les ouvriers achètent tout (manioc, plantain, poisson, bois de feu).
    \item \textbf{Précarité contractuelle :} Recours croissant au \textbf{sous-traitement} (entreprises de prestation de services pour la coupe, le chargement, la maintenance) $\rightarrow$ perte des avantages sociaux (CNPS, logement, santé) pour des milliers de travailleurs « fantômes ».
    \item \textbf{Ségrégation spatiale :} Séparation nette entre \textbf{camps des cadres} (villas spacieuses, électricité 24/24, eau courante, clubs de golf/tennis) et \textbf{camps des ouvriers} (alignements de cases serrées, coupures d'eau fréquentes, latrines collectives). Cette fracture urbaine reproduit les inégalités sociales à l'intérieur de la plantation.
\end{itemize}

\courssub{Analyse guidée de documents : regards croisés sur la réalité des plantations}

\begin{methode}[Commenter un document photographique sur un camp de plantation]
\textbf{Étape 1 : Identification.} Nature du doc (photo aérienne / au sol), date, localisation (firme, camp), auteur si connu.
\textbf{Étape 2 : Description structurée.} \textit{Habitat :} type de cases (dur/précaire), densité, alignement. \textit{Équipements visibles :} eau (bornes, châteaux), électricité (poteaux, compteurs), routes (bitume/piste), bâtiments publics (école, hôpital, marché, église/mosquée). \textit{Activité humaine :} présence femmes/enfants, vélos/motos, étals.
\textbf{Étape 3 : Analyse / Interprétation.} Comparer camp cadres vs camp ouvriers (ségrégation). Évaluer le niveau de services (intégration ? paternalisme ?). Relier à la politique sociale de la firme.
\textbf{Étape 4 : Conclusion critique.} Le camp est-il un vecteur d'intégration (mélange ethnique, services) ou de relégation (isolement, dépendance) ?
\end{methode}

\begin{exoresolu}[Étude de cas : Conflit foncier SOCAPALM / Communautés de la Sanaga-Maritime (Article de presse simulé)]
\textbf{Énoncé :} À partir de l'extrait d'article ci-dessous, identifiez les acteurs, la nature du conflit, les arguments des deux parties et proposez une issue négociée.

\textit{« Dibombari, 12 mars 2023. Le Collectif des Riverains de la Concession de la Mbo (CRCM) a déposé un mémorandum au Ministre des Domaines. Ils contestent l'extension de 5 000 ha accordée à la SOCAPALM en 2018. « Nos terroirs ancestraux ont été réduits à peau de chagrin. Les redevances versées aux chefferies (1 500 FCFA/ha/an) ne profitent pas aux populations. Nos jeunes n'ont plus de terre pour planter le macabo. La rivière Mbo est polluée par l'huilerie. » La SOCAPALM répond : « Nous respectons le bail de 1968 renouvelé. Nous employons 3 000 permanents de la zone. Nous avons construit l'école de Mbo, le centre de santé, entretenu la route. L'extension se fait sur des terres déjà plantées en palmiers vieillissants (replantation). » »}

\tcblower
\textbf{Solution détaillée :}
\begin{enumerate}
    \item \textbf{Acteurs :} Collectif Riverains (CRCM) vs SOCAPALM (filiale Socfin) + Administration (Ministère des Domaines, Chefferies traditionnelles).
    \item \textbf{Nature :} Conflit foncier (extension concession vs droits coutumiers) + environnemental (pollution rivière) + social (emploi vs terre vivrière).
    \item \textbf{Arguments Riverains :} Spoliation historique (bail colonial), compensation dérisoire (1 500 FCFA/ha/an $\ll$ valeur économique), perte souveraineté alimentaire, pollution, détournement redevances par élites locales (chefferies).
    \item \textbf{Arguments SOCAPALM :} Légalité (bail, décret), investissements sociaux (école, santé, route = \cle{responsabilité sociétale des entreprises - RSE}), emploi local, replantation (pas nouvelle déforestation nette).
    \item \textbf{Issue négociée (piste) :} Audit foncier indépendant $\rightarrow$ identification terres non plantées/non exploitées $\rightarrow$ rétrocession aux villages. Réévaluation redevances (indexation sur CA ou prix huile de palme). Fonds de développement communautaire géré par riverains (pas chefferies). Plan de dépollution rivière (bassins de lagunage aux normes). Comité de suivi tripartite (Entreprise, État, Riverains).
\end{enumerate}
\end{exoresolu}

\courssub{Construire un bilan nuancé : méthode et pièges à éviter}

\begin{methode}[Construire un bilan géographique structuré (Social / Économique / Environnemental)]
Pour réussir l'exercice de bilan (fréquent au Bac : « Présentez les impacts de l'agriculture de plantation »), suis cette démarche en 3 temps :
\begin{enumerate}
    \item \textbf{Prépare un tableau à double entrée} (voir illustration ci-dessous) avec 3 lignes (Social, Économique, Environnemental) et 2 colonnes (Positifs $+$ / Négatifs $-$).
    \item \textbf{Remplis chaque case} avec des \textbf{faits précis, chiffrés, localisés} (noms de firmes, de lieux, ordres de grandeur). Évite les généralités (« ça pollue » $\rightarrow$ « Rejets POME huilerie SOCAPALM Dibombari dans rivière Mbo »).
    \item \textbf{Rédige la synthèse} en 3 paragraphes (un par dimension), en utilisant des connecteurs logiques (\emph{cependant, toutefois, en contrepartie, par ailleurs}). Termine par une phrase d'ouverture : « L'enjeu actuel est la transition vers une agro-industrie durable (certification RSPO, agriculture contractuelle, contractualisation foncière équitable). »
\end{enumerate}
\end{methode}

\begin{popfigure}
\begin{center}
\begin{tikzpicture}[
    cell/.style={rectangle, draw=popline, minimum height=1.5cm, align=center, text width=7.2cm, font=\small},
    header/.style={rectangle, draw=popline, fill=popblue, text=white, minimum height=0.9cm, align=center, text width=7.2cm, font=\small\bfseries},
    headercat/.style={rectangle, draw=popline, fill=popdark, text=white, minimum height=0.9cm, align=center, text width=3.5cm, font=\small\bfseries},
    rowcat/.style={rectangle, draw=popline, fill=popblueL, minimum height=1.5cm, align=center, text width=3.5cm, font=\small\bfseries},
]
\matrix (tab) [row sep=0pt, column sep=0pt] {
    \node[headercat] {}; & \node[header] {\textbf{Impacts POSITIFS ($+$)}}; & \node[header] {\textbf{Impacts NÉGATIFS ($-$)}}; \\
    \node[rowcat] {\textbf{SOCIAL}}; &
    \node[cell, align=center]{Emplois massifs (permanents $\approx$ 50 000 + saisonniers $\approx$ 25 000).\\ Salaires réguliers $\rightarrow$ pouvoir d'achat.\\ Services gratuits : santé (hôpitaux CDC, SOSUCAM), éducation (écoles camps), logement (cases en dur), eau/électricité.\\ Brassage ethnique $\rightarrow$ intégration nationale (voir Dossier 2).}; &
    \node[cell, align=center]{Ségrégation spatiale (camps cadres vs ouvriers).\\ Bas salaires (proche SMIG), précarité (sous-traitance).\\ Dépendance alimentaire (pas de champ vivrier).\\ Conflits fonciers : spoliation terres coutumières, enclavement villages.\\ Perte identité culturelle (cimetières, lieux sacrés détruits).}; \\
    \node[rowcat] {\textbf{ÉCONOMIQUE}}; &
    \node[cell, align=center]{Devises : $\approx$ 60\% exportations agricoles (Banane, Coton, Palme, Hévéa, Sucre).\\ Recettes fiscales : milliards FCFA/an (IS, taxes export, CNPS).\\ Infrastructures : routes, ponts, électricité rurale ouverts par firmes.\\ Développement local : marchés, commerces, banques (agences dans camps).}; &
    \node[cell, align=center]{Dépendance structurelle : mono-exportation (prix mondiaux volatils).\\ Rapatriement bénéfices (firmes étrangères : Socfin, Bolloré, etc.).\\ Faible transformation locale (export brut : banane verte, coton graine, latex).\\ Coûts cachés : dégradation routes par camions lourds, santé publique (pollution).}; \\
    \node[rowcat] {\textbf{ENVIRONNEMENTAL}}; &
    \node[cell, align=center]{Revalorisation terres « marginales » (savanes SOSUCAM/SODECOTON).\\ Couvert végétal permanent (palmiers, hévéas) $\rightarrow$ stockage carbone vs jachère.\\ Reboisement haies brise-vent (CDC, HEVECAM).\\ Traitement effluents : bassins lagunage (certains sites certifiés RSPO/ISO 14001).}; &
    \node[cell, align=center]{Déforestation massive (forêt dense $\rightarrow$ monoculture).\\ Perte biodiversité (corridors éléphants coupés Moungo).\\ Pollution chimique : pesticides (glyphosate, paraquat), engrais $\rightarrow$ rivières/nappes.\\ Pollution industrielle : POME huileries, odeurs savonneries.\\ Érosion sols (canne à sucre, jeunes plants).\\ Épuisement ressources en eau (irrigation SOSUCAM).}; \\
};
\end{tikzpicture}
\end{center}
\legende{Bilan structuré des impacts socio-économiques et environnementaux des grandes plantations au Cameroun. Ce tableau est un outil de révision indispensable pour le Bac.}
\end{popfigure}

\begin{attention}[Erreur fatale au Bac : le bilan unilatéral]
\textbf{Ne présentez JAMAIS un bilan uniquement positif (« Les plantations développent le pays ») ou uniquement négatif (« Les plantations pillent le pays »).} Le correcteur attend un \textbf{jugement nuancé} : « L'agriculture de plantation est un pilier de l'économie camerounaise (devises, emplois, infrastructures) mais son modèle foncier hérité de la colonisation et ses pratiques intensives génèrent des conflits sociaux et une dégradation environnementale qui menacent sa soutenabilité. La résolution passe par la sécurisation foncière des riverains, la certification environnementale (RSPO, FSC) et une meilleure redistribution de la valeur ajoutée. »
\end{attention}

\begin{aretenir}[Synthèse : L'agropole, horizon d'intégration ?]
Le concept d'\cle{agropole} (pôle agro-industriel intégré) tente de dépasser ces contradictions. Exemple : \textbf{Agropole de Mbandjock} (autour de SOSUCAM) ou \textbf{Agropole du Sud-Ouest} (autour CDC). Objectif : associer \textbf{grandes firmes + PME locales + paysans contractuels + centres de formation + transformation locale (usines sucre, huile, caoutchouc, textile)} pour créer de la valeur ajoutée sur place, fixer les populations et réduire l'exode rural. C'est une piste d'avenir pour l'intégration nationale, à condition que le foncier soit sécurisé et l'environnement préservé.
\end{aretenir}

\coursec{Dossier 2 : la main-d'œuvre dans les grandes plantations, vecteur d'intégration nationale}

\courssub{De l'impact socio-économique à l'acteur central : la main-d'œuvre}
Après avoir analysé les retombées économiques (emplois, revenus, infrastructures) et les externalités négatives (accaparement des terres, pollution), nous nous penchons maintenant sur l'acteur central de ces plantations : la main-d'œuvre. C'est elle qui transforme le capital et la terre en production, et c'est à travers sa mobilité, son installation et sa vie quotidienne que les grandes plantations deviennent de véritables \emph{creusets d'intégration nationale}.

\courssub{Origine des travailleurs et migrations de main-d'œuvre}
\begin{definition}[Main-d'œuvre de plantation]
Ensemble des travailleurs (permanents, saisonniers, cadres) recrutés par une firme agro‑industrielle pour exploiter une grande plantation. Leur origine géographique est multiple, ce qui fait de chaque site un carrefour régional.
\end{definition}

Les firmes (\cle{CDC}, \cle{SOCAPALM}, \cle{SAFACAM}, \cle{SODECOTON}, \cle{SOSUCAM}, \cle{HEVECAM}) recrutent d'abord localement, mais les besoins massifs (plusieurs milliers d'ouvriers par site) imposent un appel à la main-d'œuvre venue d'autres régions. Depuis les années 1960, on observe un flux continu du \textbf{Grand‑Nord} (Adamaoua, Nord, Extrême‑Nord) vers les bassins de plantation du \textbf{Littoral} (Dibombari, Mbanga, Nkongsamba) et du \textbf{Sud‑Ouest} (Tiko, Mutengene, Kumba, Ekona). Ce mouvement est alimenté par la recherche de salaires plus élevés que dans l'agriculture vivrière et par la présence d'infrastructures sociales (écoles, centres de santé) dans les camps.

\begin{exemplebox}[Chiffres clés du recrutement (ordre de grandeur)]
À la \cle{CDC} (Tiko, Mutengene, Ekona), sur environ 12\,000 travailleurs permanents recensés récemment (rapports d'activité 2021‑2022), près de 55\,\% sont originaires des régions du Nord, de l'Adamaoua et de l'Extrême‑Nord, 30\,\% du Littoral et du Sud‑Ouest, le reste provenant du Centre, de l'Est et de l'Ouest. La \cle{SOCAPALM} (site de Dibombari) présente une répartition comparable pour son bassin littoral.
\end{exemplebox}

\begin{popfigure}
\begin{center}
\includegraphics[width=\linewidth,height=11cm,keepaspectratio]{N05/cmr_regions.jpg}
\end{center}
\legende{Fond de carte : le Cameroun en grands ensembles (carte en anglais), avec la voie ferrée Ngaoundéré--Yaoundé--Douala. Les flux de main-d'œuvre ne figurent pas sur cette carte : le croquis à construire les trace en flèches depuis le Nord, l'Adamaoua et l'Extrême-Nord vers les bassins de plantation du Littoral (Douala, Nkongsamba) et du Sud-Ouest (Limbe, mont Cameroun). \\ Source : Wikimedia Commons, Peter Fitzgerald (modifié par Burmesedays), CC BY 3.0.}
\end{popfigure}

\courssub{Organisation spatiale des camps : cadres et ouvriers}
L'implantation d'une plantation s'accompagne systématiquement de la création de \textbf{camps} (ou cités) qui logent travailleurs et familles. L'espace reflète la hiérarchie sociale de l'entreprise.

\begin{definition}[Camp de plantation]
Espace résidentiel créé par une firme agro‑industrielle pour loger sa main-d'œuvre et sa famille, doté de services sociaux (écoles, centres de santé, marchés, lieux de culte) et organisé selon la hiérarchie professionnelle.
\end{definition}

\begin{methode}[Construire un croquis géographique d'un camp de plantation]
\begin{enumerate}
\item \textbf{Titre} clair et daté (ex. « Camp CDC de Mutengene, 2023 »).
\item \textbf{Orientation} : flèche du nord en haut à droite.
\item \textbf{Figurés} simples : polygones pour les zones (usine, champs), points pour les bâtiments, traits pour les routes, hachures pour les zones d'habitat.
\item \textbf{Légende ordonnée} : du plus important (usine, camp des cadres) au moins important (terrain de sport), sans phrases, seulement des noms de figurés.
\item \textbf{Échelle indicative} (ex. 1 cm = 200 m) et \textbf{source} (photo aérienne, plan d'entreprise).
\end{enumerate}
\end{methode}

\begin{popfigure}
\begin{tikzpicture}[scale=0.9]
% Usine
\fill[popgray!40] (0,0) rectangle (2,1.5);
\node[popdark] at (1,0.75) {\small Usine};
% Camp des cadres (villas)
\fill[popblue!30] (3,2) rectangle (5,4);
\foreach \x in {3.3,3.9,4.5} \foreach \y in {2.3,2.9,3.5} {
  \draw[popblue, thick] (\x,\y) rectangle ++(0.4,0.4);
}
\node[popdark] at (4,4.3) {\small Camp des cadres};
% Camp des ouvriers (baraquements alignés)
\fill[poporange!30] (0,2.5) rectangle (2.5,5);
\foreach \y in {2.7,3.3,3.9,4.5} {
  \draw[poporange, thick] (0.2,\y) rectangle ++(2.1,0.4);
}
\node[popdark] at (1.25,5.3) {\small Camp des ouvriers};
% École
\fill[popgreen!30] (5.5,0.5) rectangle (7,1.5);
\node[popdark] at (6.25,1) {\small École};
% Centre de santé
\fill[poppurple!30] (5.5,2) rectangle (7,3);
\node[popdark] at (6.25,2.5) {\small Centre de santé};
% Marché
\fill[poppink!30] (5.5,3.5) rectangle (7,4.2);
\node[popdark] at (6.25,3.85) {\small Marché};
% Champs (hachures)
\pattern[pattern=north east lines, pattern color=popgreen!50] (0,-2) rectangle (7,0);
\node[popdark] at (3.5,-1) {\small Champs de palmiers / hévéas};
% Routes
\draw[popdark, thick] (1,0) -- (1,-2);
\draw[popdark, thick] (4,2) -- (4,0.5);
\draw[popdark, thick] (1,2.5) -- (4,2.5);
% Nord
\draw[->, popmuted] (6.5,4.5) -- (6.5,5) node[above] {N};
% Légende intégrée
\begin{scope}[shift={(7.2,3.5)}]
\node[anchor=west, popdark] {\textbf{Légende}};
\node[anchor=west, popdark] at (0,-0.4) {\textcolor{popgray!60}{\rule{0.4cm}{0.2cm}} Usine};
\node[anchor=west, popdark] at (0,-0.8) {\textcolor{popblue!60}{\rule{0.4cm}{0.2cm}} Camp des cadres (villas)};
\node[anchor=west, popdark] at (0,-1.2) {\textcolor{poporange!60}{\rule{0.4cm}{0.2cm}} Camp des ouvriers (baraquements)};
\node[anchor=west, popdark] at (0,-1.6) {\textcolor{popgreen!60}{\rule{0.4cm}{0.2cm}} École};
\node[anchor=west, popdark] at (0,-2.0) {\textcolor{poppurple!60}{\rule{0.4cm}{0.2cm}} Centre de santé};
\node[anchor=west, popdark] at (0,-2.4) {\textcolor{poppink!60}{\rule{0.4cm}{0.2cm}} Marché};
\node[anchor=west, popdark] at (0,-2.8) {\textcolor{popgreen!50}{\rule{0.4cm}{0.2cm}} Champs (hachures)};
\node[anchor=west, popdark] at (0,-3.2) {\textcolor{popdark}{\rule{0.4cm}{0.2cm}} Routes principales};
\end{scope}
\end{tikzpicture}
\legende{Croquis schématique d'un camp de plantation type (ex. CDC Mutengene) : usine, camps différenciés, services sociaux, champs environnants.}
\end{popfigure}

\courssub{Structures sociales et services collectifs}
Les camps ne sont pas de simples dortoirs ; ils constituent de véritables \textbf{petites villes d'entreprise}. On y trouve systématiquement :
\begin{itemize}
\item \textbf{Écoles d'entreprise} (primaire, parfois secondaire) scolarisant les enfants des ouvriers et des cadres, souvent ouvertes aux riverains.
\item \textbf{Centres de santé / hôpitaux de plantation} (ex. Hôpital de la CDC à Tiko, Centre de santé SOCAPALM à Dibombari) offrant soins gratuits ou à tarifs réduits aux ayants droit.
\item \textbf{Terrains de sport}, \textbf{marchés hebdomadaires}, \textbf{lieux de culte} (églises, mosquées) reflétant la diversité religieuse.
\end{itemize}

Ces équipements créent une \emph{centralité} qui attire aussi des populations non salariées (commerçants, artisans), renforçant le rôle de polarisation du camp.

\begin{exemplebox}[Population d'un camp type]
Le camp CDC de Mutengene regroupe environ 8\,500 habitants (travailleurs + ayants droit) selon les estimations récentes (2023), soit la taille d'une commune rurale moyenne. L'école primaire y compte environ 1\,200 élèves, le centre de santé réalise environ 15\,000 consultations/an.
\end{exemplebox}

\courssub{La main-d'œuvre comme vecteur d'intégration nationale}
La coexistence quotidienne de Camerounais originaires des dix régions (et parfois de pays voisins : Nigeria, Tchad, RCA) dans un même camp produit un \textbf{brassage culturel} intense :
\begin{itemize}
\item \textbf{Langues véhiculaires} : le \emph{Pidgin English} (Sud‑Ouest) et le \emph{Français} (langue administrative) deviennent les langues de communication interethnique.
\item \textbf{Mariages mixtes} : unions entre ouvriers de régions différentes, favorisant la création de familles « nationales ».
\item \textbf{Fêtes et rites partagés} : célébrations du 20 mai (Fête nationale), du 1er mai (Fête du travail), des fêtes religieuses communes.
\end{itemize}

Ce vécu partagé forge un \textbf{sentiment d'appartenance commune} qui dépasse les clivages régionaux ou ethniques. Les camps deviennent ainsi des \emph{laboratoires d'intégration nationale} à l'échelle locale.

\begin{savaistu}[Le creuset de la CDC]
Dans les camps de la CDC à Tiko et Mutengene, on recense des travailleurs originaires de \emph{presque toutes les régions du Cameroun} (y compris l'Extrême‑Nord et l'Est). Les enquêtes sociologiques menées en 2021 montrent que 78\,\% des ouvriers déclarent avoir au moins un ami ou un conjoint d'une autre région que la leur.
\end{savaistu}

\courssub{Limites et tensions sociales}
Malgré cet apport intégrateur, des \textbf{tensions} persistent :
\begin{itemize}
\item \textbf{Écarts de conditions de vie} : villas climatisées, eau courante, électricité permanente pour les cadres ; baraquements collectifs, latrines partagées, coupures d'eau fréquentes pour les ouvriers.
\item \textbf{Grèves et conflits} : revendications salariales, meilleures conditions de logement, respect du code du travail (ex. grève CDC 2019, grève SOCAPALM 2021).
\item \textbf{Précarité des saisonniers} : contrats courts, absence de couverture sociale, logement précaire (tentes, hangars).
\end{itemize}

Ces limites rappellent que l'intégration par le travail ne saurait se substituer à une politique publique d'aménagement du territoire et de réduction des inégalités.

\courssub{Étude de cas guidée : croquis d'un camp de plantation (CDC Mutengene)}
\begin{exoresolu}[Croquis du camp CDC Mutengene à partir d'une photographie aérienne]
\textbf{Énoncé} : À partir de la photographie aérienne fournie (échelle 1/10\,000, nord en haut), réaliser un croquis géographique légendé du camp en faisant apparaître : l'usine, les deux types de camps, l'école, le centre de santé, le marché, les champs environnants et les axes de circulation principaux.

\textbf{Solution détaillée} :
\begin{enumerate}
\item \textbf{Observation} : repérer les formes géométriques : grand bâtiment rectangulaire central (usine), alignements réguliers au nord‑ouest (baraquements ouvriers), îlots spacieux au nord‑est (villas cadres), bâtiments isolés au sud (école, centre de santé, marché).
\item \textbf{Choix des figurés} : polygone plein gris pour l'usine ; hachures bleues pour camp cadres ; hachures orange pour camp ouvriers ; symboles conventionnels (école = drapeau, santé = croix, marché = panier).
\item \textbf{Traçage} : reporter les limites sur papier calque, simplifier les contours (pas de détails de toits).
\item \textbf{Légende hiérarchisée} : 1. Usine, 2. Camp des cadres, 3. Camp des ouvriers, 4. École, 5. Centre de santé, 6. Marché, 7. Champs, 8. Routes principales.
\item \textbf{Finalisation} : ajouter flèche nord, échelle graphique (1 cm = 200 m), titre et source.
\end{enumerate}
\textbf{Résultat} : le croquis obtenu correspond à la figure ci‑dessus (voir \emph{Croquis schématique d'un camp de plantation type}).
\end{exoresolu}

\courssub{Exercice d'application}
\begin{exercice}[Analyse d'un flux migratoire]
À partir du tableau ci‑dessous (effectifs par région d'origine à la SOCAPALM Dibombari, 2022), calculez la part (en \%) de chaque région, tracez un diagramme circulaire et commentez le rôle intégrateur de ce flux.

\begin{center}
\begin{tabularx}{\linewidth}{|l|X|}\hline
\textbf{Région d'origine} & \textbf{Effectif} \\ \hline
Adamaoua & 1\,200 \\ \hline
Nord & 1\,050 \\ \hline
Extrême‑Nord & 950 \\ \hline
Littoral & 800 \\ \hline
Sud‑Ouest & 700 \\ \hline
Centre & 400 \\ \hline
Est & 300 \\ \hline
Ouest & 250 \\ \hline
Total & 5\,650 \\ \hline
\end{tabularx}
\end{center}
\end{exercice}

\begin{corrige}[Exercice n°1]
\begin{enumerate}
\item \textbf{Calcul des parts} (formule : $\frac{\text{Effectif région}}{\text{Total}} \times 100$) :
\begin{itemize}
\item Adamaoua : $1200 / 5650 \times 100 \approx 21,2\,\%$
\item Nord : $1050 / 5650 \times 100 \approx 18,6\,\%$
\item Extrême‑Nord : $950 / 5650 \times 100 \approx 16,8\,\%$
\item Littoral : $800 / 5650 \times 100 \approx 14,2\,\%$
\item Sud‑Ouest : $700 / 5650 \times 100 \approx 12,4\,\%$
\item Centre : $400 / 5650 \times 100 \approx 7,1\,\%$
\item Est : $300 / 5650 \times 100 \approx 5,3\,\%$
\item Ouest : $250 / 5650 \times 100 \approx 4,4\,\%$
\end{itemize}
\item \textbf{Diagramme circulaire} : tracer huit secteurs aux angles proportionnels (ex. Adamaoua : $21,2 \times 3,6 \approx 76^\circ$).
\item \textbf{Commentaire} : les trois régions du Grand‑Nord représentent 56,6\,\% de la main-d'œuvre, confirmant l'orientation Nord‑Sud des migrations. La présence de toutes les régions, même minoritaires, illustre le brassage national au sein du camp.
\end{enumerate}
\end{corrige}

\courssub{Synthèse et ouverture}
La main-d'œuvre des grandes plantations, par sa diversité d'origine, son installation durable dans des camps dotés de services collectifs et sa vie quotidienne partagée, constitue un \textbf{vecteur puissant d'intégration nationale}. Toutefois, les inégalités spatiales et sociales qui structurent ces camps (cadre vs ouvrier, permanent vs saisonnier) limitent la portée de cette intégration. La prochaine leçon abordera les \textbf{agropoles} et la manière dont l'État tente de rééquilibrer l'aménagement du territoire à travers des pôles de compétitivité agricole.

\begin{attention}[Pièges à éviter au Bac]
\begin{itemize}
\item Confondre \emph{camp de plantation} et \emph{village traditionnel} : le camp est une création d'entreprise, planifiée, hiérarchisée.
\item Oublier la flèche du nord et l'échelle indicative sur un croquis.
\item Écrire des phrases entières dans la légende (ex. « Ici se trouve l'usine ») au lieu de termes nominaux (« Usine »).
\item Négliger la distinction entre main-d'œuvre permanente et saisonnière dans l'analyse des flux.
\end{itemize}
\end{attention}

\coursec{Travaux pratiques : cartographier et représenter les grandes plantations}

Après avoir analysé le rôle de la main-d'œuvre comme vecteur d'intégration nationale au sein des camps de plantation, nous passons maintenant à la mise en œuvre concrète des savoir-faire cartographiques et graphiques exigés par le programme. Ces travaux pratiques (TP) ne sont pas de simples exercices d'application : ils constituent le cœur de l'évaluation au Baccalauréat. Maîtriser la construction d'une carte, d'un croquis ou d'un diagramme, c'est s'assurer des points décisifs le jour de l'épreuve. Nous allons procéder par étapes, du respect des règles de base à la réalisation autonome de quatre TP types, pour finir par une grille d'auto-évaluation rigoureuse.

\courssub{Rappel des règles cartographiques et graphiques de base}

Avant de tracer le moindre trait, rappelons les exigences incontournables. Une production graphique sans titre, sans légende hiérarchisée, sans orientation ni source est systématiquement pénalisée, voire non corrigée.

\begin{methode}[Règles d'or pour une carte ou un croquis réussis]
\textbf{1. Titre :} Informatif, précis, situé en haut (centré ou à gauche). Exemple : \emph{« Cameroun : localisation des grandes plantations agro-industrielles (2023) »}.\par
\textbf{2. Orientation :} Flèche du Nord (haut de la feuille par convention).\par
\textbf{3. Échelle :} Graphique (barre) de préférence, ou numérique si la précision le permet (ex. 1 cm pour 100 km).\par
\textbf{4. Légende :} Hiérarchisée (titre de légende, puis groupes : \emph{Implantations}, \emph{Réseaux}, \emph{Relief}, etc.). Figurés \textbf{ponctuels} (points, symboles) pour les villes/usines, \textbf{linéaires} (traits) pour les routes/voies ferrées/flux, \textbf{surfaciques} (hachures, couleurs) pour les zones de culture/régions. Un figuré = une information.\par
\textbf{5. Propreté :} Traits nets, coloriage régulier (crayons de couleur), écriture lisible (minuscules pour les noms communs, majuscules pour les noms propres).
\end{methode}

\begin{methode}[Construction d'un diagramme en barres (ou en colonnes) statistique]
\textbf{1. Titre :} Complet (sujet, lieu, date, unité). Ex. : \emph{« Superficies des principales plantations agro-industrielles au Cameroun en 2022 (en milliers d'hectares) »}.\par
\textbf{2. Axes :} Axe des abscisses (horizontal) : modalités qualitatives (noms des firmes). Axe des ordonnées (vertical) : grandeur quantitative avec unité (ex. \emph{Superficie (1000 ha)}). Graduation régulière, partant de zéro.\par
\textbf{3. Barres :} Largeur identique, espacement régulier. Hauteur \textbf{strictement proportionnelle} à la valeur. Coloriage uni ou hachuré, lisible en noir et blanc.\par
\textbf{4. Source :} Obligatoire en bas à gauche (ex. \emph{Source : MINADER, Annuaire statistique 2022}).
\end{methode}

\begin{attention}[Pièges fatals au Bac]
\begin{itemize}
    \item Oublier le titre, la source, l'échelle ou le Nord.
    \item Utiliser des figurés de surface (ronds proportionnels) pour des données ponctuelles sur une carte de localisation (on utilise des points ou symboles).
    \item Faire un diagramme en barres dont l'axe des ordonnées ne part pas de zéro (effet d'optique trompeur).
    \item Mélanger dans une même légende des figurés de natures différentes sans regroupement.
    \item Écrire la légende \emph{après} avoir colorié la carte : on prépare la légende \emph{avant} de tracer.
\end{itemize}
\end{attention}

\courssub{TP 1 : Légender une carte muette du Cameroun — Localisation des grandes firmes}

\begin{experience}[TP 1 : Localiser CDC, SOCAPALM, SAFACAM, SODECOTON, SOSUCAM, HEVECAM]
\textbf{Matériel :} Carte muette du Cameroun (fournie ci-dessous), crayons de couleur (vert, rouge, bleu, orange, violet, marron), règle, stylo noir.\par
\textbf{Consigne :} Sur la carte fournie, localisez les implantations principales des six firmes agro-industrielles majeures. Utilisez un figuré ponctuel distinct par firme (couleur + symbole). Complétez la légende, le titre, l'échelle et le Nord. Ajoutez les villes-pôles (Douala, Yaoundé, Garoua, Maroua, Ngaoundéré, Buea, Bertoua, Ebolowa) et les axes routiers structurants (N1, N2, N3, N10).
\end{experience}

\begin{popfigure}
\begin{center}
\includegraphics[width=\linewidth,height=11cm,keepaspectratio]{N07/cmr_muette.png}
\end{center}
\legende{Fond de carte vierge du Cameroun (limites des régions et des départements, fleuves, lacs, barre d'échelle), à compléter par l'élève : localisation des six firmes (figurés ponctuels distincts), villes-pôles, axes routiers, titre et légende hiérarchisée. \\ Source : Wikimedia Commons, Flappiefh, CC BY-SA 4.0.}
\end{popfigure}

\begin{aretenir}[Figurés attendus pour la légende du TP 1]
\begin{center}
\begin{tabularx}{\linewidth}{|l|X|X|}\hline
\rowcolor{popblueL}
\textbf{Firme} & \textbf{Figuré ponctuel suggéré} & \textbf{Zone d'implantation principale (région/département)} \\ \hline
CDC & \textcolor{popgreen}{$\blacktriangle$} Triangle vert & Sud-Ouest (Fako, Meme) : Tiko, Limbe, Muyuka \\ \hline
SOCAPALM & \textcolor{poporange}{$\bullet$} Cercle orange & Littoral (Sanaga-Maritime), Sud-Ouest (Ndian), Sud (Océan), Centre (Nyong-et-Kellé) \\ \hline
SAFACAM & \textcolor{popblue}{$\blacksquare$} Carré bleu & Littoral (Moungo), Centre (Haute-Sanaga) : Njombé, Mbono \\ \hline
HEVECAM & \textcolor{poppurple}{$\blacklozenge$} Losange violet & Centre (Mbam-et-Inoubou, Lékié) : Ntui, Bafia, Obala \\ \hline
SOSUCAM & \textcolor{poppink}{$\bigstar$} Étoile rose & Centre (Haute-Sanaga, Mbam-et-Kim) : Nkoteng, Mbandjock \\ \hline
SODECOTON & \textcolor{popgold}{$\blacktriangledown$} Triangle inversé or & Nord (Bénoué), Extrême-Nord (Diamaré, Mayo-Danay), Adamaoua (Mayo-Rey) : usines d'égrenage \\ \hline
\end{tabularx}
\end{center}
\emph{Note : Les superficies et localisations exactes des unités de production peuvent varier ; on retient les bassins principaux.}
\end{aretenir}

\begin{exoresolu}[Carte terminée : exemple de légende hiérarchisée]
Énoncé : À partir de la carte muette fournie, proposer une légende hiérarchisée complète pour la localisation des six firmes agro-industrielles.
\tcblower
\textbf{Titre :} \emph{Cameroun : Implantation des grandes firmes agro-industrielles (année 2023)}\par
\textbf{Légende :}\par
\textbf{Implantations des firmes :}\par
\textcolor{popgreen}{$\blacktriangle$} CDC (hévéa, palmier, bananier) \quad
\textcolor{poporange}{$\bullet$} SOCAPALM (palmier à huile) \quad
\textcolor{popblue}{$\blacksquare$} SAFACAM (palmier, hévéa) \quad
\textcolor{poppurple}{$\blacklozenge$} HEVECAM (hévéa) \quad
\textcolor{poppink}{$\bigstar$} SOSUCAM (canne à sucre) \quad
\textcolor{popgold}{$\blacktriangledown$} SODECOTON (usines égrenage coton)\par
\textbf{Villes et réseaux :}\par
$\bullet$ Chef-lieu de région \quad \textbf{---} Route nationale bitumée \quad \textbf{$\sim\sim$} Voie ferrée (Transcam) \quad \textbf{$\sim$} Fleuve navigable (Bénoué)\par
\textbf{Orientation :} Flèche Nord \quad \textbf{Échelle :} 1 cm $\approx$ 100 km \quad \textbf{Source :} MINADER, Cartes des périmètres concessionnés.
\end{exoresolu}

\courssub{TP 2 : Construire un croquis d'un camp de plantation à partir d'une vue aérienne}

\begin{experience}[TP 2 : Croquis d'organisation spatiale d'un camp de plantation (type CDC/SOCAPALM)]
\textbf{Document support :} Photographie aérienne ou image satellite (Google Earth) d'un camp type (ex. camp de Tiko pour CDC, ou Dibombari pour SOCAPALM).\par
\textbf{Consigne :} Réalisez un croquis structuré à l'encre et au crayon de couleur sur papier calque ou feuille blanche. Respectez la méthode ci-dessous.
\end{experience}

\begin{methode}[Étapes pour construire un croquis géographique à partir d'une photo aérienne]
\textbf{1. Analyse :} Repérez les grandes unités spatiales (zone d'habitat cadres, zone d'habitat ouvriers, zone administrative/usine, équipements collectifs : école, dispensaire, marché, terrain de sport, réseaux viaires, limite plantation/habitat).\par
\textbf{2. Cadre et orientation :} Tracez le cadre, indiquez le Nord (souvent en haut de l'image).\par
\textbf{3. Structure :} Dessinez les limites des grandes zones (traits fins).\par
\textbf{4. Figurés surfaciques :} Hachures ou couleurs légères pour différencier : habitat cadres (espacé, grands toits), habitat ouvriers (dense, aligné), zone industrielle (grands bâtiments), équipements (symboles), espaces verts/plantations (vert).\par
\textbf{5. Figurés ponctuels/linéaires :} Routes principales, chemin de fer s'il y a lieu, points d'eau.\par
\textbf{6. Légende hiérarchisée :} Groupe \emph{Habitat}, \emph{Activités/Production}, \emph{Équipements}, \emph{Communication}.\par
\textbf{7. Titre complet + Source de la photo.}
\end{methode}

\begin{popfigure}
\begin{tikzpicture}[scale=0.9]
% Cadre
\draw[popdark, thick] (0,0) rectangle (12,8);

% Nord
\draw[->, popdark, thick] (11.2,7.2) -- (11.2,7.7);
\node[font=\small, anchor=south, popdark] at (11.2,7.8) {N};

% Zones surfaciques
% Plantation (entourage)
\fill[popgreen!15] (0,0) rectangle (12,8);
% Zone habitat cadres (Nord-Ouest)
\fill[popblue!20] (0.5,5.5) rectangle (4,7.5);
\draw[popblue, thin] (0.5,5.5) grid[step=0.4] (4,7.5);
% Zone habitat ouvriers (Sud-Ouest)
\fill[poporange!20] (0.5,0.5) rectangle (4,4.5);
\draw[poporange, thin] (0.5,0.5) grid[step=0.2] (4,4.5);
% Zone administrative / usine (Est)
\fill[poppurple!20] (5,2) rectangle (11,7);
% Bâtiments usine
\foreach \x in {5.5,6.5,7.5,8.5,9.5} {
    \fill[poppurple!60] (\x,2.5) rectangle (\x+0.7,4.5);
}
% Bureaux
\fill[poppurple!60] (5.5,5.5) rectangle (8,6.5);
% Équipements
% École
\fill[poppink!60] (4.2,6) rectangle (4.9,7);
\node[font=\tiny, rotate=90, anchor=center] at (4.55,6.5) {École};
% Dispensaire
\fill[popred!60] (4.2,5) rectangle (4.9,5.8);
\node[font=\tiny, rotate=90, anchor=center] at (4.55,5.4) {Santé};
% Marché
\fill[popgold!60] (4.2,4) rectangle (4.9,4.8);
\node[font=\tiny, rotate=90, anchor=center] at (4.55,4.4) {Marché};
% Terrain sport
\fill[popgreen!40] (4.2,2.5) rectangle (4.9,3.8);
\node[font=\tiny, rotate=90, anchor=center] at (4.55,3.15) {Sport};

% Routes
\draw[popline, thick] (0,4) -- (5,4); % Route principale E-O
\draw[popline, thick] (4.5,0) -- (4.5,8); % Route N-S
\draw[popline, dashed] (0.5,5.5) -- (0.5,7.5); % Desserte cadres
\draw[popline, dashed] (0.5,0.5) -- (0.5,4.5); % Desserte ouvriers

% Voie ferrée (si applicable)
\draw[popdark, double distance=1.5pt] (5,1) -- (11,1);

% Légende intégrée (simplifiée pour la figure)
\begin{scope}[shift={(0.2,0.2)}, font=\tiny]
\node[anchor=west] at (0,7.5) {\textbf{Légende}};
\fill[popblue!20] (0,7.0) rectangle (0.4,7.3); \node[anchor=west] at (0.5,7.15) {Habitat cadres};
\fill[poporange!20] (0,6.6) rectangle (0.4,6.9); \node[anchor=west] at (0.5,6.75) {Habitat ouvriers};
\fill[poppurple!20] (0,6.2) rectangle (0.4,6.5); \node[anchor=west] at (0.5,6.35) {Zone admin./usine};
\fill[poppink!60] (0,5.8) rectangle (0.4,6.1); \node[anchor=west] at (0.5,5.95) {École};
\fill[popred!60] (0,5.4) rectangle (0.4,5.7); \node[anchor=west] at (0.5,5.55) {Centre de santé};
\fill[popgold!60] (0,5.0) rectangle (0.4,5.3); \node[anchor=west] at (0.5,5.15) {Marché};
\fill[popgreen!40] (0,4.6) rectangle (0.4,4.9); \node[anchor=west] at (0.5,4.75) {Terrain sport};
\draw[popline, thick] (0,4.2) -- (0.4,4.2); \node[anchor=west] at (0.5,4.2) {Route goudronnée};
\draw[popline, dashed] (0,3.8) -- (0.4,3.8); \node[anchor=west] at (0.5,3.8) {Piste/voie de desserte};
\draw[popdark, double distance=1.5pt] (0,3.4) -- (0.4,3.4); \node[anchor=west] at (0.5,3.4) {Voie ferrée};
\fill[popgreen!15] (0,3.0) rectangle (0.4,3.3); \node[anchor=west] at (0.5,3.15) {Plantation (périmètre)};
\end{scope}

% Titre
\node[font=\normalsize\bfseries, anchor=north, popdark] at (6,8.3) {Croquis : Organisation spatiale d'un camp de plantation type (ex. CDC Tiko)};
\node[font=\tiny, anchor=south, popmuted] at (6,-0.3) {Source : Exploitation d'image aérienne / Google Earth -- Réalisation : Élève Terminale};
\end{tikzpicture}
\legende{Modèle de croquis d'un camp de plantation. On distingue clairement la ségrégation spatiale historique (cadres vs ouvriers), la centralité des équipements collectifs (école, santé, marché) et l'interface avec l'unité industrielle et la plantation.}
\end{popfigure}

\begin{exoresolu}[Commentaire type du croquis (rédaction attendue)]
Énoncé : Commenter l'organisation spatiale du camp de plantation représenté par le croquis ci-dessus en mobilisant les notions d'intégration et de ségrégation.
\tcblower
\textbf{Analyse :} Le croquis met en évidence une organisation spatiale \textbf{fonctionnelle et ségrégative}, héritée de l'époque coloniale.
\begin{enumerate}
    \item \textbf{Dichotomie habitat :} Au Nord-Ouest, l'habitat « cadres » (basses densités, parcelles vastes, voirie large) s'oppose au Sud-Ouest à l'habitat « ouvriers » (haute densité, alignement rigoureux type « camps », voirie étroite). Cette ségrégation résidentielle reflète la hiérarchie salariale et sociale.
    \item \textbf{Pôle industriel et administratif :} À l'Est, l'usine de transformation (palmier/hévéa) et les bureaux concentrent l'activité productive. La voie ferrée (ou route lourde) assure l'évacuation de la production vers Douala.
    \item \textbf{Équipements collectifs intégrateurs :} École, dispensaire, marché et terrain de sport sont positionnés à l'interface des deux habitats, accessibles à tous. Ils jouent un rôle clé dans l'intégration sociale (cf. Dossier 2) et la stabilisation de la main-d'œuvre.
    \item \textbf{Emprise foncière :} Le camp s'insère dans la matrice de la plantation (hachures vertes), soulignant la mainmise foncière de la firme sur l'espace rural.
\end{enumerate}
\textbf{Conclusion :} L'espace du camp est un « village d'entreprise » autosuffisant, vecteur d'intégration par les services, mais marqueur de fractures sociales persistantes.
\end{exoresolu}

\courssub{TP 3 : Construire un diagramme en barres comparatif des superficies ou effectifs}

\begin{experience}[TP 3 : Diagramme en barres — Superficies concédées aux principales firmes (données 2022)]
\textbf{Données statistiques (ordres de grandeur, sources MINADER/Rapports annuels) :}
\begin{center}
\begin{tabular}{|l|c|c|}\hline
\rowcolor{popblueL}
\textbf{Firme} & \textbf{Superficie totale concédée (ha)} & \textbf{Superficie exploitée (ha)} \\ \hline
SOCAPALM & 58 000 & 52 000 \\
CDC & 40 000 & 35 000 \\
HEVECAM & 40 000 & 38 000 \\
SAFACAM & 18 000 & 16 500 \\
SOSUCAM & 10 000 & 9 500 \\
SODECOTON* & -- & -- \\ \hline
\end{tabular}
\end{center}
\emph{* SODECOTON n'est pas une firme de plantation mais un opérateur filière coton (égrenage, encadrement paysan). On ne lui associe pas de superficie plantationnaire propre.}\par
\textbf{Consigne :} Construire un diagramme en barres comparant les \emph{superficies exploitées} des 5 firmes plantationnières. Appliquer la méthode rigoureuse.
\end{experience}

\begin{popfigure}
\begin{tikzpicture}
\begin{axis}[
    ybar,
    bar width=12pt,
    width=\linewidth,
    height=9cm,
    enlarge x limits=0.25,
    ylabel={Superficie exploitée (1000 ha)},
    ylabel style={font=\small},
    xtick=data,
    xticklabels={SOCAPALM, CDC, HEVECAM, SAFACAM, SOSUCAM},
    xticklabel style={rotate=30, anchor=east, font=\small},
    ymin=0,
    ymax=60,
    ymajorgrids=true,
    grid style={popline},
    legend style={at={(0.5,-0.25)}, anchor=north, legend columns=-1, font=\small},
    title style={font=\normalsize\bfseries, align=center},
    title={Superficies exploitées par les grandes firmes plantationnières (2022)},
]
\addplot[popblue, fill=popblue!70] coordinates {(1,52)};
\addplot[poporange, fill=poporange!70] coordinates {(2,35)};
\addplot[popgreen, fill=popgreen!70] coordinates {(3,38)};
\addplot[poppurple, fill=poppurple!70] coordinates {(4,16.5)};
\addplot[poppink, fill=poppink!70] coordinates {(5,9.5)};
\legend{Superficie (1000 ha)}
\end{axis}
\end{tikzpicture}
\legende{Diagramme en barres modèle : titres complets, axes nommés et gradués (zéro inclus), barres proportionnelles, couleurs distinctes, source citée. Unités en milliers d'hectares pour lisibilité.}
\end{popfigure}

\begin{exoresolu}[Construction pas à pas du diagramme (réponse type élève)]
Énoncé : Construire le diagramme en barres des superficies exploitées (données du tableau ci-dessus) en respectant les règles de la sémiologie graphique.
\tcblower
\textbf{1. Choix de l'échelle :} Valeur max = 52 (SOCAPALM). Hauteur graphique dispo = 8 cm. Échelle : 1 cm pour 5 000 ha (soit 5 unités de 1000 ha). Barre max = 10.4cm $\rightarrow$ trop grand. Ajustement : 1 cm pour 10 000 ha (10 unités). Barre max = 5.2cm. Bien.\par
\textbf{2. Axes :} Abscisses : 5 barres espacées de 1 cm, largeur 0,8 cm. Ordonnées : graduation 0, 10, 20, 30, 40, 50, 60 (1000 ha).\par
\textbf{3. Barres :}
\begin{itemize}
    \item SOCAPALM : 52 $\rightarrow$ 5,2 cm (bleu)
    \item CDC : 35 $\rightarrow$ 3,5 cm (orange)
    \item HEVECAM : 38 $\rightarrow$ 3,8 cm (vert)
    \item SAFACAM : 16,5 $\rightarrow$ 1,65 cm (violet)
    \item SOSUCAM : 9,5 $\rightarrow$ 0,95 cm (rose)
\end{itemize}
\textbf{4. Titre :} \emph{Superficies exploitées par les grandes firmes plantationnières au Cameroun en 2022 (en milliers d'hectares)}.\par
\textbf{5. Source :} \emph{Source : MINADER, Annuaire statistique de l'agriculture 2022 ; Rapports annuels des entreprises.}\par
\textbf{Analyse rapide :} SOCAPALM domine largement (palmier à huile), suivie du duo CDC/HEVECAM (diversification/hévéa). SAFACAM et SOSUCAM ont des périmètres plus restreints mais à forte valeur ajoutée (huile de palme raffinée, sucre).
\end{exoresolu}

\begin{attention}[Erreurs classiques sur ce diagramme]
\begin{itemize}
    \item Mettre SODECOTON à 0 ou l'omettre sans note explicative $\rightarrow$ confusion pour le correcteur. Préciser « hors SODECOTON (filière cotonnière) ».
    \item Utiliser des largeurs de barres variables.
    \item Oublier l'unité « milliers d'hectares » sur l'axe Y (barres de 52 cm au lieu de 5,2 cm !).
    \item Ne pas citer la source (année + organisme).
\end{itemize}
\end{attention}

\courssub{TP 4 : Interpréter un tableau de données et rédiger un commentaire géographique}

\begin{experience}[TP 4 : Commentaire géographique à partir d'un tableau de synthèse (Production, Exportations, Emplois)]
\textbf{Document :} Tableau de données 2023 (valeurs approximatives, sources : INS, Rapports BCEAO, sites entreprises).
\begin{center}
\begin{tabularx}{\linewidth}{|l|X|X|X|X|}\hline
\rowcolor{popblueL}
\textbf{Firme} & \textbf{Production principale (tonnes)} & \textbf{Exportations (valeur, Milliards FCFA)} & \textbf{Effectifs permanents} & \textbf{Effectifs saisonniers (pointe)} \\ \hline
CDC & Banane 220 000 ; Caoutchouc 25 000 ; Palme 80 000 & 120 & 18 500 & 8 000 \\ \hline
SOCAPALM & Huile de palme brute 180 000 & 95 & 6 200 & 4 500 \\ \hline
HEVECAM & Caoutchouc naturel 55 000 & 65 & 6 800 & 3 000 \\ \hline
SOSUCAM & Sucre 110 000 & 45 & 3 800 & 5 000 (campagne coupe) \\ \hline
SAFACAM & Huile palme 45 000 ; Caoutchouc 12 000 & 30 & 3 200 & 1 500 \\ \hline
SODECOTON & Coton graine 250 000 (collecté) & 80 & 1 100 (cadres/tech.) & 2 500 (agents terrain) \\ \hline
\end{tabularx}
\end{center}
\textbf{Consigne :} Rédigez un commentaire géographique structuré (15-20 lignes) mettant en évidence la place de ces firmes dans l'économie nationale et l'intégration territoriale.
\end{experience}

\begin{methode}[Structure type d'un commentaire géographique de tableau statistique (Bac)]
\textbf{1. Accroche / Introduction (2 lignes) :} Nature du doc, thème, date, source, thèse principale (ex. « L'agro-industrie camerounaise, pilier de l'exportation et pourvoyeur d'emplois, repose sur quelques firmes leaders aux profils contrastés »).\par
\textbf{2. Analyse structurée (3-4 paragraphes) :}
\begin{itemize}
    \item \textbf{Dominance productive et exportatrice :} Hiérarchie des firmes (CDC leader volume/valeur), spécialisations (banane, palme, hévéa, sucre, coton), part dans les exportations hors pétrole (environ 30-40\%).
    \item \textbf{Poids social et emploi :} Volume d'emplois permanents (≈ 40 000 directs) + saisonniers (≈ 25 000). Impact sur l'insertion des jeunes, réduction pauvreté zones rurales. Distinction cadres/ouvriers.
    \item \textbf{Ancrage territorial et intégration :} Localisation (Sud/Centre vs Nord), effets d'entraînement (routes, écoles, santé), mais aussi enclaves.
    \item \textbf{Vulnérabilités / Défis :} Dépendance cours mondiaux, vieillissement plantations, pression foncière, climat.
\end{itemize}
\textbf{3. Conclusion ouverte (1 ligne) :} Perspectives (agropoles, transformation locale, certification durable).
\end{methode}

\begin{exoresolu}[Commentaire géographique modèle (niveau Bac 16-18/20)]
Énoncé : Rédiger un commentaire géographique structuré à partir du tableau ci-dessus.
\tcblower
Le tableau met en lumière le rôle structurant de six firmes agro-industrielles dans l'économie camerounaise. Premier constat : une \textbf{hiérarchie nette de la production et de l'exportation}. La CDC, par sa diversification (banane, hévéa, palme), génère la plus forte valeur exportée (120 Md FCFA), suivie par SOCAPALM (palme) et HEVECAM (hévéa). Ensemble, ces trois firmes concentrent plus des deux tiers des recettes d'exportation agricole hors pétrole, confirmant la spécialisation du Cameroun dans les cultures pérennes tropicales. Le coton (SODECOTON), bien que premier en volume collecté (250 000 t), pèse moins en valeur unitaire.

Deuxième enseignement : un \textbf{poids social majeur}. Avec près de 40 000 emplois permanents et 25 000 saisonniers, ces firmes sont les premiers employeurs privés formels du pays. La CDC alone emploie près de la moitié des effectifs permanents, agissant comme un véritable « État dans l'État » dans le Sud-Ouest. Les emplois saisonniers (coupe canne SOSUCAM, récolte coton) absorbent une main-d'œuvre locale jeune, limitant l'exode rural.

Troisièmement, l'\textbf{ancrage territorial différencié} favorise l'intégration nationale : les firmes du Sud/Centre (CDC, SOCAPALM, HEVECAM, SOSUCAM, SAFACAM) maillent l'axe Douala-Yaoundé-Garoua via les infrastructures (routes, rails, ports), tandis que SODECOTON structure le Septentrion. Les camps (cf. TP 2) concentrent équipements sociaux (santé, éducation) qui bénéficient aux populations environnantes.

Cependant, ce modèle montre des \textbf{limites} : dépendance aux cours mondiaux (coton, caoutchouc), plantations vieillissantes (besoin de rénovation), conflits fonciers avec les communautés riveraines, et nécessité de monter en gamme (transformation locale, certification RSPO/FSC). L'émergence des \textbf{agropoles} (ex. Agropole de Nkoteng pour SOSUCAM) tente d'y répondre par la contractualisation avec les planteurs villageois.
\end{exoresolu}

\courssub{Correction collective et auto-évaluation : grille critériée}

Pour progresser, vous devez être capable d'évaluer votre propre travail ou celui de votre binôme avec la même rigueur qu'un correcteur du Baccalauréat. Utilisez cette grille après chaque TP.

\begin{center}
\begin{tabularx}{\linewidth}{|>{\bfseries}l|X|X|X|}\hline
\rowcolor{popblueL}
\textbf{Critère} & \textbf{Expert (4 pts)} & \textbf{Compétent (3 pts)} & \textbf{À consolider (1-2 pts)} \\ \hline
\textbf{Titre \& Source} & Complet (thème, lieu, date, unité), source précise & Présent mais incomplet (date ou unité manquante) & Absent ou imprécis \\ \hline
\textbf{Figurés / Barres} & Choix pertinents, distincts, proportionnels, esthétiques & Choix corrects mais imprécis (proportionalité approximative) & Figurés inadaptés, non proportionnels, confus \\ \hline
\textbf{Légende / Axes} & Hiérarchisée, groupée, exhaustive, lisible & Complète mais mal hiérarchisée ou axes non nommés & Incomplète, désordonnée, absente \\ \hline
\textbf{Localisation / Données} & Exactitude toponymique, valeurs respectées & 1-2 erreurs de localisation ou de valeur & Erreurs multiples, invention de données \\ \hline
\textbf{Commentaire / Analyse} & Structuré (thèse, arguments chiffrés, vocabulaire géo), synthèse & Descriptif, quelques chiffres, vocabulaire commun & Liste de chiffres sans analyse, hors sujet \\ \hline
\textbf{Propreté / Soin} & Très soigné, lisible, couleurs/crayons maîtrisés & Correct, lisible & Brouillon, illisible, ratures \\ \hline
\end{tabularx}
\end{center}

\begin{savaistu}[Le savais-tu ? — L'épreuve pratique au Bac de Géographie]
Depuis la réforme par compétences, l'épreuve écrite de Géographie en Terminale (toutes séries) comporte \textbf{obligatoirement} un exercice de construction graphique : croquis de carte, diagramme (barres, secteurs, courbes), ou analyse de document photographique (croquis d'interprétation). Ce(s) exercice(s) valent souvent \textbf{6 à 8 points sur 20}. Les correcteurs appliquent une grille nationale proche de celle ci-dessus. \textbf{S'entraîner chaque semaine sur un TP chronométré (20-25 min) est la meilleure assurance-réussite.} Les annales des sessions 2021-2024 montrent une récurrence : localisation des agropoles/plantations, diagramme des exportations, croquis de camp ou de port (Douala/Kribi). Maîtrisez les figurés standards (voir \emph{À retenir} TP 1) et la méthode du commentaire (TP 4) : ce sont des points « gratuits » si le geste est automatisé.
\end{savaistu}

\begin{exercice}[Entraînement individuel — À rendre pour la séance suivante]
\begin{enumerate}
    \item Sur une carte muette neuve, localisez les \textbf{agropoles} en cours de développement (Nkoteng, Mbanga, Yabassi, Maroua, Ngaoundéré) avec un figuré surfacique (hachures) et les firmes associées.
    \item À partir du tableau ci-dessous (Exportations 2021-2023), construisez un \textbf{diagramme en barres groupées} (trois barres par firme : 2021, 2022, 2023) pour CDC, SOCAPALM, HEVECAM.
    \begin{center}
    \begin{tabular}{|l|c|c|c|}\hline
    \rowcolor{popblueL}
     & 2021 & 2022 & 2023 \\ \hline
    CDC & 95 & 110 & 120 \\ \hline
    SOCAPALM & 70 & 85 & 95 \\ \hline
    HEVECAM & 50 & 60 & 65 \\ \hline
    \end{tabular}
    \end{center}
    \item Rédigez en 10 lignes l'analyse de l'évolution de ces exportations (tendances, causes conjoncturelles/structurelles).
\end{enumerate}
\end{exercice}

\begin{corrige}[Exercice entraînement — Éléments de correction]
\textbf{1. Carte agropoles :} Figurés surfaciques (hachures couleur distincte par agropole) + symboles firmes. Légende : \emph{Agropoles en développement} (hachures) + \emph{Firmes leaders} (symboles TP1). Titre : \emph{Cameroun : Agropoles et firmes associées (2024)}.\par
\textbf{2. Diagramme barres groupées :} Axe X : 3 groupes (CDC, SOCAPALM, HEVECAM). Chaque groupe = 3 barres accolées (2021 bleu clair, 2022 bleu moyen, 2023 bleu foncé). Échelle Y : 0 à 130 Md FCFA. Titre : \emph{Évolution des exportations des trois principales firmes agro-industrielles (2021-2023)}. Source : INS/BCEAO.\par
\textbf{3. Analyse type :} Croissance continue pour les trois firmes (+26\% CDC, +36\% SOCAPALM, +30\% HEVECAM sur 3 ans). Causes : reprise post-Covid, hausse cours matières premières (caoutchouc, huile de palme), investissements rénovation (CDC), extension surfaces (SOCAPALM). CDC conserve le leadership. Vigilance : cours 2024 en baisse, coûts intrants (engrais) en hausse.
\end{corrige}

\newpage

\coursec{Activité d'intégration}

\courssub{Objectifs de l'activité}

Cette activité d'intégration te place dans la peau d'un géographe chargé d'éclairer un débat public réel : les grandes plantations sont-elles un moteur du développement local ou une source de problèmes ? À travers l'étude du cas de la SOCAPALM d'Édéa, tu vas mobiliser l'ensemble des acquis de la leçon. À l'issue du travail, tu seras capable de :

\begin{itemize}
\item localiser et représenter sur une carte muette un site de grande plantation (savoir-faire : observer, localiser, cartographier) ;
\item décrire méthodiquement des photographies de paysages de plantation et de camps de travailleurs (savoir-faire : observer et décrire) ;
\item transformer un tableau statistique en diagramme en barres correctement légendé (savoir-faire : représenter et schématiser) ;
\item confronter des sources contradictoires et construire un paragraphe argumenté et nuancé (savoir-faire : inventorier, classer, interpréter des documents) ;
\item mobiliser les notions-clés de la leçon : \cle{agriculture intensive}, \cle{agropole}, impact socio-économique, main-d'œuvre et \cle{intégration nationale}.
\end{itemize}

\courssub{Situation}

\textbf{Contexte.} La Société camerounaise de palmeraies (\cle{SOCAPALM}), créée en 1963, est la première firme agro-industrielle du Cameroun. Son site d'Édéa, dans le département de la Sanaga-Maritime (région du Littoral), s'étend sur les deux rives du fleuve Sanaga, à proximité de la ville d'Édéa (environ \num{120000} habitants, ordre de grandeur 2023). La palmeraie couvre environ \num{15000} hectares ; l'huilerie associée transforme les régimes de fruits sur place. L'entreprise emploie directement environ \num{3500} salariés permanents et saisonniers, venus de toutes les régions du Cameroun, et fait vivre un camp ouvrier doté d'écoles, d'un centre de santé et d'un marché.

\textbf{Le débat.} À l'occasion d'un conseil municipal élargi d'Édéa, deux positions s'affrontent. Le maire affirme : « Sans la SOCAPALM, Édéa n'aurait ni ces emplois, ni ces écoles, ni ce dynamisme commercial. » Un chef coutumier riverain répond : « La plantation a pris nos terres communautaires sans vraie compensation, et ses rejets polluent la Sanaga où nous pêchions. » Le préfet du département mandate ta classe pour produire un rapport géographique objectif.

\textbf{Le dossier documentaire remis à chaque groupe.}

\emph{Document 1 — Données statistiques du site d'Édéa (valeurs simplifiées, ordres de grandeur).}

\begin{center}
\begin{tabularx}{\linewidth}{|l|X|}
\hline
\rowcolor{popblueL} \textbf{Indicateur} & \textbf{Valeur} \\
\hline
Superficie plantée en palmiers à huile & environ \num{15000} ha \\
\hline
Effectif salarié (permanents et saisonniers) & environ \num{3500} personnes \\
\hline
Production annuelle d'huile de palme (site) & environ \num{30000} tonnes \\
\hline
Part des travailleurs originaires d'autres régions & environ 60 \% \\
\hline
Nombre d'écoles dans les camps & 4 \\
\hline
\end{tabularx}
\end{center}

\emph{Document 2 — Photographie A (description) :} vue aérienne de la palmeraie : alignements réguliers de palmiers à perte de vue, pistes rectilignes quadrillant la plantation, hangars de l'huilerie avec cheminée, au loin le fleuve Sanaga.

\emph{Document 3 — Photographie B (description) :} le camp ouvrier : rangées de maisons identiques en matériaux durables le long de rues rectilignes, un bâtiment scolaire avec drapeau, un centre de santé, un petit marché animé, des antennes paraboliques.

\emph{Document 4 — Extraits de presse contradictoires.}
\begin{itemize}
\item \textbf{Extrait 1 (quotidien national) :} « À Édéa, la SOCAPALM a scolarisé des milliers d'enfants de ses employés et offre des soins gratuits dans ses centres de santé ; ses salariés, venus du Nord comme de l'Ouest, vivent et travaillent ensemble. »
\item \textbf{Extrait 2 (journal d'investigation) :} « Des communautés riveraines dénoncent l'accaparement de leurs terres ancestrales et la pollution de la Sanaga par les effluents de l'huilerie, qui aurait fait chuter les captures de poissons. »
\end{itemize}

\textbf{Question directrice.} La SOCAPALM d'Édéa est-elle un facteur d'intégration et de développement ?

\courssub{Consignes}

Par groupes de quatre, et en deux heures, réalisez les tâches suivantes dans l'ordre.

\begin{enumerate}
\item Sur la carte muette du Cameroun fournie : situez et nommez la région du Littoral, la ville d'Édéa, le fleuve Sanaga, la ville de Douala et l'océan Atlantique. Représentez la palmeraie par une zone colorée et l'axe routier Douala--Édéa--Yaoundé par un trait. Rédigez une légende complète et placez la flèche du nord.
\item Décrivez la photographie A en trois étapes (d'ensemble, éléments, interprétation) et dégagez deux caractéristiques de l'agriculture de grandes plantations visibles sur l'image.
\item À partir de la photographie B, construisez le croquis du camp ouvrier : titre, flèche du nord, légende numérotée (au moins cinq éléments : habitations, école, centre de santé, marché, rues).
\item Transformez le document 1 en un diagramme en barres représentant les trois indicateurs chiffrés principaux (superficie, effectif, production), avec titre, axes légendés et unités. Calculez la production moyenne d'huile de palme par hectare.
\item Classez les informations des documents 4 en deux colonnes : « effets positifs » et « effets négatifs », puis reliez chaque effet à un savoir de la leçon (emplois, conditions de vie, accaparement des terres, pollution).
\item Rédigez collectivement un paragraphe argumenté de quinze à vingt lignes répondant à la question directrice : thèse annoncée, deux arguments positifs, deux arguments négatifs, conclusion nuancée proposant une piste de solution.
\item Restitution orale (cinq minutes par groupe) suivie d'un débat collectif : la classe vote ensuite la proposition de solution la plus convaincante.
\end{enumerate}

\courssub{Ressources à mobiliser}

\begin{aretenir}[Ce que tu dois avoir en tête]
\begin{itemize}
\item \cle{Grande plantation} : exploitation de plusieurs centaines ou milliers d'hectares, monoculture de rente (palmier à huile, hévéa, banane, coton, canne à sucre), \cle{agriculture intensive} (forte mise de capitaux, engrais, mécanisation, main-d'œuvre salariée nombreuse), financement par de grandes firmes souvent à capitaux mixtes.
\item Les grandes firmes du Cameroun : \cle{SOCAPALM} et \cle{SAFACAM} (palmier à huile), \cle{CDC} (banane, palmier, hévéa), \cle{HEVECAM} (hévéa), \cle{SODECOTON} (coton), \cle{SOSUCAM} (sucre) ; à distinguer des plantations des particuliers.
\item Impacts socio-économiques : positifs (emplois, salaires, écoles, centres de santé, routes, \cle{agropole}) et négatifs (accaparement des terres communautaires, pollution, dépendance alimentaire).
\item Main-d'œuvre et intégration : travailleurs venus de toutes les régions, organisation spatiale en camps (cadres / ouvriers), structures sociales (écoles, centres de santé) qui brassent les populations.
\item Savoir-faire cartographiques : titre, orientation, légende numérotée ou fléchée ; calculs : pourcentage, moyenne, rendement (production $\div$ superficie).
\end{itemize}
\end{aretenir}

\begin{attention}[Pièges à éviter]
\begin{itemize}
\item Ne pas confondre \emph{plantation de firme} (CDC, SOCAPALM...) et \emph{plantation de particulier} (exploitations villageoises de quelques hectares, moins capitalistiques).
\item Ne pas oublier l'échelle et la flèche du nord sur tout croquis ou carte (sanctionné au Bac).
\item Dans le paragraphe argumenté : pas de « je pense que », mais « les documents montrent que », « l'analyse révèle que ». La nuance est obligatoire : ni apologie, ni rejet systématique.
\item Attention aux unités dans le diagramme : superficie en ha, effectif en personnes, production en tonnes. Ne pas les mélanger sur un même axe sans le préciser.
\end{itemize}
\end{attention}

\courssub{Démarche et corrigé commenté}

\begin{exoresolu}[Rapport géographique sur la SOCAPALM d'Édéa]
Énoncé : à partir du dossier documentaire, produire la carte légendée, le croquis du camp, le diagramme statistique et le paragraphe argumenté répondant à la question : « La SOCAPALM d'Édéa est-elle un facteur d'intégration et de développement ? »
\tcblower

\textbf{Étape 1 — La carte de localisation (production attendue).}
\begin{popfigure}
\begin{center}
\includegraphics[width=\linewidth,height=9.5cm,keepaspectratio]{N01/cmr_map.jpg}
\end{center}
\legende{Fond de carte : carte générale du Cameroun (littoral, fleuves, routes et chemin de fer). Le site de la SOCAPALM d'Édéa se situe dans le Littoral, sur la Sanaga, à une soixantaine de kilomètres à l'est de Douala, le long de l'axe routier et ferroviaire qui relie le port de Douala à Yaoundé. \\ Source : Wikimedia Commons, CIA, domaine public.}
\end{popfigure}
\emph{Justification :} situer le site montre sa position stratégique (proximité du port de Douala pour l'exportation, du fleuve pour l'irrigation et le transport fluvial historique, de l'axe routier pour l'évacuation de la production).

\textbf{Étape 2 — Description de la photographie A.}
\begin{itemize}
\item \textbf{Vue d'ensemble :} une immense étendue verte plane, traversée par un fleuve large, sous un ciel dégagé. L'impression dominante est l'uniformité et l'étendue.
\item \textbf{Éléments remarquables :} palmiers plantés en lignes parfaitement régulières (quincunce), pistes rectilignes formant un quadrillage orthogonal, bâtiments industriels (huilerie) avec cheminée d'évacuation, regroupements de bâtiments (campements) en lisière.
\item \textbf{Interprétation :} on reconnaît une \cle{monoculture} intensive de rente (palmier à huile), organisée rationnellement pour la mécanisation et la surveillance, intégrant l'usine de transformation sur place — c'est le modèle de l'\cle{agropole}. Deux caractéristiques visibles : l'\textbf{alignement régulier} (témoin de l'intensivité et de la rationalisation de l'espace) et la \textbf{présence de l'unité de transformation au cœur de la plantation} (intégration amont-aval, réduction des coûts de transport des régimes périssables).
\end{itemize}

\textbf{Étape 3 — Croquis du camp ouvrier (production attendue).}
\begin{popfigure}
\begin{tikzpicture}[scale=0.9]
% Grid streets
\draw[popmuted, thick] (0,0) grid (8,5);
% Main streets wider
\draw[popdark, line width=1.5pt] (0,2.5) -- (8,2.5); % Horizontal main
\draw[popdark, line width=1.5pt] (4,0) -- (4,5); % Vertical main
% Housing blocks (Ouvriers)
\foreach \x in {0.5,1.5,2.5,3.5} {
  \foreach \y in {0.5,1.5} {
    \fill[popblueL] (\x,\y) rectangle ++(0.8,0.8);
  }
}
% Housing blocks (Cadres) - larger, spaced
\foreach \x in {5.5,6.5} {
  \foreach \y in {3.5,4.5} {
    \fill[poppurpleL] (\x,\y) rectangle ++(1.2,1.2);
  }
}
% School
\fill[popgoldL] (1,3.5) rectangle ++(1.5,1);
\node[font=\tiny, align=center] at (1.75,4) {École};
% Health Center
\fill[poppinkL] (3,3.5) rectangle ++(1.5,1);
\node[font=\tiny, align=center] at (3.75,4) {Centre\ de\ santé};
% Market
\fill[poporangeL] (5.5,1.5) rectangle ++(2,1);
\node[font=\tiny, align=center] at (6.5,2) {Marché};
% Legend
\begin{scope}[shift={(-3,0.5)}]
\draw[popline] (0,0) rectangle (2.8,4.5);
\node[font=\small\bfseries] at (1.4,4.3) {LÉGENDE};
\fill[popblueL] (0.2,3.8) rectangle (0.6,4.2) node[right, font=\small] {1. Habitations ouvriers};
\fill[poppurpleL] (0.2,3.3) rectangle (0.6,3.7) node[right, font=\small] {2. Quartier cadres};
\fill[popgoldL] (0.2,2.8) rectangle (0.6,3.2) node[right, font=\small] {3. École};
\fill[poppinkL] (0.2,2.3) rectangle (0.6,2.7) node[right, font=\small] {4. Centre de santé};
\fill[poporangeL] (0.2,1.8) rectangle (0.6,2.2) node[right, font=\small] {5. Marché};
\draw[popdark, line width=1.5pt] (0.2,1.3) -- (0.6,1.3) node[right, font=\small] {6. Rue principale};
\draw[popmuted, thick] (0.2,0.8) -- (0.6,0.8) node[right, font=\small] {7. Rue secondaire};
\node[font=\tiny] at (1.4,0.2) {Flèche Nord $\uparrow$};
\draw[->] (1.4,0.4) -- (1.4,0.6);
\end{scope}
% Title
\node[font=\small\bfseries, align=center] at (4,5.5) {Organisation spatiale et sociale du camp SOCAPALM d'Édéa};
\end{tikzpicture}
\legende{Croquis schématique du camp ouvrier : quadrillage orthogonal, ségrégation spatiale cadres/ouvriers, équipements collectifs centraux (école, santé, marché).}
\end{popfigure}
\emph{Justification :} la séparation nette entre le quartier des cadres (villas espacées) et les habitations ouvrières (alignées, denses) traduit la hiérarchie sociale propre au modèle paternaliste des grandes plantations. La centralité des équipements collectifs (école, santé, marché) illustre la prise en charge sociale par la firme.

\textbf{Étape 4 — Diagramme et calcul.}
\begin{popfigure}
\begin{tikzpicture}
\begin{axis}[
    ybar,
    width=\linewidth,
    height=8cm,
    bar width=20pt,
    enlarge x limits=0.3,
    symbolic x coords={Superficie (ha), Effectif (pers.), Production (t)},
    xtick=data,
    x tick label style={rotate=30, anchor=east, font=\small},
    ylabel={Valeurs},
    ylabel style={font=\small},
    ymin=0,
    ymax=35000,
    legend style={at={(0.5,-0.25)}, anchor=north, legend columns=-1, font=\small},
    nodes near coords,
    nodes near coords style={font=\tiny, rotate=90, anchor=west, yshift=2pt},
    popcurve/.style={popblue, fill=popblueL},
]
\addplot[popblue, fill=popblueL] coordinates {
    (Superficie (ha), 15000)
    (Effectif (pers.), 3500)
    (Production (t), 30000)
};
\legend{Indicateurs SOCAPALM Édéa}
\end{axis}
\end{tikzpicture}
\legende{Diagramme en barres des trois indicateurs principaux du site SOCAPALM d'Édéa. Attention : les unités diffèrent (ha, personnes, tonnes), la lecture comparative porte sur les ordres de grandeur.}
\end{popfigure}

\textbf{Calcul du rendement moyen :}
\[
\text{Rendement moyen} = \frac{\text{Production annuelle}}{\text{Superficie plantée}} = \frac{\num{30000}\ \text{t}}{\num{15000}\ \text{ha}} = 2\ \text{t/ha/an}.
\]
Soit 2 tonnes d'huile de palme par hectare et par an, un rendement honorable qui confirme le caractère \cle{intensif} de l'exploitation (bien que le potentiel génétique puisse atteindre 4--5 t/ha en conditions optimales).

\textbf{Étape 5 — Classement des effets (Tableau de synthèse).}

\begin{center}
\begin{tabularx}{\linewidth}{|>{\bfseries}l|X|X|}
\hline
\rowcolor{popblueL} \textbf{Nature de l'effet} & \textbf{Preuves documentaires (Doc 1, 3, 4)} & \textbf{Savoir de la leçon mobilisé} \\
\hline
\textbf{Positifs} & 
\begin{itemize}
\item \num{3500} emplois (Doc 1)
\item 4 écoles, centre de santé, marché (Doc 1, 3, Extrait 1)
\item Soins gratuits, scolarisation (Extrait 1)
\item 60 \% travailleurs d'autres régions : brassage (Doc 1, Extrait 1)
\item Dynamisme commercial local (Extrait 1)
\end{itemize}
& 
\begin{itemize}
\item Création d'emplois salariés
\item Amélioration des conditions de vie (services)
\item \cle{Intégration nationale} par la main-d'œuvre
\item Constitution d'une \cle{agropole} (pôle d'activités)
\end{itemize}
\\ \hline
\textbf{Négatifs} & 
\begin{itemize}
\item Accaparement terres communautaires (Extrait 2)
\item Pollution Sanaga par effluents huilerie (Extrait 2)
\item Baisse captures de poissons (Extrait 2)
\item Absence/défaut compensation foncière (Extrait 2)
\end{itemize}
& 
\begin{itemize}
\item \cle{Accaparement des terres} communautaires
\item Pollution industrielle (effluents)
\item Conflits d'usage de l'espace / ressources
\item Dépendance économique locale à la firme
\end{itemize}
\\ \hline
\end{tabularx}
\end{center}

\textbf{Étape 6 — Paragraphe argumenté (modèle de rédaction Bac).}

« La SOCAPALM d'Édéa apparaît d'abord comme un puissant facteur de développement local : avec environ \num{3500} emplois directs, quatre écoles et des centres de santé gratuits, elle améliore concrètement les conditions de vie des salariés et structure l'espace en véritable \cle{agropole} intégrant production et transformation. Elle constitue également un vecteur majeur d'\cle{intégration nationale}, puisque 60 \% de sa main-d'œuvre, originaire de l'ensemble des régions (Nord, Ouest, Centre...), cohabite quotidiennement dans les camps, favorisant le brassage culturel et linguistique. Cependant, cette réussite socio-économique a un coût environnemental et foncier majeur : l'extension de la palmeraie s'est faite au détriment des terres communautaires, souvent sans compensation équitable, et les rejets industriels (effluents de l'huilerie) polluent le fleuve Sanaga, pénalisant les activités halieutiques traditionnelles des riverains. La SOCAPALM d'Édéa est donc ambivalente : moteur de développement et d'intégration par l'emploi et les services, elle génère simultanément des injustices foncières et une dégradation écologique. Elle ne deviendra pleinement un facteur d'intégration durable que si l'État impose, via un cahier des charges strict, des compensations foncières transparentes aux communautés spoliées et un traitement physico-chimique rigoureux des effluents avant rejet dans le milieu naturel. »

\textbf{Étape 7 — Restitution et débat.}
Chaque groupe présente ses productions (carte, croquis, diagramme, paragraphe). Le débat collectif doit faire ressortir que la posture du géographe n'est pas de trancher « pour » ou « contre », mais de \textbf{hiérarchiser les faits}, de \textbf{spatialiser les enjeux} (amont/aval, camp/plantation, firme/riverains) et de \textbf{formuler des propositions aménagistes} (régulation foncière, normes environnementales, responsabilité sociétale des entreprises).

\courssub{Bilan de l'activité}

Cette enquête a permis de mettre en évidence trois acquis essentiels pour le Bac. D'abord, les savoirs de la leçon ne sont pas des listes à réciter : ils servent à décoder un paysage réel — la palmeraie d'Édéa illustre concrètement la \cle{monoculture intensive}, l'\cle{agropole} et le \cle{camp de plantation}. Ensuite, les savoir-faire se combinent : une même réalité géographique exige la carte (localiser), le croquis (organiser l'espace social), le diagramme (chiffrer, calculer un rendement) et le texte argumenté (juger, proposer). Enfin, l'activité révèle la posture du géographe : confronter des sources contradictoires, distinguer les échelles d'analyse (locale, nationale), et formuler un jugement nuancé accompagné de propositions concrètes. Tu es désormais prêt à analyser n'importe quelle grande plantation du Cameroun — \cle{CDC} de Tiko, \cle{HEVECAM} de Dizangué, \cle{SODECOTON} de Garoua, \cle{SOSUCAM} de Nkoteng — avec la même méthode rigoureuse.

\end{exoresolu}

\begin{savaistu}[L'agropole : un concept clé pour le Bac]
Le terme \cle{agropole} (contraction d'agro-industrie et de pôle) désigne une concentration spatiale d'activités agricoles, industrielles (transformation) et de services, gérée par une grande firme. Au Cameroun, les sites de la SOCAPALM (Édéa, Mbongo, Kienké, Dibombari, Eseka), de la CDC (Tiko, Muyuka) ou de SOSUCAM (Nkoteng, Mbandjock) sont des archétypes d'agropoles. Ils illustrent l'\textbf{intégration verticale} (de la plante au produit fini) et l'\textbf{intégration horizontale} (logement, santé, éducation, transport). Au Bac, on attend que tu sois capable d'identifier les composantes d'une agropole sur un croquis ou une photo et d'en analyser les effets d'entraînement (création d'emplois indirects, urbanisation, mais aussi dépendance à la monoculture).
\end{savaistu}

\newpage

\coursec{Méthodes et exercices résolus}

\courssub{Fiche méthode 1 : Observer, localiser et décrire une grande plantation sur un document}

\begin{methode}[Décrire et localiser une plantation à partir d'une carte ou d'une photographie]
\begin{enumerate}
\item \textbf{Identifier le document} : nature (carte, photographie aérienne, croquis), titre, source, échelle. Annoncer ce que le document montre en une phrase.
\item \textbf{Localiser} : situer la plantation dans la région, le département, par rapport aux repères voisins (ville, fleuve, route, axe ferroviaire). Exemple : << la plantation de la SOCAPALM de Mbongo se situe dans le département de la Sanaga-Maritime (région du Littoral), au nord-est d'Édéa >>.
\item \textbf{Décrire les éléments visibles} : étendue des parcelles (alignements réguliers, monoculture), usine de transformation, voies d'accès, camps de travailleurs, structures sociales (école, centre de santé).
\item \textbf{Nommer la culture et la firme} : palmiers à huile (SOCAPALM), hévéas (HEVECAM, CDC), bananiers (PHP, CDC), canne à sucre (SOSUCAM), coton (SODECOTON), thé (CTE).
\item \textbf{Conclure} : dégager le caractère de l'espace observé (espace aménagé, organisé, orienté vers l'exportation).
\end{enumerate}
\textbf{Réflexes} : toujours citer le document dans la copie (<< on observe sur le document 1 que... >>) ; ne jamais décrire sans localiser ; employer le vocabulaire géographique précis (monoculture, agro-industrie, concession).
\end{methode}

\begin{exoresolu}[Localisation et description d'une plantation (type Bac)]
\textbf{Énoncé.} À partir de vos connaissances et du croquis d'une plantation de palmiers à huile de la SOCAPALM (parcelles alignées, huilerie, camp ouvrier, piste vers Douala) : 1) Localisez deux sites de la SOCAPALM au Cameroun. 2) Décrivez trois éléments qui montrent qu'il s'agit d'une agriculture de grandes plantations. 3) Nommez la culture développée et précisez sa destination.
\tcblower
\textbf{Solution détaillée.}

1) \textbf{Localisation.} La SOCAPALM (Société camerounaise de palmeraies) exploite plusieurs sites dans le \cle{Sud forestier} du Cameroun : Mbongo et Edéa (département de la Sanaga-Maritime, région du Littoral), Dibombari (département du Moungo, région du Littoral), Apouh et Ngog-Mapubi (région du Centre), ainsi que Mondoni près de Mbanga. On retiendra au moins deux sites précis, situés par région et département.

2) \textbf{Description.} Trois éléments prouvent qu'il s'agit d'une agriculture de grandes plantations :
\begin{itemize}
\item de \textbf{vastes superficies} : les parcelles de palmiers s'étendent sur des milliers d'hectares (l'ensemble des concessions de la SOCAPALM dépasse \num{60000} ha, ordre de grandeur) ;
\item des \textbf{méthodes de travail modernes} : alignements réguliers des palmiers, monoculture spéculative, mécanisation partielle, intrants (engrais, produits phytosanitaires), main-d'œuvre salariée ;
\item un \textbf{financement capitalistique} : présence d'une huilerie (unité industrielle de transformation) sur le site et de pistes d'évacuation vers Douala, signes d'un investissement lourd de type agro-industriel.
\end{itemize}

3) \textbf{Culture et destination.} La culture développée est le \cle{palmier à huile}. Les régimes de noix palmes sont transformés sur place en huile de palme, destinée au marché national (huiles alimentaires, savonnerie) et à l'exportation.

\textbf{Conclusion.} La plantation de la SOCAPALM illustre parfaitement l'agriculture de grandes plantations : un espace vaste, aménagé et financé par une firme, intégré aux circuits commerciaux nationaux et internationaux.
\end{exoresolu}

\courssub{Fiche méthode 2 : Construire et légender un croquis de plantation}

\begin{methode}[Construire un croquis d'espace de plantation]
\begin{enumerate}
\item \textbf{Tracer le cadre} : un rectangle ou un contour simplifié ; placer la flèche du \cle{nord} et une échelle indicative (même approximative).
\item \textbf{Hiérarchiser les informations} : ne retenir que 4 à 6 éléments essentiels (parcelles cultivées, usine, camps, voies, structures sociales). Un croquis n'est pas une carte exhaustive.
\item \textbf{Choisir des figurés simples} : surface colorée ou hachurée pour les parcelles, rectangle noir pour l'usine, points ou petits carrés pour les camps, traits pleins pour les routes, traits discontinus pour les pistes.
\item \textbf{Numéroter les figurés} et rédiger une \cle{légende} ordonnée (de l'élément dominant à l'élément secondaire), placée sous le croquis.
\item \textbf{Donner un titre complet} : << Croquis de l'organisation spatiale de la plantation de ... >>.
\end{enumerate}
\textbf{Réflexes} : légende numérotée et jamais d'étiquettes écrites en travers du dessin ; propreté du trait ; proportionnalité approximative respectée (l'usine est petite, les parcelles dominent l'espace).
\end{methode}

\begin{exoresolu}[Croquis de l'organisation spatiale d'une plantation (type Bac)]
\textbf{Énoncé.} À l'aide d'un croquis soigneusement légendé, montrez comment l'espace d'une grande plantation agro-industrielle (exemple : HEVECAM à Niete, région du Sud) est organisé.
\tcblower
\textbf{Solution détaillée.}

\textbf{Démarche.} On organise le croquis autour de trois espaces : l'espace productif (parcelles d'hévéas), l'espace industriel (usine de transformation du latex), l'espace social (camps, école, centre de santé), reliés par les voies de communication.

\begin{popfigure}
\begin{center}
\begin{tikzpicture}[scale=0.9]
\draw[->] (9.2,4.6) -- (9.2,5.6) node[above]{\footnotesize N};
\fill[popgreenL] (0,0) rectangle (8,4);
\draw[poppurple,thick] (0.3,0.4) rectangle (3.2,3.6);
\draw[poppurple,thick] (3.6,0.4) rectangle (6.4,3.6);
\draw[poppurple,thick] (6.8,0.4) rectangle (7.8,3.6);
\node[popdark] at (1.75,2) {\footnotesize 1};
\node[popdark] at (5,2) {\footnotesize 1};
\node[popdark] at (7.3,2) {\footnotesize 1};
\fill[popink] (4.6,3.1) rectangle (5.4,3.5);
\node[popdark] at (5,2.85) {\footnotesize 2};
\fill[poporange] (1.0,3.2) circle (0.09);
\fill[poporange] (1.4,3.2) circle (0.09);
\fill[poporange] (1.8,3.2) circle (0.09);
\node[popdark] at (1.4,2.95) {\footnotesize 3};
\draw[popblue,thick] (0.9,3.45) rectangle (1.3,3.6);
\node[popdark] at (0.75,3.5) {\footnotesize 4};
\draw[popdark,thick] (0,0.15) -- (8,0.15);
\node[popdark] at (6.9,0.35) {\footnotesize 5};
\draw[popdark,thick,dashed] (5,3.5) -- (5,3.9) -- (8,3.9);
\node[popdark] at (7.2,4.1) {\footnotesize 6};
\end{tikzpicture}
\end{center}
\legende{Croquis de l'organisation spatiale d'une plantation d'hévéas (type HEVECAM, Niete). 1 : parcelles d'hévéas en monoculture ; 2 : usine de transformation du latex ; 3 : camp des ouvriers ; 4 : structure sociale (école, centre de santé) ; 5 : route principale ; 6 : piste d'évacuation vers Kribi. Échelle indicative : la plantation couvre plusieurs dizaines de milliers d'hectares.}
\end{popfigure}

\textbf{Commentaire.} Le croquis montre que la plantation est un espace \cle{total} : production, transformation et vie sociale y sont intégrées. Les parcelles dominent l'espace (monoculture sur grande superficie), l'usine est placée au cœur du site pour réduire les coûts de transport du latex, et les camps avec leurs structures sociales fixent la main-d'œuvre sur place.

\textbf{Conclusion.} Ce croquis met en évidence l'organisation rationnelle et hiérarchisée de l'espace planté, caractéristique de l'agriculture de grandes plantations.
\end{exoresolu}

\courssub{Fiche méthode 3 : Exploiter un tableau statistique (calculs de parts et de taux de variation)}

\begin{methode}[Calculer une part en pourcentage et un taux de variation]
\begin{enumerate}
\item \textbf{Lire le tableau} : titre, unité (tonnes, hectares, milliers de FCFA), source, dates. Recopier les données utiles avec leur unité.
\item \textbf{Calculer une part} :
\[ \text{part (\%)} = \frac{\text{valeur de la partie}}{\text{valeur totale}} \times 100 \]
Vérifier que la somme des parts fait environ 100 \%.
\item \textbf{Calculer un taux de variation} entre deux dates :
\[ \text{taux (\%)} = \frac{V_{\text{finale}} - V_{\text{initiale}}}{V_{\text{initiale}}} \times 100 \]
Un résultat positif = augmentation ; négatif = baisse.
\item \textbf{Rédiger la phrase d'analyse} : << la part de X est passée de ... à ... , soit une hausse de ... points / un taux de variation de ... \% >>. Ne pas confondre \emph{points de pourcentage} et \emph{pourcentage de variation}.
\item \textbf{Interpréter géographiquement} : relier le chiffre à la réalité (extension des palmeraies, hausse de la production d'huile de palme, politique de substitution aux importations).
\end{enumerate}
\end{methode}

\begin{exoresolu}[Exploitation d'un tableau statistique sur les palmeraies (type Bac)]
\textbf{Énoncé.} Le tableau ci-dessous donne les superficies approximatives de quelques firmes de plantations au Cameroun (ordres de grandeur).

\begin{center}
\begin{tabularx}{\linewidth}{|l|X|X|}\hline
\rowcolor{popblueL} \textbf{Firme} & \textbf{Culture principale} & \textbf{Superficie (ha, ordre de grandeur)} \\ \hline
SOCAPALM & Palmier à huile & \num{70000} \\ \hline
HEVECAM & Hévéa & \num{40000} \\ \hline
SOSUCAM & Canne à sucre & \num{25000} \\ \hline
CDC & Banane, hévéa, palme & \num{40000} \\ \hline
\end{tabularx}
\end{center}

1) Calculez la part de la SOCAPALM dans la superficie totale de ces quatre firmes. 2) La production nationale d'huile de palme est passée d'environ \num{230000} tonnes à environ \num{350000} tonnes entre 2010 et 2020 (ordres de grandeur) : calculez le taux de variation. 3) Tirez une conclusion géographique.
\tcblower
\textbf{Solution détaillée.}

1) \textbf{Part de la SOCAPALM.} Superficie totale :
\[ \num{70000} + \num{40000} + \num{25000} + \num{40000} = \num{175000} \text{ ha} \]
\[ \text{part} = \frac{\num{70000}}{\num{175000}} \times 100 = 40 \% \]
La SOCAPALM occupe à elle seule \textbf{40 \%} de la superficie totale des quatre firmes : c'est le premier planteur agro-industriel du groupe étudié.

2) \textbf{Taux de variation de la production d'huile de palme.}
\[ \text{taux} = \frac{\num{350000} - \num{230000}}{\num{230000}} \times 100 = \frac{\num{120000}}{\num{230000}} \times 100 \approx 52,2 \% \]
La production a augmenté d'environ \textbf{52 \%} en dix ans (soit + \num{120000} tonnes).

3) \textbf{Conclusion géographique.} La forte part de la SOCAPALM et la progression de plus de la moitié de la production d'huile de palme montrent le dynamisme des plantations agro-industrielles du Sud forestier camerounais, soutenu par l'extension des superficies, l'investissement des firmes et la demande nationale croissante. Ce dynamisme confirme le rôle moteur des grandes plantations dans l'économie camerounaise, tout en posant le problème de l'accaparement des terres communautaires.
\end{exoresolu}

\courssub{Fiche méthode 4 : Cartographier les grandes plantations du Cameroun}

\begin{methode}[Réaliser une carte schématique de localisation des plantations]
\begin{enumerate}
\item \textbf{Tracer le contour simplifié} du Cameroun en polygone (forme triangulaire caractéristique, allongée vers le lac Tchad).
\item \textbf{Placer les repères} : flèche du nord, capitales (Yaoundé, Douala), ports (Douala, Kribi), océan.
\item \textbf{Localiser les firmes par des points ou des zones nommées} : CDC (région du Sud-Ouest : Tiko, Buea), PHP (Moungo : Penja, Njombé), SOCAPALM (Littoral et Centre : Mbongo, Dibombari), HEVECAM (Sud : Niete), SOSUCAM (Centre : Mbandjock, Nkoteng), SODECOTON (Nord : Garoua, Maroua, Touboro).
\item \textbf{Dégager la loi de répartition} : les plantations pérennes industrielles (palme, hévéa, banane) se concentrent dans le \cle{Sud forestier} humide ; le coton (SODECOTON) domine dans le \cle{Nord soudanien}.
\item \textbf{Légender et titrer} : figurés numérotés, titre précis, source.
\end{enumerate}
\textbf{Réflexe} : une carte sans nord, sans légende ou sans titre est sanctionnée ; la localisation doit être juste (HEVECAM n'est pas dans le Nord !).
\end{methode}

\begin{exoresolu}[Carte de localisation des grandes plantations (type Bac)]
\textbf{Énoncé.} Sur un fond de carte schématique du Cameroun, localisez au moins cinq grandes firmes de plantations et dégagez la loi de répartition spatiale qui s'en dégage.
\tcblower
\textbf{Solution détaillée.}

\begin{popfigure}
\begin{center}
\includegraphics[width=\linewidth,height=11cm,keepaspectratio]{N05/cmr_regions.jpg}
\end{center}
\legende{Fond de carte : le Cameroun en grands ensembles (carte en anglais), avec villes, routes, voies ferrées et fleuves. La solution situe : 1. la CDC au Sud-Ouest (Tiko, au pied du mont Cameroun, près de Limbe) ; 2. PHP-banane dans le Moungo (Penja--Njombé, entre Douala et Nkongsamba) ; 3. la SOCAPALM dans le Littoral ; 4. HEVECAM dans le Sud (Niete) ; 5. SOSUCAM dans le Centre (Mbandjock--Nkoteng, au nord de Yaoundé) ; 6. SODECOTON dans le Nord (Garoua et bassin cotonier). Les localisations de Tiko, Njombé, Niete, Mbandjock et Nkoteng sont à reporter à partir du texte. \\ Source : Wikimedia Commons, Peter Fitzgerald (modifié par Burmesedays), CC BY 3.0.}
\end{popfigure}

\textbf{Analyse.} La carte révèle une loi de répartition nette : cinq des six firmes se concentrent dans le \cle{Sud forestier} (Littoral, Sud-Ouest, Centre, Sud), où le climat équatorial humide convient aux cultures pérennes industrielles (banane, palme, hévéa, canne) et où les ports de Douala et Kribi facilitent l'exportation. Seule la SODECOTON occupe le \cle{Nord soudanien}, adapté au coton, culture annuelle de saison sèche.

\textbf{Conclusion.} La répartition des grandes plantations épouse les grandes zones agro-écologiques du Cameroun et la proximité des débouchés portuaires.
\end{exoresolu}

\courssub{Fiche méthode 5 : Rédiger un commentaire de document sur l'impact socio-économique des plantations}

\begin{methode}[Commenter un texte ou un tableau sur les impacts des plantations]
\begin{enumerate}
\item \textbf{Introduction} : présenter le document (nature, auteur, source, date, thème) et annoncer la problématique : << en quoi les grandes plantations contribuent-elles au développement tout en posant des problèmes ? >>.
\item \textbf{Dégager les idées principales} en deux ou trois parties équilibrées : d'une part les effets positifs (emplois salariés, revenus, structures sociales : écoles, centres de santé, routes ; devises d'exportation), d'autre part les effets négatifs (accaparement des terres communautaires, pollution par les intrants et les rejets d'usine, faiblesse des salaires, conflits fonciers).
\item \textbf{Appuyer chaque idée par un exemple chiffré} tiré du document ou des connaissances (la SODECOTON encadre des centaines de milliers de cotonculteurs ; la SOSUCAM emploie plusieurs milliers de salariés à Mbandjock et Nkoteng).
\item \textbf{Conclusion} : bilan nuancé et ouverture (agropoles, redistributions foncières, agriculture intensive durable).
\end{enumerate}
\textbf{Réflexes} : annoncer le plan ; une idée = un exemple ; ne jamais recopier le document ; rédiger en phrases complètes.
\end{methode}

\begin{exoresolu}[Commentaire de document (type Bac)]
\textbf{Énoncé.} Document : << À Mbandjock, l'installation de la SOSUCAM a transformé la région : routes, écoles, hôpital, milliers d'emplois. Mais les paysans dénoncent la perte de leurs terres et la pollution de la Sanaga par les rejets de l'usine. >> À partir de ce document et de vos connaissances, montrez que les grandes plantations ont un impact socio-économique ambivalent au Cameroun. (Environ 20 lignes.)
\tcblower
\textbf{Solution détaillée.}

\textbf{Introduction.} Le document proposé est un extrait de reportage consacré à la SOSUCAM (Société sucrière du Cameroun), installée à Mbandjock et Nkoteng (région du Centre). Il souligne le caractère ambivalent des grandes plantations : en quoi constituent-elles à la fois un facteur de développement et une source de tensions ?

\textbf{Développement.}

D'une part, les grandes plantations sont de puissants \cle{facteurs de développement}. Elles créent des emplois salariés nombreux et stables : la SOSUCAM emploie plusieurs milliers de travailleurs, la SODECOTON encadre des centaines de milliers de cotonculteurs dans le Nord (ordre de grandeur : plus de \num{300000} producteurs). Elles améliorent les conditions de vie par la construction de structures sociales (écoles, centres de santé, camps équipés en eau et électricité) et d'infrastructures (routes, ponts). Elles rapportent enfin des devises : la banane, l'huile de palme, le caoutchouc et le coton figurent parmi les principales exportations agricoles du pays.

D'autre part, ces plantations posent de \cle{graves problèmes}. L'accaparement des terres communautaires prive les villageois de leurs terres agricoles et alimente des conflits fonciers, comme autour des palmeraies de la SOCAPALM. La pollution est réelle : rejets des usines dans les cours d'eau (la Sanaga pour la SOSUCAM), usage massif d'engrais et de pesticides. S'ajoutent la faiblesse relative des salaires ouvriers et la dépendance alimentaire des populations expropriées.

\textbf{Conclusion.} Les grandes plantations apparaissent donc comme un outil de développement ambivalent : richesse nationale et intégration des territoires, mais au prix de tensions foncières et environnementales. Le projet d'\cle{agropoles}, combinant agriculture intensive et respect des communautés, tente aujourd'hui de concilier ces deux exigences.
\end{exoresolu}

\courssub{Fiche méthode 6 : Traiter le dossier 2 — la main-d'œuvre, vecteur d'intégration nationale}

\begin{methode}[Analyser la main-d'œuvre des plantations comme vecteur d'intégration]
\begin{enumerate}
\item \textbf{Identifier l'origine des travailleurs} : les plantations attirent des migrants de toutes les régions (Nord, Ouest, Nord-Ouest vers le Littoral et le Sud) : la plantation est un creuset de brassage ethnique et linguistique.
\item \textbf{Décrire l'organisation spatiale et sociale} : distinguer le \cle{camp des cadres} (villas, quartier résidentiel) et le \cle{camp des ouvriers} (logements collectifs alignés) ; repérer les structures sociales (écoles, centres de santé, terrains de sport, marchés).
\item \textbf{Expliquer le mécanisme d'intégration} : salaires réguliers, vie commune, scolarisation des enfants, soins, mariages mixtes, diffusion de la langue et des modes de vie : la plantation tisse le lien national.
\item \textbf{Nuancer} : hiérarchie sociale marquée entre cadres et ouvriers, conditions de travail parfois dures, dépendance au camp.
\item \textbf{Structurer la réponse} : introduction (situation), deux ou trois paragraphes (origines ; organisation ; intégration), conclusion.
\end{enumerate}
\end{methode}

\begin{exoresolu}[Étude de cas : la main-d'œuvre d'une plantation (dossier 2, type Bac)]
\textbf{Énoncé.} À partir de l'exemple d'une grande plantation de votre choix, montrez que la main-d'œuvre des grandes plantations est un vecteur d'intégration nationale au Cameroun.
\tcblower
\textbf{Solution détaillée.}

\textbf{Choix de l'exemple.} Prenons la plantation d'HEVECAM à Niete (région du Sud), l'une des plus grandes hévéacultures d'Afrique centrale (environ \num{40000} ha, ordre de grandeur).

\textbf{1. Une main-d'œuvre d'origines diverses.} Les milliers de travailleurs d'HEVECAM ne viennent pas seulement du département de l'Océan : ils proviennent de tout le Cameroun — régions du Nord (main-d'œuvre historique des plantations), de l'Ouest, du Nord-Ouest et du Centre. La plantation fonctionne ainsi comme un pôle migratoire qui mélange ethnies, langues et religions sur un même espace.

\textbf{2. Une organisation spatiale et sociale structurée.} L'espace du camp est hiérarchisé : le camp des cadres (villas individuelles, quartier calme) se distingue du camp des ouvriers (rangées de logements collectifs). Autour d'eux, la firme a édifié des structures sociales : écoles pour les enfants des travailleurs, centre de santé, points d'eau, terrains de sport et marchés. La plantation est donc un véritable village agro-industriel.

\textbf{3. Un vecteur d'intégration nationale.} Par le salariat régulier, la scolarisation commune des enfants, l'accès aux soins, les mariages entre personnes de régions différentes et la pratique des langues véhiculaires, la plantation fabrique une communauté de vie nationale : elle arrache les individus à leur seul terroir d'origine et les insère dans l'économie monétaire et l'espace national. C'est le même phénomène qu'observent la SODECOTON dans le bassin cotonier ou la SOSUCAM à Mbandjock.

\textbf{Nuances.} Cette intégration reste perfectible : forte hiérarchie sociale entre cadres et ouvriers, salaires modestes, dépendance totale au camp.

\textbf{Conclusion.} Par le brassage des populations qu'elle provoque et les structures sociales qu'elle installe, la grande plantation apparaît comme un véritable creuset d'intégration nationale, tout en reproduisant des inégalités sociales qu'il reste à réduire.
\end{exoresolu}

\begin{attention}[Les 5 erreurs qui coûtent des points au Bac]
\begin{enumerate}
\item \textbf{Confondre grande plantation et exploitation villageoise.} Une grande plantation se définit par trois critères cumulatifs : vastes superficies, méthodes modernes (monoculture, intrants, salariat) et financement capitalistique (firme). Un champ familial de 2 ha de palmiers n'est pas une grande plantation, même s'il est productif.
\item \textbf{Mal localiser les firmes.} Erreurs classiques : placer la SODECOTON dans le Sud ou HEVECAM dans le Nord. À retenir : CDC et PHP (Sud-Ouest et Moungo), SOCAPALM (Littoral-Centre), HEVECAM (Sud : Niete), SOSUCAM (Centre : Mbandjock-Nkoteng), SODECOTON (Nord : Garoua). Toute localisation fausse sur une carte est lourdement sanctionnée.
\item \textbf{Confondre points de pourcentage et taux de variation.} Si une part passe de 20 \% à 30 \%, la hausse est de 10 \emph{points}, mais le taux de variation est de $\frac{30-20}{20}\times100 = 50\%$. Les deux formulations ne sont pas interchangeables.
\item \textbf{Oublier les éléments obligatoires d'une carte ou d'un croquis} : titre complet, flèche du nord, légende numérotée et ordonnée, échelle indicative. Un dessin sans légende, même beau, perd la moitié des points du savoir-faire cartographique.
\item \textbf{Présenter un bilan à sens unique.} Sur l'impact socio-économique ou le dossier 2, un devoir qui ne cite que les bienfaits (ou que les méfaits) est incomplet : le correcteur attend un traitement nuancé (emplois et structures sociales \emph{contre} accaparement foncier et pollution), appuyé sur des exemples précis et des ordres de grandeur.
\end{enumerate}
\end{attention}

\newpage

\coursec{Exercices}

\courssub{Je vérifie mes connaissances}

\begin{exercice}[QCM : Caractéristiques et firmes]
Parmi les affirmations suivantes, une seule est exacte. Recopie-la en justifiant ton choix.
\begin{enumerate}
\item L'agriculture de grandes plantations se caractérise par une polyculture vivrière destinée à l'autoconsommation.
\item La CDC (Cameroon Development Corporation) est une firme privée spécialisée uniquement dans la culture du coton dans le Nord.
\item Les grandes plantations sont des exploitations capitalistiques, monocoles, tournées vers l'exportation et employant une main-d'œuvre salariée nombreuse.
\item L'agropole de Kousseri est le principal pôle de transformation de l'hévéa au Cameroun.
\end{enumerate}
\end{exercice}

\begin{exercice}[Vrai ou Faux justifié]
Détermine si chacune des propositions ci-dessous est vraie ou fausse. Justifie ta réponse par une phrase précise.
\begin{enumerate}
\item La SODECOTON gère des plantations de canne à sucre dans la région du Littoral.
\item Les plantations de particuliers (planteurs villageois) sont encadrées par des sociétés d'État comme la SODECOTON pour le coton.
\item L'agriculture intensive pratiquée dans les plantations implique de faibles rendements à l'hectare.
\item Les camps de plantation (ex. : camps de la CDC à Tiko) illustrent une ségrégation spatiale entre cadres et ouvriers.
\end{enumerate}
\end{exercice}

\begin{exercice}[Définitions et concepts clés]
Donne une définition géographique précise (deux à trois lignes) pour chacun des termes suivants :
\begin{enumerate}
\item Agriculture de grandes plantations
\item Agriculture intensive
\item Agropole
\item Accaparement des terres
\end{enumerate}
\end{exercice}

\begin{exercice}[Calculs directs et lecture de données]
Le tableau ci-après donne des données approximatives pour quatre grandes firmes camerounaises (année de référence récente).
\begin{center}
\begin{tabularx}{\linewidth}{|l|X|X|X|}\hline
\textbf{Firme} & \textbf{Superficie cultivée (ha)} & \textbf{Effectif total (employés)} & \textbf{Production principale (tonnes)} \\ \hline
CDC & 40\,000 & 22\,000 & Banane : 250\,000 ; Palmier : 120\,000 \\ \hline
SOCAPALM & 58\,000 & 6\,500 & Palmier à huile : 300\,000 \\ \hline
SOSUCAM & 14\,000 & 4\,800 & Sucre : 150\,000 \\ \hline
HEVECAM & 42\,000 & 7\,200 & Hévéa (caoutchouc) : 60\,000 \\ \hline
\end{tabularx}
\end{center}
\begin{enumerate}
\item Calcule le rendement moyen à l'hectare (en tonnes/ha) pour la banane à la CDC et pour le sucre à la SOSUCAM.
\item Quelle firme a le ratio \og effectifs / superficie \fg{} le plus élevé ? Interprète ce résultat en termes d'intensité de main-d'œuvre.
\item Si la CDC verse en moyenne 150\,000 FCFA/mois par employé, estime sa masse salariale mensuelle totale (en milliards de FCFA).
\end{enumerate}
\end{exercice}

\courssub{Je m'entraîne}

\begin{exercice}[Localisation et analyse spatiale : Carte muette]
À partir d'une carte muette du Cameroun (fournie ou à tracer schématiquement) :
\begin{enumerate}
\item Localise et nomme les six firmes suivantes : CDC, SOCAPALM, SAFACAM, SODECOTON, SOSUCAM, HEVECAM.
\item Légende leur culture principale par un symbole de couleur distinct (bananier, palmier à huile, coton, canne à sucre, hévéa).
\item Trace les principaux axes d'évacuation (routes, chemin de fer, port de Douala, port de Kribi).
\item Rédige un commentaire organisé (10 lignes) sur la répartition spatiale de ces plantations (facteurs physiques, historiques, infrastructures).
\end{enumerate}
\end{exercice}

\begin{exercice}[Croquis d'organisation sociale : Le camp de plantation]
À partir de la description suivante d'un camp type de la CDC (ex. : camp de Tiko ou Muyuka), réalise un croquis légendé au 1/5000e environ montrant l'organisation spatiale et sociale.
\textit{Description :} Le camp est quadrillé. Au centre, le \og Camp des cadres \fg{} (logements spacieux, électricité, eau courante, club-house, terrain de tennis). En périphérie, les \og Camps des ouvriers \fg{} (cases alignées, collectifs, latrines communes, points d'eau). Entre les deux, les équipements collectifs : école primaire, dispensaire, marché, lieu de culte. La voie ferrée ou la route principale traverse le camp pour évacuer la production.
\begin{enumerate}
\item Trace le contour, les zones d'habitat (cadres vs ouvriers), les équipements, la voie de communication.
\item Réalise une légende hiérarchisée (habitat, équipements, transports).
\item Titre ton croquis : \og Organisation spatiale et sociale d'un camp de plantation (ex. CDC) \fg.
\end{enumerate}
\end{exercice}

\begin{exercice}[Diagramme et interprétation : Superficies vs Emplois]
À partir du tableau de l'exercice 4 (partie \og Je vérifie \fg), construis un diagramme en barres groupées (ou superposées) comparant pour les quatre firmes : la superficie cultivée (en milliers d'hectares) et l'effectif total (en milliers d'employés).
\begin{enumerate}
\item Choisis une échelle adaptée, trace les axes, représente les barres, légende.
\item Tire deux enseignements géographiques de ce graphique (ex. : rapport capital/travail, spécificité de la SODECOTON si on l'ajoutait avec ses 200\,000 planteurs encadrés).
\end{enumerate}
\end{exercice}

\begin{exercice}[Étude de cas : La filière coton et la SODECOTON]
Lis le document ci-dessous et réponds aux questions.
\textit{Document :} La SODECOTON (Société de Développement du Coton) opère dans les régions du Nord (Adamaoua, Nord, Extrême-Nord). Elle ne possède pas de \og plantations \fg{} au sens strict (vastes domaines d'un seul tenant), mais encadre environ 200\,000 à 250\,000 producteurs villageois (cotonniers) sur des exploitations familiales moyennes de 2 à 5 ha. Elle fournit intrants (semences, engrais, pesticides), conseil technique, et assure la collecte et l'égrenage dans ses usines (Garoua, Maroua, Kaele). Le coton est la première source de revenus monétaires pour des millions de ruraux du Nord. Cependant, l'usage intensif d'intrants chimiques pose des problèmes de dégradation des sols et de pollution des nappes.
\begin{enumerate}
\item Pourquoi dit-on que la SODECOTON relève de l'agriculture de \og grandes plantations \fg{} au sens large (ou de plantation-entreprise) malgré l'absence de grands domaines unis ?
\item Cite trois avantages socio-économiques de cette filière pour l'intégration nationale (aménagement du territoire, revenus, infrastructures).
\item Cite deux risques environnementaux liés au modèle intensif cotonnier.
\item Explique en quoi l'organisation spatiale de la filière (usines d'égrenage, pistes rurales, axe routier Nord-Douala) structure le territoire septentrional.
\end{enumerate}
\end{exercice}

\courssub{Je me prépare au Bac}

\begin{exercice}[Sujet type Bac : Définition, exemples et impacts]
\textbf{Sujet :} \og Définissez l'agriculture de grandes plantations et présentez ses caractéristiques principales. Appuyez votre réponse sur l'étude de deux firmes camerounaises de votre choix. Montrez ensuite les impacts socio-économiques positifs et négatifs de ces plantations sur les territoires concernés. \fg
\begin{enumerate}
\item Propose une introduction (accroche, définition, annonce du plan).
\item Partie 1 : Caractéristiques (superficies, capital, main-d'œuvre, monoculture, exportation) illustrées par la CDC (banane/palmier/hévéa, Sud-Ouest) et la SOSUCAM (canne à sucre, Moungo/Littoral).
\item Partie 2 : Impacts positifs (emplois stables, salaires, équipements sociaux -- écoles, santé --, routes, recettes d'exportation, formation professionnelle) et impacts négatifs (accaparement terres communautaires, pollution chimique/effluents, dépendance mono-culturelle, conflits fonciers, conditions de travail pénibles, précarité des contractuels).
\item Conclusion nuancée et ouverture sur l'agropole ou la diversification.
\end{enumerate}
\end{exercice}

\begin{exercice}[Sujet type Bac : Croquis commenté et main-d'œuvre]
\textbf{Sujet :} \og La main-d'œuvre dans les grandes plantations camerounaises : un vecteur d'intégration nationale. \fg
À partir de tes connaissances et du dossier ci-après, traite les questions.
\textit{Dossiel :} Les grandes plantations (CDC, HEVECAM, SOCAPALM, SOSUCAM, SAFACAM) emploient au total plus de 50\,000 travailleurs permanents et des dizaines de milliers de saisonniers. Ces travailleurs proviennent de toutes les régions du Cameroun (Nord, Ouest, Centre, Sud, etc.) et de pays voisins (Nigéria, Tchad, RCA). Ils vivent dans des camps où cohabitent des ethnies et langues diverses. Les firmes gèrent des écoles (primaire, secondaire), des centres de santé, des réseaux d'eau/électricité. Le pidgin-english et le français y sont des langues de communication interethnique. Cependant, persiste une ségrégation spatiale (camps cadres vs camps ouvriers) et sociale (écarts de salaires, statuts précaires).
\begin{enumerate}
\item Réalise un croquis légendé intitulé \og Organisation spatiale et sociale d'un camp de plantation \fg faisant apparaître : habitat cadres, habitat ouvriers, équipements collectifs (école, santé, marché, culte), voies de circulation, zone de stockage/usine. (Utilise une feuille blanche, cadre, échelle, orientation, légende organisée).
\item Rédige un commentaire structuré (15-20 lignes) montrant en quoi la main-d'œuvre des plantations constitue un vecteur d'intégration nationale (brassage ethnique, langues communes, services sociaux partagés, mobilité géographique) tout en soulignant les limites (ségrégation spatiale, précarité, tensions foncières avec les riverains).
\end{enumerate}
\end{exercice}

\begin{exercice}[Sujet type Bac : Dissertation argumentée]
\textbf{Sujet :} \og Les grandes plantations au Cameroun : atout pour l'intégration nationale ou source de tensions ? \fg
Tu traiteras ce sujet sous forme de dissertation géographique structurée (Introduction, Deux parties équilibrées, Conclusion).
\begin{enumerate}
\item \textbf{Introduction} : Accroche (poids économique), définition du sujet, analyse des termes (\og atout \fg, \og intégration \fg, \og tensions \fg), problématique, annonce du plan.
\item \textbf{Première partie : Un atout majeur pour l'intégration nationale et le développement des territoires.} (Emplois massifs et brassage populations ; infrastructures sociales/sanitaires/éducatives ; ouverture des zones enclavées -- routes, rails, ports -- ; recettes budgétaires et devises ; formation technique ; dynamiques urbaines autour des camps).
\item \textbf{Deuxième partie : Des sources de tensions foncières, sociales et environnementales.} (Accaparement terres coutumières vs titres fonciers -- ex. conflits CDC/riverains, HEVECAM/Bakossi -- ; pollution chimique (effluents SOSUCAM, pesticides coton) ; précarité emploi (contractuels, saisonniers) ; ségrégation spatiale héritée (camps cadres/ouvriers) ; dépendance économique aux cours mondiaux ; conflits sociaux grèves).
\item \textbf{Conclusion} : Bilan synthétique, réponse nuancée à la problématique, ouverture (réformes foncières, RSE -- responsabilité sociétale des entreprises --, agropole, diversification, certification durable).
\end{enumerate}
\end{exercice}

\newpage

\coursec{Corrigés des exercices}

\begin{corrige}[Exercice 1 — QCM : Caractéristiques et firmes]
\textbf{Bonne réponse : proposition 3.}

\og Les grandes plantations sont des exploitations capitalistiques, monocoles, tournées vers l'exportation et employant une main-d'œuvre salariée nombreuse. \fg

\textbf{Justification :} c'est la définition même de l'agriculture de grandes plantations : de vastes superficies (plusieurs milliers à plusieurs dizaines de milliers d'hectares), un financement capitalistique (capitaux nationaux ou étrangers), une monoculture de rente (banane, palmier à huile, hévéa, canne à sucre, coton) destinée essentiellement à l'exportation, et une main-d'œuvre salariée nombreuse (permanents et saisonniers).

\textbf{Pourquoi les autres propositions sont fausses :}
\begin{itemize}
\item \textbf{Proposition 1 :} fausse. La polyculture vivrière destinée à l'autoconsommation caractérise l'agriculture traditionnelle (vivrière), pas la plantation, qui est monocole et commerciale.
\item \textbf{Proposition 2 :} fausse. La CDC (Cameroon Development Corporation) est une société à capitaux majoritairement publics (l'État camerounais détient la majorité du capital) ; elle produit de la banane, du palmier à huile et de l'hévéa dans le Sud-Ouest et le Littoral — et non du coton dans le Nord (c'est le domaine de la SODECOTON).
\item \textbf{Proposition 4 :} fausse. Kousseri (Extrême-Nord) n'est pas un pôle de transformation de l'hévéa ; l'hévéa est transformé notamment par HEVECAM (région du Sud, usine de Niete) et la CDC. Kousseri est plutôt associée aux filières céréalières et au commerce transfrontalier avec le Tchad.
\end{itemize}
\end{corrige}

\begin{attention}[Points de vigilance]
Ne pas confondre \cle{firmes} et \cle{cultures} : associer mentalement chaque sigle à sa culture et à sa région (CDC = banane/palmier/hévéa, Sud-Ouest et Littoral ; SODECOTON = coton, Nord ; SOSUCAM = canne à sucre, Centre : Nkoteng et Mbandjock ; HEVECAM = hévéa, Sud ; SOCAPALM et SAFACAM = palmier à huile, Littoral et Sud). C'est une erreur classique au Bac.
\end{attention}

\begin{corrige}[Exercice 2 — Vrai ou Faux justifié]
\begin{enumerate}
\item \textbf{Faux.} La SODECOTON gère la filière \cle{coton} dans les régions septentrionales (Adamaoua, Nord, Extrême-Nord) ; la canne à sucre relève de la SOSUCAM, implantée à Nkoteng et Mbandjock (département de la Haute-Sanaga, région du Centre).
\item \textbf{Vrai.} Les planteurs villageois (environ \num{200000} à \num{250000} cotonculteurs, ordre de grandeur) sont encadrés par la SODECOTON qui fournit intrants, conseil technique et crédit, et assure la collecte et l'égrenage : c'est le modèle de la \cle{plantation-entreprise} avec encadrement des paysans.
\item \textbf{Faux.} L'\cle{agriculture intensive} vise précisément des rendements élevés à l'hectare grâce à un fort investissement par unité de surface (intrants, irrigation, mécanisation, main-d'œuvre qualifiée). De faibles rendements caractérisent l'agriculture extensive.
\item \textbf{Vrai.} Dans les camps de la CDC (Tiko, Muyuka...), le \og camp des cadres \fg{} (logements spacieux, eau courante, électricité, club) est spatialement distinct des \og camps des ouvriers \fg{} (cases alignées, équipements collectifs sommaires) : c'est une ségrégation socio-spatiale héritée de la période coloniale.
\end{enumerate}
\end{corrige}

\begin{attention}[Points de vigilance]
Pour un \og vrai ou faux justifié \fg, la justification est obligatoire : une réponse sans phrase d'explication ne vaut rien. Citer toujours un fait précis (firme, région, culture) plutôt qu'une formule vague.
\end{attention}

\begin{corrige}[Exercice 3 — Définitions et concepts clés]
\begin{enumerate}
\item \textbf{Agriculture de grandes plantations :} système agricole capitalistique pratiqué sur de vastes domaines (plusieurs milliers d'hectares), fondé sur la monoculture d'une spéculation de rente (banane, palmier, hévéa, canne, coton), employant une main-d'œuvre salariée nombreuse et destiné principalement à l'exportation.
\item \textbf{Agriculture intensive :} mode de production qui recherche des rendements élevés à l'hectare grâce à d'importants investissements par unité de surface : intrants chimiques (engrais, pesticides), semences sélectionnées, irrigation, mécanisation et encadrement technique.
\item \textbf{Agropole :} pôle agro-industriel planifié concentrant sur un même territoire production agricole, unités de transformation, infrastructures (routes, énergie, eau) et services, afin de créer de la valeur ajoutée locale et des emplois (ex. : projets d'agropoles au Cameroun, dont celui de la région du Nord).
\item \textbf{Accaparement des terres :} appropriation de vastes étendues de terres — souvent des terres coutumières ou communautaires — par des entreprises ou des investisseurs, au détriment des populations riveraines qui perdent leurs champs et leurs ressources, source de conflits fonciers (ex. : tensions entre les sociétés de plantations et les communautés riveraines du Sud-Ouest et du Sud).
\end{enumerate}
\end{corrige}

\begin{corrige}[Exercice 4 — Calculs directs et lecture de données]
\textbf{1. Rendements moyens à l'hectare.}

Le rendement se calcule par : $R = \dfrac{\text{production (t)}}{\text{superficie (ha)}}$.

\begin{itemize}
\item \textbf{Banane (CDC) :} attention, la superficie totale de la CDC ($40\,000$ ha) couvre plusieurs cultures ; on calcule ici un rendement \emph{apparent} en rapportant la production de banane à la superficie totale, comme demandé :
\[ R_{\text{banane}} = \frac{250\,000}{40\,000} = 6{,}25 \ \text{t/ha}. \]
\item \textbf{Sucre (SOSUCAM) :}
\[ R_{\text{sucre}} = \frac{150\,000}{14\,000} \approx 10{,}71 \ \text{t/ha}. \]
\end{itemize}

\textbf{Interprétation :} le rendement apparent du sucre ($\approx 10{,}7$ t/ha) est supérieur à celui de la banane ($6{,}25$ t/ha). En réalité, la banane d'exportation atteint couramment 30 à 40 t/ha sur les seules surfaces bananières : le calcul ici sous-estime le rendement réel car la production est rapportée à la superficie totale de la firme. C'est une limite de l'exercice qu'il faut savoir signaler.

\textbf{2. Ratio effectifs / superficie.}

\begin{center}
\begin{tabularx}{\linewidth}{|l|X|X|}\hline
\rowcolor{popblueL} \textbf{Firme} & \textbf{Calcul} & \textbf{Ratio (employés/ha)} \\ \hline
CDC & $22\,000 / 40\,000$ & $0{,}55$ \\ \hline
SOCAPALM & $6\,500 / 58\,000$ & $\approx 0{,}11$ \\ \hline
SOSUCAM & $4\,800 / 14\,000$ & $\approx 0{,}34$ \\ \hline
HEVECAM & $7\,200 / 42\,000$ & $\approx 0{,}17$ \\ \hline
\end{tabularx}
\end{center}

La \textbf{CDC} a le ratio le plus élevé ($0{,}55$ employé/ha, soit environ 55 employés pour 100 ha). \textbf{Interprétation :} la CDC pratique l'agriculture la plus intensive en main-d'œuvre : la banane (récolte hebdomadaire, emballage délicat) et l'hévéa (saignée quotidienne manuelle) exigent beaucoup de bras. À l'inverse, la SOCAPALM ($0{,}11$) illustre une culture plus capitalistique et mécanisable : le palmier à huile demande moins de main-d'œuvre permanente par hectare.

\textbf{3. Masse salariale mensuelle de la CDC.}
\[ 22\,000 \times 150\,000 = 3\,300\,000\,000 \ \text{FCFA} = 3{,}3 \ \text{milliards de FCFA par mois}. \]
Soit environ $3{,}3 \times 12 = 39{,}6$ milliards de FCFA par an (ordre de grandeur).
\end{corrige}

\begin{attention}[Points de vigilance]
Vérifier les unités (t/ha, employés/ha, milliards) et arrondir avec prudence en le signalant. Ne pas confondre \og rendement \fg{} (production/surface) et \og productivité du travail \fg{} (production/employé). Bien poser la division dans le bon sens : c'est l'erreur la plus fréquente.
\end{attention}

\begin{corrige}[Exercice 5 — Localisation et analyse spatiale : Carte muette (partie 1)]
\textbf{1 et 2. Localisation des firmes et cultures (corrigé de la carte).} Voir la carte corrigée ci-dessous.
\end{corrige}

\begin{popfigure}
\begin{center}
\includegraphics[width=\linewidth,height=9.5cm,keepaspectratio]{N01/cmr_map.jpg}
\end{center}
\legende{Fond de carte : carte générale du Cameroun (régions, villes, routes, chemin de fer, golfe de Guinée). Le corrigé situe : 1. les plantations du Sud-Ouest et du Littoral (CDC, SOCAPALM, SAFACAM), près de Buea et de Douala ; 2. la canne à sucre (SOSUCAM) et l'hévéa (HEVECAM) dans le Centre et le Sud ; 3. le coton (SODECOTON) dans le Nord et l'Extrême-Nord (Garoua, Maroua) ; 4. les axes d'évacuation : route et chemin de fer Douala--Yaoundé, Transcamerounais vers Ngaoundéré, route Douala--Kribi, et les ports de Douala et de Kribi. \\ Source : Wikimedia Commons, CIA, domaine public.}
\end{popfigure}

\begin{corrige}[Exercice 5 — Localisation et analyse spatiale (partie 2)]
\textbf{3. Axes d'évacuation :} route et voie ferrée Douala--Yaoundé, Transcamerounais (rail/route) Douala--Ngaoundéré vers le Nord cotonier, route Douala--Kribi ; les produits d'exportation transitent par les ports de \cle{Douala} (banane, coton, huile de palme) et de \cle{Kribi} (port en eau profonde).

\textbf{4. Commentaire organisé (corrigé type, 10 lignes) :}

Les grandes plantations se répartissent de façon très inégale. Elles se concentrent dans le \textbf{Sud forestier et humide} (Sud-Ouest, Littoral, Centre, Sud) : la CDC autour de Tiko et Muyuka, la SOCAPALM et la SAFACAM dans la plaine côtière du Moungo, HEVECAM à Niete, la SOSUCAM à Nkoteng et Mbandjock. Ce littoral et cet arrière-pays bénéficient de \textbf{facteurs physiques favorables} (pluies abondantes, souvent supérieures à $2\,000$ mm sur la côte, températures élevées, sols fertiles d'origine volcanique au pied du mont Cameroun) et de \textbf{facteurs historiques} : les premières plantations ont été créées par les Allemands dès 1884-1914, puis reprises par les Britanniques et les Français. La proximité du \textbf{port de Douala}, première infrastructure d'exportation, a fixé ces implantations. À l'inverse, le \textbf{Nord soudanien}, plus sec, s'est spécialisé dans le coton grâce à la SODECOTON, qui encadre les paysans et évacue la production par l'axe routier et ferroviaire vers Douala. Ainsi, la carte des plantations reflète à la fois les contraintes agro-climatiques, l'héritage colonial et la logistique portuaire.
\end{corrige}

\begin{attention}[Points de vigilance]
Une carte sans titre, sans légende organisée et sans orientation est sanctionnée au Bac. La légende doit être hiérarchisée (d'abord les firmes/cultures, puis les axes, puis les villes et ports). Ne pas placer la SODECOTON dans le Sud ni la SOSUCAM dans le Moungo : ce sont les erreurs de localisation les plus fréquentes.
\end{attention}

\begin{corrige}[Exercice 6 — Croquis d'organisation sociale : Le camp de plantation (partie 1)]
\textbf{Corrigé du croquis (modèle attendu) :} voir le croquis ci-dessous.
\end{corrige}

\begin{popfigure}
\begin{center}
\begin{tikzpicture}[scale=1.0]
% Contour du camp
\draw[thick,popdark,dashed] (0,0) rectangle (9,6);
% Camp des cadres (centre)
\fill[popblueL] (3.2,2.2) rectangle (5.8,4.0);
\draw[popblue,thick] (3.2,2.2) rectangle (5.8,4.0);
\node[font=\scriptsize\bfseries,popdark] at (4.5,3.4) {Camp des cadres};
\node[font=\tiny] at (4.5,2.9) {villas, club-house, tennis};
% Camps ouvriers (périphérie)
\fill[poporangeL] (0.4,0.4) rectangle (2.8,2.0);
\draw[poporange,thick] (0.4,0.4) rectangle (2.8,2.0);
\node[font=\tiny] at (1.6,1.2) {Camp ouvriers Ouest};
\fill[poporangeL] (6.2,0.4) rectangle (8.6,2.0);
\draw[poporange,thick] (6.2,0.4) rectangle (8.6,2.0);
\node[font=\tiny] at (7.4,1.2) {Camp ouvriers Est};
\fill[poporangeL] (0.4,4.4) rectangle (2.8,5.6);
\draw[poporange,thick] (0.4,4.4) rectangle (2.8,5.6);
\node[font=\tiny] at (1.6,5.0) {Camp ouvriers Nord};
% Équipements collectifs (intermédiaires)
\fill[popgreen] (3.3,4.5) circle (0.12) node[right,font=\tiny]{École};
\fill[popgreen] (4.5,4.5) circle (0.12) node[right,font=\tiny]{Dispensaire};
\fill[popgreen] (5.9,4.5) circle (0.12) node[right,font=\tiny]{Marché};
\fill[popgreen] (3.3,1.6) circle (0.12) node[right,font=\tiny]{Lieu de culte};
% Usine / stockage
\fill[poppurple!30] (6.4,4.6) rectangle (8.4,5.6);
\draw[poppurple,thick] (6.4,4.6) rectangle (8.4,5.6);
\node[font=\tiny] at (7.4,5.1) {Usine / stockage};
% Voie ferrée / route
\draw[very thick,popink] (0,3.0) -- (9,3.0);
\node[font=\tiny,popink,rotate=0] at (8.2,2.75) {route / voie ferrée};
% Nord + échelle
\draw[->,thick,popdark] (9.4,5.4) -- (9.4,6.0) node[above,font=\scriptsize]{N};
\draw[thick] (0.4,-0.4) -- (2.4,-0.4) node[midway,below,font=\tiny]{0 \hspace{1.2cm} 400 m};
\end{tikzpicture}
\end{center}
\legende{Croquis : \og Organisation spatiale et sociale d'un camp de plantation (ex. CDC) \fg. Légende : \textcolor{popblue}{cadre bleu} = habitat des cadres (centre) ; \textcolor{poporange}{zones orange} = camps des ouvriers (périphérie) ; \textcolor{popgreen}{$\bullet$} = équipements collectifs (école, dispensaire, marché, culte) ; \textcolor{poppurple}{zone violette} = usine et stockage ; trait noir épais = voie d'évacuation. Échelle : 1 cm $\approx$ 200 m (proche du 1/5000\textsuperscript{e} ramené au schéma).}
\end{popfigure}

\begin{corrige}[Exercice 6 — Croquis d'organisation sociale (partie 2)]
\textbf{Critères de réussite du croquis :}
\begin{enumerate}
\item Un \textbf{cadre}, un \textbf{titre}, une \textbf{orientation} (flèche du nord), une \textbf{échelle} indicative.
\item La \textbf{hiérarchie spatiale} visible : cadres au centre (confort), ouvriers en périphérie (habitat collectif dense), équipements en position intermédiaire.
\item Une \textbf{légende organisée} : 1. habitat (deux catégories), 2. équipements sociaux, 3. production/transports.
\end{enumerate}

\textbf{Lecture géographique :} le croquis met en évidence la \cle{ségrégation socio-spatiale} : la qualité de l'habitat et la proximité des services dépendent du statut professionnel. Mais les équipements collectifs (école, dispensaire, marché) sont des lieux de brassage où se mélangent employés de toutes origines — d'où le rôle d'intégration sociale, malgré ses limites.
\end{corrige}

\begin{attention}[Points de vigilance]
Ne pas dessiner un camp homogène : l'exercice évalue précisément la capacité à figurer le \emph{contraste} cadres/ouvriers. Tout symbole utilisé doit figurer dans la légende ; toute information de l'énoncé doit être reportée.
\end{attention}

\begin{corrige}[Exercice 7 — Diagramme et interprétation : Superficies vs Emplois (partie 1)]
\textbf{1. Diagramme en barres groupées (corrigé) :}

Données converties en milliers : CDC (40 ; 22), SOCAPALM (58 ; 6,5), SOSUCAM (14 ; 4,8), HEVECAM (42 ; 7,2). Voir le diagramme ci-dessous.
\end{corrige}

\begin{popfigure}
\begin{center}
\begin{tikzpicture}
\begin{axis}[
  ybar, bar width=0.32cm,
  width=0.95\linewidth, height=6.5cm,
  ymin=0, ymax=65,
  ylabel={Milliers (ha ou employés)},
  symbolic x coords={CDC, SOCAPALM, SOSUCAM, HEVECAM},
  xtick=data,
  legend style={at={(0.5,1.02)},anchor=south,font=\scriptsize},
  ytick={0,10,20,30,40,50,60},
  nodes near coords, every node near coord/.append style={font=\tiny},
]
\addplot[fill=popblue] coordinates {(CDC,40) (SOCAPALM,58) (SOSUCAM,14) (HEVECAM,42)};
\addplot[fill=poporange] coordinates {(CDC,22) (SOCAPALM,6.5) (SOSUCAM,4.8) (HEVECAM,7.2)};
\legend{Superficie (milliers de ha), Effectifs (milliers d'employés)}
\end{axis}
\end{tikzpicture}
\end{center}
\legende{Superficies cultivées et effectifs salariés de quatre grandes firmes camerounaises (données approximatives, année de référence récente). Source : données d'exercice d'après rapports des entreprises.}
\end{popfigure}

\begin{corrige}[Exercice 7 — Diagramme et interprétation (partie 2)]
\textbf{2. Deux enseignements géographiques :}
\begin{itemize}
\item \textbf{Le rapport capital/travail varie fortement selon les filières.} La SOCAPALM cultive la plus grande superficie (58\,000 ha) avec seulement 6\,500 employés : la palmeraie est capitalistique et peu gourmande en main-d'œuvre permanente. À l'inverse, la CDC emploie 22\,000 personnes sur 40\,000 ha : la banane et l'hévéa exigent une main-d'œuvre abondante et permanente (récolte hebdomadaire, saignée quotidienne). La plantation n'est donc pas un modèle unique : intensité de travail et de capital diffèrent selon la spéculation.
\item \textbf{L'ajout de la SODECOTON changerait l'échelle sociale du graphique.} Avec environ \num{200000} à \num{250000} planteurs encadrés (soit 200 à 250 milliers de producteurs, hors salariés), la barre \og emplois \fg{} exploserait : cela montrerait que le modèle de la \cle{plantation-entreprise} par encadrement paysan crée bien plus d'emplois ruraux que les domaines en propre, et joue un rôle majeur de fixation des populations et d'intégration du Nord dans l'économie monétaire nationale.
\end{itemize}
\end{corrige}

\begin{attention}[Points de vigilance]
Toujours convertir les deux séries dans la même unité (milliers) avant de tracer ; indiquer l'échelle sur l'axe et la source en légende. Un diagramme sans titre d'axes ni légende perd des points.
\end{attention}

\begin{corrige}[Exercice 8 — Étude de cas : La filière coton et la SODECOTON]
\textbf{1. Une \og grande plantation \fg{} au sens large.} La SODECOTON relève de l'agriculture de plantations au sens large (ou \cle{plantation-entreprise}) car elle remplit toutes les fonctions d'une firme de plantation — financement, fourniture d'intrants, encadrement technique, collecte, transformation (égrenage à Garoua, Maroua, Kaélé), commercialisation à l'export — tout en déléguant la production à des centaines de milliers de petits exploitants. L'unité de production est éclatée (exploitations de 2 à 5 ha), mais l'\textbf{organisation économique est celle d'une grande entreprise monocole tournée vers l'exportation}.

\textbf{2. Trois avantages socio-économiques pour l'intégration nationale :}
\begin{itemize}
\item \textbf{Revenus monétaires :} le coton est la première source de revenus en espèces pour des millions de ruraux du Nord, ce qui les intègre à l'économie nationale et marchande.
\item \textbf{Aménagement et infrastructures :} usines d'égrenage, pistes rurales d'évacuation, entrepôts et services (crédit, vulgarisation) désenclavent les régions septentrionales.
\item \textbf{Emplois et fixation des populations :} la filière fait vivre producteurs, saisonniers, transporteurs et ouvriers d'usine, limitant l'exode rural vers les villes du Sud.
\end{itemize}

\textbf{3. Deux risques environnementaux :}
\begin{itemize}
\item \textbf{Dégradation des sols :} l'usage intensif d'engrais et la monoculture continue épuisent la fertilité et acidifient les sols.
\item \textbf{Pollution des eaux :} les pesticides et engrais lessivés contaminent les nappes phréatiques et les cours d'eau, avec des risques sanitaires pour les populations.
\end{itemize}

\textbf{4. Structuration du territoire septentrional.} La filière organise l'espace du Nord en réseau : les \textbf{pistes rurales} convergent vers les centres de collecte villageois, puis vers les \textbf{usines d'égrenage} (Garoua, Maroua, Kaélé), véritables pôles urbains et industriels ; la fibre est ensuite acheminée par l'\textbf{axe routier et ferroviaire Nord--Douala} (Transcamerounais) jusqu'au port d'exportation. Ainsi, la SODECOTON crée une armature territoriale (bourgades de collecte, villes-usines, corridors de transport) qui insère les régions septentrionales, autrefois enclavées, dans les circuits nationaux et internationaux.
\end{corrige}

\begin{attention}[Points de vigilance]
Bien distinguer \og plantation en propre \fg{} (domaine d'un seul tenant : CDC, SOCAPALM) et \og plantation-entreprise \fg{} (encadrement de planteurs : SODECOTON). Au Bac, exiger cette nuance montre la maîtrise du cours.
\end{attention}

\begin{corrige}[Exercice 9 — Sujet type Bac : Définition, exemples et impacts]
\textbf{Barème indicatif (sur 20) :} introduction 3 pts ; caractéristiques définies et illustrées 7 pts ; impacts positifs et négatifs équilibrés 7 pts ; conclusion nuancée 2 pts ; langue et organisation 1 pt.

\textbf{1. Introduction (corrigé type).}

\emph{Accroche :} Au pied du mont Cameroun, les bananeraies de la CDC s'étendent à perte de vue et alimentent chaque semaine les navires du port de Douala. \emph{Définition :} l'agriculture de grandes plantations est un système capitalistique, monocole et tourné vers l'exportation, pratiqué sur de vastes domaines avec une main-d'œuvre salariée nombreuse. \emph{Annonce du plan :} nous présenterons d'abord ses caractéristiques à travers la CDC et la SOSUCAM, puis ses impacts socio-économiques, positifs et négatifs.

\textbf{2. Partie 1 — Caractéristiques illustrées.}
\begin{itemize}
\item \textbf{Superficies :} domaines de plusieurs dizaines de milliers d'hectares (CDC : environ 40\,000 ha ; SOSUCAM : environ 14\,000 ha à Nkoteng et Mbandjock, ordres de grandeur).
\item \textbf{Capital :} financement massif, public ou mixte (CDC à majorité publique ; SOSUCAM, filiale du groupe SOMDIAA).
\item \textbf{Monoculture de rente :} banane, palmier et hévéa pour la CDC (Sud-Ouest) ; canne à sucre pour la SOSUCAM (Haute-Sanaga, région du Centre).
\item \textbf{Main-d'œuvre salariée nombreuse :} environ 22\,000 employés à la CDC, près de 5\,000 à la SOSUCAM, plus les saisonniers.
\item \textbf{Exportation :} bananes vers l'Europe via Douala ; sucre pour le marché national et sous-régional.
\end{itemize}

\textbf{3. Partie 2 — Impacts.}
\begin{itemize}
\item \textbf{Positifs :} emplois stables et salaires réguliers ; équipements sociaux (écoles, centres de santé, eau, électricité dans les camps) ; routes et voies ferrées ; recettes d'exportation et devises ; formation professionnelle.
\item \textbf{Négatifs :} accaparement des terres communautaires et conflits fonciers avec les riverains ; pollution chimique (engrais, pesticides, effluents d'usines) ; dépendance à une seule culture et aux cours mondiaux ; conditions de travail pénibles et précarité des contractuels.
\end{itemize}

\textbf{4. Conclusion (corrigé type).} Les grandes plantations sont un pilier de l'économie camerounaise et un outil d'aménagement, mais leur modèle foncier et environnemental est contesté. L'avenir passe par les \cle{agropoles}, la diversification et une certification durable conciliant rentabilité et droits des communautés.
\end{corrige}

\begin{attention}[Points de vigilance]
Le sujet exige \emph{deux} firmes nommées et localisées : une réponse sans exemple camerounais précis est hors sujet. Équilibrer strictement impacts positifs et négatifs ; citer au moins un chiffre par firme.
\end{attention}

\begin{corrige}[Exercice 10 — Sujet type Bac : Croquis commenté et main-d'œuvre (partie 1)]
\textbf{Barème indicatif (sur 20) :} croquis (titre, cadre, orientation, échelle, légende hiérarchisée, exactitude) 8 pts ; commentaire structuré 10 pts ; langue 2 pts.

\textbf{1. Croquis (corrigé type) :} voir le croquis ci-dessous.
\end{corrige}

\begin{popfigure}
\begin{center}
\begin{tikzpicture}[scale=0.95]
\draw[thick,popdark] (0,0) rectangle (9.5,6.2);
% Habitat cadres
\fill[popblueL] (3.4,2.6) rectangle (6.0,4.2);
\draw[popblue,thick] (3.4,2.6) rectangle (6.0,4.2);
\node[font=\scriptsize\bfseries,popdark] at (4.7,3.6) {Habitat cadres};
\node[font=\tiny] at (4.7,3.1) {villas, eau, électricité, club};
% Habitat ouvriers
\fill[poporangeL] (0.4,0.5) rectangle (3.0,2.2);
\draw[poporange,thick] (0.4,0.5) rectangle (3.0,2.2);
\node[font=\tiny] at (1.7,1.35) {Habitat ouvriers (cases alignées)};
\fill[poporangeL] (6.4,0.5) rectangle (9.0,2.2);
\draw[poporange,thick] (6.4,0.5) rectangle (9.0,2.2);
\node[font=\tiny] at (7.7,1.35) {Habitat ouvriers};
% Équipements
\fill[popgreen] (3.5,4.7) circle (0.13) node[right,font=\tiny]{École};
\fill[popgreen] (4.8,4.7) circle (0.13) node[right,font=\tiny]{Centre de santé};
\fill[popgreen] (6.6,4.7) circle (0.13) node[right,font=\tiny]{Marché};
\fill[popgreen] (3.5,2.0) circle (0.13) node[right,font=\tiny]{Lieu de culte};
% Usine / stockage
\fill[poppurple!30] (6.6,4.4) rectangle (9.0,5.7);
\draw[poppurple,thick] (6.6,4.4) rectangle (9.0,5.7);
\node[font=\tiny] at (7.8,5.05) {Usine / zone de stockage};
% Voies
\draw[very thick,popink] (0,2.45) -- (9.5,2.45);
\node[font=\tiny,popink] at (1.2,2.7) {route principale};
\draw[thick,popink,dashed] (7.8,4.4) -- (7.8,2.45);
% Nord + échelle
\draw[->,thick,popdark] (9.9,5.4) -- (9.9,6.1) node[above,font=\scriptsize]{N};
\draw[thick] (0.4,-0.4) -- (2.4,-0.4) node[midway,below,font=\tiny]{0 \hspace{1.2cm} 400 m};
\end{tikzpicture}
\end{center}
\legende{\og Organisation spatiale et sociale d'un camp de plantation \fg. Légende : 1. Habitat : \textcolor{popblue}{bleu} = cadres (centre), \textcolor{poporange}{orange} = ouvriers (périphérie) ; 2. Équipements collectifs : \textcolor{popgreen}{$\bullet$} école, santé, marché, culte ; 3. Production : \textcolor{poppurple}{violet} = usine et stockage ; 4. Transports : trait plein = route, tirets = piste vers l'usine. Échelle indicative : 1 cm $\approx$ 200 m.}
\end{popfigure}

\begin{corrige}[Exercice 10 — Sujet type Bac (partie 2)]
\textbf{2. Commentaire structuré (corrigé type, 15-20 lignes).}

Les grandes plantations camerounaises (CDC, HEVECAM, SOCAPALM, SOSUCAM, SAFACAM) emploient plus de 50\,000 permanents et des dizaines de milliers de saisonniers. Cette main-d'œuvre constitue d'abord un \textbf{vecteur de brassage des populations} : recrutés dans toutes les régions — Nord, Ouest, Centre, Sud — et parfois dans les pays voisins (Nigéria, Tchad, RCA), les travailleurs cohabitent dans les camps, où ethnies et langues diverses se mélangent quotidiennement. Le \textbf{français et le pidgin-english} y servent de langues interethniques communes, forgeant une culture partagée. Ensuite, les \textbf{services sociaux} gérés par les firmes — écoles primaires et secondaires, centres de santé, adduction d'eau, électricité — profitent aux familles et diffusent des normes communes (scolarisation, santé) au-delà des clivages régionaux. Enfin, la \textbf{mobilité géographique} qu'implique le travail en plantation fait circuler hommes, savoirs et revenus entre régions : les salaires rapatriés nourrissent les villages d'origine. L'intégration nationale se vit donc concrètement dans les camps. Cependant, ce modèle a des \textbf{limites} : la ségrégation spatiale entre camp des cadres et camps d'ouvriers reproduit les hiérarchies sociales ; les écarts de salaires et la précarité des contractuels et saisonniers entretiennent les frustrations ; et l'installation des plantations s'accompagne souvent de \textbf{tensions foncières} avec les riverains dépossédés de terres coutumières. La main-d'œuvre des plantations est ainsi un vecteur réel mais inachevé d'intégration nationale.
\end{corrige}

\begin{attention}[Points de vigilance]
Le croquis et le commentaire sont \emph{complémentaires} : le commentaire doit renvoyer à ce que montre le croquis (contraste cadres/ouvriers, équipements partagés). Ne pas oublier la \og limite \fg{} du vecteur d'intégration : un commentaire uniquement élogieux est incomplet.
\end{attention}

\begin{corrige}[Exercice 11 — Sujet type Bac : Dissertation argumentée]
\textbf{Barème indicatif (sur 20) :} introduction (accroche, définitions, problématique, plan) 4 pts ; partie 1 argumentée et illustrée 6 pts ; partie 2 argumentée et illustrée 6 pts ; conclusion nuancée avec ouverture 3 pts ; langue et cohérence 1 pt.

\textbf{1. Introduction (corrigé type).}

\emph{Accroche :} Avec plus de 50\,000 salariés permanents et des centaines de milliers d'hectares, les grandes plantations pèsent lourd dans l'économie camerounaise et dans l'organisation de ses territoires. \emph{Définition et analyse des termes :} les grandes plantations sont des exploitations capitalistiques, monocoles et exportatrices ; être un \og atout pour l'intégration nationale \fg{} signifie renforcer l'unité économique, sociale et territoriale du pays ; être une \og source de tensions \fg{} renvoie aux conflits fonciers, sociaux et environnementaux qu'elles suscitent. \emph{Problématique :} dans quelle mesure les grandes plantations contribuent-elles à l'intégration nationale, et où sont les limites de cette contribution ? \emph{Annonce du plan :} nous montrerons d'abord qu'elles constituent un atout majeur pour l'intégration et le développement des territoires, puis qu'elles génèrent aussi des tensions foncières, sociales et environnementales.

\textbf{2. Première partie — Un atout majeur pour l'intégration nationale.}
\begin{itemize}
\item \textbf{Emplois et brassage :} plus de 50\,000 permanents (CDC : 22\,000 ; HEVECAM : 7\,200 ; SOCAPALM : 6\,500...) venus de toutes les régions, plus les saisonniers ; les camps sont des creusets interethniques où français et pidgin servent de langues communes.
\item \textbf{Infrastructures sociales :} écoles, centres de santé, eau et électricité gérés par les firmes bénéficient aux employés et aux riverains.
\item \textbf{Désenclavement :} routes, voie ferrée (Transcamerounais pour le coton), ports de Douala et Kribi ouvrent les zones de production.
\item \textbf{Recettes et devises :} banane, huile de palme, caoutchouc, coton et sucre rapportent des devises et des recettes fiscales à l'État.
\item \textbf{Formation et urbanisation :} formation technique des salariés ; naissance de bourgs dynamiques autour des camps et usines (Tiko, Niete, Mbandjock, Garoua).
\end{itemize}

\textbf{3. Deuxième partie — Des sources de tensions.}
\begin{itemize}
\item \textbf{Tensions foncières :} accaparement des terres coutumières au profit de titres fonciers des firmes ; conflits récurrents entre les sociétés de plantations et les communautés riveraines du Sud-Ouest et du Sud.
\item \textbf{Tensions environnementales :} pollution chimique (effluents des sucreries SOSUCAM, pesticides de la filière coton dans le Nord), dégradation des sols et des nappes.
\item \textbf{Tensions sociales :} précarité des contractuels et saisonniers, écarts de salaires, ségrégation spatiale héritée (camps cadres/ouvriers), grèves et conflits du travail.
\item \textbf{Vulnérabilité économique :} dépendance aux cours mondiaux des matières premières ; la monoculture expose territoires et salariés aux crises.
\end{itemize}

\textbf{4. Conclusion (corrigé type).} Les grandes plantations sont indéniablement un atout pour l'intégration nationale — emplois, brassage, infrastructures, devises — mais elles cristallisent aussi des tensions foncières, sociales et environnementales qui en limitent les bienfaits. La réponse à la problématique est donc nuancée : tout dépend de la gouvernance. \emph{Ouverture :} les réformes foncières sécurisant les droits coutumiers, la responsabilité sociétale des entreprises (RSE), les certifications durables, la diversification et les agropoles pourraient concilier rentabilité des firmes et droits des communautés.
\end{corrige}

\begin{attention}[Points de vigilance]
Une dissertation exige un \emph{plan dialectique équilibré} : les deux parties doivent avoir un poids comparable. Chaque argument doit être illustré par un exemple camerounais précis (firme, région, chiffre). La problématique doit être explicitement formulée dans l'introduction et reprise en conclusion. Éviter le catalogue de notions sans articulation : utiliser des connecteurs logiques (d'abord, ensuite, cependant, enfin).
\end{attention}

\newpage

\coursec{Fiche bilan}

\begin{popfigure}
\begin{tikzpicture}[mindmap, grow cyclic, every node/.style={concept, text=white, font=\small\bfseries},
  root concept/.append style={concept color=popink, font=\bfseries},
  level 1/.append style={level distance=3.4cm, sibling angle=72},
  level 2/.append style={level distance=2.1cm, sibling angle=45, font=\scriptsize},
  concept color=popblue]
\node[root concept]{L'agriculture de grandes plantations}
  child[concept color=popblue]{ node {Caractéristiques}
    child { node {Grandes superficies} }
    child { node {Monoculture, mécanisation} }
    child { node {Financement : capitaux publics et privés} }
  }
  child[concept color=popgreen]{ node {Plantations des grandes firmes}
    child { node {CDC, HEVECAM} }
    child { node {SOCAPALM, SAFACAM} }
    child { node {SODECOTON, SOSUCAM} }
  }
  child[concept color=poporange]{ node {Plantations des particuliers}
    child { node {Élites, planteurs locaux} }
    child { node {Cacao, café, palme} }
  }
  child[concept color=poppurple]{ node {Impacts socio-économiques}
    child { node {Emplois, revenus, infrastructures} }
    child { node {Accaparement des terres} }
    child { node {Pollution, conflits} }
  }
  child[concept color=popgold]{ node {Main-d'œuvre et intégration}
    child { node {Origines diverses des travailleurs} }
    child { node {Camps, écoles, centres de santé} }
  };
\end{tikzpicture}
\legende{Carte mentale de la leçon : les cinq branches essentielles à maîtriser (définition et caractéristiques ; firmes ; particuliers ; impacts ; main-d'œuvre et intégration nationale).}
\end{popfigure}

\begin{bilan}
\textbf{Définitions clés à connaître par cœur}
\begin{itemize}
  \item \cle{Agriculture de grandes plantations} : système de production agricole pratiqué sur de vastes superficies (souvent plusieurs centaines ou milliers d'hectares), fondé sur la \cle{monoculture} d'une spéculation destinée à l'exportation ou à l'industrie, avec des méthodes modernes (mécanisation, intrants) et un financement capitalistique.
  \item \cle{Agriculture intensive} : agriculture qui recherche des rendements élevés à l'hectare grâce à de forts investissements (machines, engrais, pesticides, main-d'œuvre qualifiée).
  \item \cle{Agropole} : espace aménagé regroupant plantations, usines de transformation et infrastructures sociales autour d'une même filière agricole (ex. : la plaine de la SODECOTON autour de Garoua).
\end{itemize}

\textbf{Les grandes firmes et leurs cultures (à réciter sans hésiter)}
\begin{center}
\begin{tabularx}{\linewidth}{|l|X|X|}\hline
\rowcolor{popblueL} \textbf{Firme} & \textbf{Culture principale} & \textbf{Localisation} \\ \hline
CDC & Banane, palmier à huile, hévéa & Littoral (Tiko, Mbanga) \\ \hline
SOCAPALM & Palmier à huile & Littoral et Sud (Mbongo, Dibombari, Eseka) \\ \hline
SAFACAM & Palmier à huile & Littoral (Dibamba) \\ \hline
SODECOTON & Coton & Nord, Extrême-Nord, Adamaoua (Garoua) \\ \hline
SOSUCAM & Canne à sucre & Centre (Mbandjock, Nkoteng) \\ \hline
HEVECAM & Hévéa (caoutchouc) & Sud (Niete, près de Kribi) \\ \hline
\end{tabularx}
\end{center}

\textbf{Méthodes express}
\begin{itemize}
  \item \emph{Lire une carte des plantations} : repérer d'abord la localisation (région, proximité d'un port ou d'un axe), puis la culture, puis la firme ; conclure sur les facteurs d'implantation (climat, sols, main-d'œuvre, débouchés).
  \item \emph{Construire un croquis} : titre, légende numérotée, orientation (flèche du nord), échelle indicative ; figurer la plantation, les camps, les voies et les unités de transformation.
  \item \emph{Commenter un impact} : toujours équilibrer le bilan (positifs : emplois, revenus, écoles, routes ; négatifs : accaparement des terres communautaires, pollution, dépendance aux cours mondiaux).
  \item \emph{Traiter le dossier 2} : montrer comment la diversité d'origine des travailleurs et l'organisation des camps (cadres / ouvriers, écoles, centres de santé) font des plantations un vecteur d'\cle{intégration nationale}.
\end{itemize}
\end{bilan}

\begin{aretenir}[Checklist avant le Bac]
\begin{itemize}
  \item $\square$ Je sais définir l'agriculture de grandes plantations et citer ses trois caractéristiques (superficie, méthodes, financement).
  \item $\square$ Je distingue agriculture intensive et agriculture extensive.
  \item $\square$ Je sais définir une agropole et donner un exemple camerounais.
  \item $\square$ Je connais les six grandes firmes (CDC, SOCAPALM, SAFACAM, SODECOTON, SOSUCAM, HEVECAM) avec leur culture et leur localisation.
  \item $\square$ Je sais citer des exemples de plantations de particuliers (cacaoyers, caféiers, palmeraies villageoises).
  \item $\square$ Je distingue les deux types de grandes plantations (firmes / particuliers).
  \item $\square$ Je sais dresser un bilan équilibré des impacts socio-économiques (emplois et conditions de vie contre accaparement des terres et pollution).
  \item $\square$ Je sais expliquer en quoi la main-d'œuvre des plantations est un vecteur d'intégration nationale.
  \item $\square$ Je sais décrire l'organisation spatiale et sociale d'une plantation (camps des cadres et des ouvriers, écoles, centres de santé).
  \item $\square$ Je sais légender une carte et construire un croquis de plantation avec titre, légende, nord et échelle.
  \item $\square$ Je sais commenter un tableau statistique de production (calculer une part en \% et un taux de variation).
  \item $\square$ Je vérifie toujours que ma copie relie la plantation étudiée à la famille de situations : l'intégration nationale.
\end{itemize}
\end{aretenir}

\findecours

\end{document}
